% !TEX program = lualatex
% ---------------------------------------------------------------------------------------------
%  Readback note 28: The tower method at a place: floor measures, gauges and the bounds on each floor
%
%  Build:  python3 tools/rb.py note <slug>
%
%  One unit of a Lean artifact, rewritten as ordinary mathematics. The prose is hand-written. Lean names,
%  line ranges and counts come from data.tex; quoted Lean statements are copied from the artifact by the build.
%
%  Rules of the text: standard terminology; every symbol is explained where it first appears and listed again
%  in the notation at the end; no letter has two meanings; a proof is complete or says what it takes for granted;
%  Lean names appear only in the LEAN lines and in the concordance.
%
%  The signature rule (the build enforces it): a function, relation or predicate that the artifact defines is
%  written only through its notation macro, which the build makes from the spec (`notation`) with one slot for
%  each explicit argument of the Lean definition. No argument is fixed and dropped, also not "for brevity" and
%  not for one section. Every argument and every hypothesis of a stated declaration stands inside its
%  statement, or in a Setting that the statement cites by number; only implicit and instance arguments (the
%  ambient objects) may stand in the conventions.
% ---------------------------------------------------------------------------------------------
\documentclass[11pt,a4paper]{article}
\ifdefined\pdfvariable\pdfvariable suppressoptionalinfo \numexpr 1+2+4+8\relax\fi
\usepackage[a4paper,hmargin=35mm,top=22mm,bottom=27mm,footskip=12mm,nohead]{geometry}

\usepackage{fontspec}
\directlua{
  luaotfload.add_fallback("gfallback", {"STIXTwoMath-Regular.otf:mode=node;", "DejaVuSans.ttf:mode=node;"})
  luaotfload.add_fallback("gmonofallback", {"DejaVuSansMono.ttf:mode=node;", "STIXTwoMath-Regular.otf:mode=node;"})
}
\setmainfont{NewCM10}[Extension=.otf,UprightFont=*-Book,ItalicFont=*-BookItalic,
  BoldFont=*-Bold,BoldItalicFont=*-BoldItalic,SmallCapsFont=*-Book,SmallCapsFeatures={Letters=SmallCaps},
  RawFeature={fallback=gfallback}]
\setsansfont{NewCMSans10}[Extension=.otf,UprightFont=*-Book,ItalicFont=*-BookOblique,
  BoldFont=*-Bold,BoldItalicFont=*-BoldOblique,RawFeature={fallback=gfallback}]
\setmonofont{JetBrainsMono}[Extension=.otf,UprightFont=*-Regular,ItalicFont=*-Italic,
  BoldFont=*-SemiBold,BoldItalicFont=*-SemiBoldItalic,Scale=MatchLowercase,
  RawFeature={-calt;-liga;fallback=gmonofallback}]
\usepackage{amsmath}
\usepackage{amsthm}
\usepackage[math-style=TeX]{unicode-math}
\setmathfont{NewCMMath-Book.otf}
\emergencystretch=2.5em\relax

\usepackage{xcolor}
\usepackage{tikz}
\usetikzlibrary{calc,arrows.meta,positioning,backgrounds}
\usepackage[protrusion=true,expansion=true]{microtype}
\usepackage{enumitem}
\usepackage{needspace}
\usepackage{booktabs}
\usepackage{array}
\usepackage{longtable}
\usepackage{fancyhdr}
\usepackage{lastpage}
\usepackage{multicol}
\usepackage[hidelinks,pdfusetitle]{hyperref}

\definecolor{ink}{HTML}{0F172A}
\definecolor{inksoft}{HTML}{475569}
\definecolor{inkfaint}{HTML}{94A3B8}
\definecolor{rulegrey}{HTML}{CBD5E1}
\definecolor{arith}{HTML}{C2410C}
\definecolor{est}{HTML}{7C3AED}
\definecolor{ana}{HTML}{0891B2}
\definecolor{ok}{HTML}{047857}
\definecolor{leancol}{HTML}{3730A3}

\title{The tower method at a place: floor measures, gauges and the bounds on each floor}
\author{grh-compact}
\date{}

% ------------------------------------------------------------------------------------- macros
\newcommand{\leanfont}{\ttfamily\fontsize{8.6}{10.5}\selectfont\color{leancol}}
\newcommand{\lean}[1]{\csname LN@#1\endcsname}
\newcommand{\fillin}[1]{{\color{arith}\itshape[#1]}}  % a place of the template that is not written yet; the build reports every one that is left
\newcommand{\leanquote}[1]{\csname LQ@#1\endcsname}
% generated by rb.py note; do not edit
\newcommand{\ConcDecls}{116}
\newcommand{\Pin}{72fead10}
\newcommand{\PinText}{that of commit {\leanfont 72fead10}}
\newcommand{\ArtifactName}{grh-compact}
\expandafter\newcommand\csname LN@memvecsOfSizeofheisenbergElem\endcsname{\href{https://github.com/lsernau/grh-lean/blob/72fead102f663d5c7a407964a13dc0cf9d977207/grh-compact/src/New/RealPlace/FlatCharts.lean\#L21-L26}{{\leanfont Size.\allowbreak{}mem\_vecsOfSize\_of\_heisenbergElem}}}
\expandafter\newcommand\csname LN@isNearmuliff\endcsname{\href{https://github.com/lsernau/grh-lean/blob/72fead102f663d5c7a407964a13dc0cf9d977207/grh-compact/src/New/RealPlace/FlatCharts.lean\#L28-L29}{{\leanfont Level.\allowbreak{}isNear\_mul\_iff}}}
\expandafter\newcommand\csname LN@offloorGroupmul\endcsname{\href{https://github.com/lsernau/grh-lean/blob/72fead102f663d5c7a407964a13dc0cf9d977207/grh-compact/src/New/RealPlace/FlatCharts.lean\#L31-L33}{{\leanfont Level.\allowbreak{}IsNear.\allowbreak{}of\_floorGroup\_mul}}}
\expandafter\newcommand\csname LN@measurableSetisNear\endcsname{\href{https://github.com/lsernau/grh-lean/blob/72fead102f663d5c7a407964a13dc0cf9d977207/grh-compact/src/New/RealPlace/FlatCharts.lean\#L35-L37}{{\leanfont Level.\allowbreak{}measurableSet\_isNear}}}
\expandafter\newcommand\csname LN@upperMeasureisNearle\endcsname{\href{https://github.com/lsernau/grh-lean/blob/72fead102f663d5c7a407964a13dc0cf9d977207/grh-compact/src/New/RealPlace/FlatCharts.lean\#L41-L51}{{\leanfont Level.\allowbreak{}upperMeasure\_isNear\_le}}}
\expandafter\newcommand\csname LN@isOpenfloorValues\endcsname{\href{https://github.com/lsernau/grh-lean/blob/72fead102f663d5c7a407964a13dc0cf9d977207/grh-compact/src/New/RealPlace/FlatCharts.lean\#L53-L55}{{\leanfont Level.\allowbreak{}isOpen\_floorValues}}}
\expandafter\newcommand\csname LN@floorValuesmul\endcsname{\href{https://github.com/lsernau/grh-lean/blob/72fead102f663d5c7a407964a13dc0cf9d977207/grh-compact/src/New/RealPlace/FlatCharts.lean\#L57-L61}{{\leanfont Level.\allowbreak{}floorValues\_mul}}}
\expandafter\newcommand\csname LN@volumefloorValuesnezeronetop\endcsname{\href{https://github.com/lsernau/grh-lean/blob/72fead102f663d5c7a407964a13dc0cf9d977207/grh-compact/src/New/RealPlace/FlatCharts.lean\#L63-L66}{{\leanfont Level.\allowbreak{}volume\_floorValues\_ne\_zero\_ne\_top}}}
\expandafter\newcommand\csname LN@floorWeightsucc\endcsname{\href{https://github.com/lsernau/grh-lean/blob/72fead102f663d5c7a407964a13dc0cf9d977207/grh-compact/src/New/RealPlace/FlatCharts.lean\#L68-L70}{{\leanfont Level.\allowbreak{}floorWeight\_succ}}}
\expandafter\newcommand\csname LN@floorWeightnezeroandnetop\endcsname{\href{https://github.com/lsernau/grh-lean/blob/72fead102f663d5c7a407964a13dc0cf9d977207/grh-compact/src/New/RealPlace/FlatCharts.lean\#L72-L73}{{\leanfont Level.\allowbreak{}floorWeight\_ne\_zero\_and\_ne\_top}}}
\expandafter\newcommand\csname LN@floorWeightmulleft\endcsname{\href{https://github.com/lsernau/grh-lean/blob/72fead102f663d5c7a407964a13dc0cf9d977207/grh-compact/src/New/RealPlace/FlatCharts.lean\#L76-L77}{{\leanfont Level.\allowbreak{}floorWeight\_mul\_left}}}
\expandafter\newcommand\csname LN@invWeightmul\endcsname{\href{https://github.com/lsernau/grh-lean/blob/72fead102f663d5c7a407964a13dc0cf9d977207/grh-compact/src/New/RealPlace/FlatCharts.lean\#L80-L83}{{\leanfont Level.\allowbreak{}invWeight\_mul}}}
\expandafter\newcommand\csname LN@measurablefloorWeight\endcsname{\href{https://github.com/lsernau/grh-lean/blob/72fead102f663d5c7a407964a13dc0cf9d977207/grh-compact/src/New/RealPlace/FlatCharts.lean\#L84-L85}{{\leanfont measurable\_floorWeight}}}
\expandafter\newcommand\csname LN@measurableinvWeight\endcsname{\href{https://github.com/lsernau/grh-lean/blob/72fead102f663d5c7a407964a13dc0cf9d977207/grh-compact/src/New/RealPlace/FlatCharts.lean\#L87-L88}{{\leanfont Level.\allowbreak{}measurable\_invWeight}}}
\expandafter\newcommand\csname LN@measurableSetvecsOfSize\endcsname{\href{https://github.com/lsernau/grh-lean/blob/72fead102f663d5c7a407964a13dc0cf9d977207/grh-compact/src/New/RealPlace/FlatCharts.lean\#L92-L93}{{\leanfont Gauge.\allowbreak{}measurableSet\_vecsOfSize}}}
\expandafter\newcommand\csname LN@memvecsOfSizeoffloorBasismulVec\endcsname{\href{https://github.com/lsernau/grh-lean/blob/72fead102f663d5c7a407964a13dc0cf9d977207/grh-compact/src/New/RealPlace/FlatCharts.lean\#L95-L99}{{\leanfont Gauge.\allowbreak{}mem\_vecsOfSize\_of\_floorBasis\_mulVec}}}
\expandafter\newcommand\csname LN@HasFibreBound\endcsname{\href{https://github.com/lsernau/grh-lean/blob/72fead102f663d5c7a407964a13dc0cf9d977207/grh-compact/src/New/RealPlace/FlatCharts.lean\#L103-L106}{{\leanfont Gauge.\allowbreak{}HasFibreBound}}}
\expandafter\newcommand\csname LN@setLIntegralleofhasFibreBound\endcsname{\href{https://github.com/lsernau/grh-lean/blob/72fead102f663d5c7a407964a13dc0cf9d977207/grh-compact/src/New/RealPlace/FlatCharts.lean\#L108-L118}{{\leanfont Gauge.\allowbreak{}setLIntegral\_le\_of\_hasFibreBound}}}
\expandafter\newcommand\csname LN@hasFibreBoundfloorQuadricofle\endcsname{\href{https://github.com/lsernau/grh-lean/blob/72fead102f663d5c7a407964a13dc0cf9d977207/grh-compact/src/New/RealPlace/FlatCharts.lean\#L123-L131}{{\leanfont Gauge.\allowbreak{}hasFibreBound\_floorQuadric\_of\_le}}}
\expandafter\newcommand\csname LN@hasFibreBoundfloorQuadric\endcsname{\href{https://github.com/lsernau/grh-lean/blob/72fead102f663d5c7a407964a13dc0cf9d977207/grh-compact/src/New/RealPlace/FlatCharts.lean\#L135-L136}{{\leanfont Gauge.\allowbreak{}hasFibreBound\_floorQuadric}}}
\expandafter\newcommand\csname LN@hasFibreBoundfloorVolume\endcsname{\href{https://github.com/lsernau/grh-lean/blob/72fead102f663d5c7a407964a13dc0cf9d977207/grh-compact/src/New/RealPlace/FlatCharts.lean\#L141-L148}{{\leanfont Gauge.\allowbreak{}hasFibreBound\_floorVolume}}}
\expandafter\newcommand\csname LN@lintegralheisenbergMeasurele\endcsname{\href{https://github.com/lsernau/grh-lean/blob/72fead102f663d5c7a407964a13dc0cf9d977207/grh-compact/src/New/RealPlace/FlatCharts.lean\#L153-L173}{{\leanfont Gauge.\allowbreak{}lintegral\_heisenbergMeasure\_le}}}
\expandafter\newcommand\csname LN@setLIntegralorthonormalPartle\endcsname{\href{https://github.com/lsernau/grh-lean/blob/72fead102f663d5c7a407964a13dc0cf9d977207/grh-compact/src/New/RealPlace/FlatCharts.lean\#L175-L181}{{\leanfont Gauge.\allowbreak{}setLIntegral\_orthonormalPart\_le}}}
\expandafter\newcommand\csname LN@measurableSetpositionsOfSize\endcsname{\href{https://github.com/lsernau/grh-lean/blob/72fead102f663d5c7a407964a13dc0cf9d977207/grh-compact/src/New/RealPlace/FlatCharts.lean\#L183-L184}{{\leanfont Gauge.\allowbreak{}measurableSet\_positionsOfSize}}}
\expandafter\newcommand\csname LN@conjmulinvclosetoone\endcsname{\href{https://github.com/lsernau/grh-lean/blob/72fead102f663d5c7a407964a13dc0cf9d977207/grh-compact/src/New/RealPlace/FlatCharts.lean\#L188-L192}{{\leanfont Gauge.\allowbreak{}conj\_mul\_inv\_close\_to\_one}}}
\expandafter\newcommand\csname LN@orbitPointmemvecsOfSize\endcsname{\href{https://github.com/lsernau/grh-lean/blob/72fead102f663d5c7a407964a13dc0cf9d977207/grh-compact/src/New/RealPlace/FlatCharts.lean\#L196-L199}{{\leanfont Gauge.\allowbreak{}orbitPoint\_mem\_vecsOfSize}}}
\expandafter\newcommand\csname LN@existshasMatSizeofisNear\endcsname{\href{https://github.com/lsernau/grh-lean/blob/72fead102f663d5c7a407964a13dc0cf9d977207/grh-compact/src/New/RealPlace/FlatCharts.lean\#L203-L209}{{\leanfont Gauge.\allowbreak{}exists\_hasMatSize\_of\_isNear}}}
\expandafter\newcommand\csname LN@hasMatSizesuboneoffloorValues\endcsname{\href{https://github.com/lsernau/grh-lean/blob/72fead102f663d5c7a407964a13dc0cf9d977207/grh-compact/src/New/RealPlace/FlatCharts.lean\#L215-L233}{{\leanfont Gauge.\allowbreak{}hasMatSize\_sub\_one\_of\_floorValues}}}
\expandafter\newcommand\csname LN@existsnezeroofpivotModulus\endcsname{\href{https://github.com/lsernau/grh-lean/blob/72fead102f663d5c7a407964a13dc0cf9d977207/grh-compact/src/New/RealPlace/FloorKernels.lean\#L18-L21}{{\leanfont Level.\allowbreak{}exists\_ne\_zero\_of\_pivotModulus}}}
\expandafter\newcommand\csname LN@lintegralupperPairingsqle\endcsname{\href{https://github.com/lsernau/grh-lean/blob/72fead102f663d5c7a407964a13dc0cf9d977207/grh-compact/src/New/RealPlace/FloorKernels.lean\#L23-L26}{{\leanfont Level.\allowbreak{}lintegral\_upperPairing\_sq\_le}}}
\expandafter\newcommand\csname LN@stronglyMeasurableupperPairing\endcsname{\href{https://github.com/lsernau/grh-lean/blob/72fead102f663d5c7a407964a13dc0cf9d977207/grh-compact/src/New/RealPlace/FloorKernels.lean\#L32-L33}{{\leanfont Level.\allowbreak{}IsPairingKernel.\allowbreak{}stronglyMeasurable\_upperPairing}}}
\expandafter\newcommand\csname LN@mono\endcsname{\href{https://github.com/lsernau/grh-lean/blob/72fead102f663d5c7a407964a13dc0cf9d977207/grh-compact/src/New/RealPlace/FloorKernels.lean\#L35-L36}{{\leanfont Level.\allowbreak{}IsPairingKernel.\allowbreak{}mono}}}
\expandafter\newcommand\csname LN@integrableupperMeasure\endcsname{\href{https://github.com/lsernau/grh-lean/blob/72fead102f663d5c7a407964a13dc0cf9d977207/grh-compact/src/New/RealPlace/FloorKernels.lean\#L38-L40}{{\leanfont Level.\allowbreak{}IsPairingKernel.\allowbreak{}integrable\_upperMeasure}}}
\expandafter\newcommand\csname LN@upperPairingeqintegralsucc\endcsname{\href{https://github.com/lsernau/grh-lean/blob/72fead102f663d5c7a407964a13dc0cf9d977207/grh-compact/src/New/RealPlace/FloorKernels.lean\#L42-L48}{{\leanfont Level.\allowbreak{}IsPairingKernel.\allowbreak{}upperPairing\_eq\_integral\_succ}}}
\expandafter\newcommand\csname LN@enormupperPairingle\endcsname{\href{https://github.com/lsernau/grh-lean/blob/72fead102f663d5c7a407964a13dc0cf9d977207/grh-compact/src/New/RealPlace/FloorKernels.lean\#L50-L51}{{\leanfont Level.\allowbreak{}IsPairingKernel.\allowbreak{}enorm\_upperPairing\_le}}}
\expandafter\newcommand\csname LN@upperPairinginvariant\endcsname{\href{https://github.com/lsernau/grh-lean/blob/72fead102f663d5c7a407964a13dc0cf9d977207/grh-compact/src/New/RealPlace/FloorKernels.lean\#L53-L61}{{\leanfont Level.\allowbreak{}IsPairingKernel.\allowbreak{}upperPairing\_invariant}}}
\expandafter\newcommand\csname LN@measurableweightedSquareSum\endcsname{\href{https://github.com/lsernau/grh-lean/blob/72fead102f663d5c7a407964a13dc0cf9d977207/grh-compact/src/New/RealPlace/FloorKernels.lean\#L63-L64}{{\leanfont Level.\allowbreak{}IsPairingKernel.\allowbreak{}measurable\_weightedSquareSum}}}
\expandafter\newcommand\csname LN@weightedSquareinvariantsupport\endcsname{\href{https://github.com/lsernau/grh-lean/blob/72fead102f663d5c7a407964a13dc0cf9d977207/grh-compact/src/New/RealPlace/FloorKernels.lean\#L66-L71}{{\leanfont Level.\allowbreak{}IsPairingKernel.\allowbreak{}weightedSquare\_invariant\_support}}}
\expandafter\newcommand\csname LN@weightedSquareSuminvariantsupport\endcsname{\href{https://github.com/lsernau/grh-lean/blob/72fead102f663d5c7a407964a13dc0cf9d977207/grh-compact/src/New/RealPlace/FloorKernels.lean\#L73-L79}{{\leanfont Level.\allowbreak{}IsPairingKernel.\allowbreak{}weightedSquareSum\_invariant\_support}}}
\expandafter\newcommand\csname LN@pivotModulussqle\endcsname{\href{https://github.com/lsernau/grh-lean/blob/72fead102f663d5c7a407964a13dc0cf9d977207/grh-compact/src/New/RealPlace/FloorKernels.lean\#L84-L91}{{\leanfont Level.\allowbreak{}IsPairingKernel.\allowbreak{}pivotModulus\_sq\_le}}}
\expandafter\newcommand\csname LQ@pivotModulussqle\endcsname{\leanquoteline{0}{lemma pivotModulus\_sq\_le \{u v\} (hK : F.FloorMassLe l K u) (h :}\leanquoteline{2}{∀ σ ∈ GT l, ‖F.upperPairing (l + 1) k (σ * g) (σ * y)‖ₑ \textasciicircum{} 2 ≤ u (σ}\leanquoteline{2}{* g) * v σ := by exact fun \_ \_ ↦ (one\_mul \_).ge) (hu : Measurable u := by fun\_prop) (hv :}\leanquoteline{2}{Measurable v := by fun\_prop) : F.pivotModulus l k g y \textasciicircum{} 2 ≤ K * ∫⁻ σ, v σ ∂F.floorMeasure l}}
\expandafter\newcommand\csname LN@selfNormeqaddsucc\endcsname{\href{https://github.com/lsernau/grh-lean/blob/72fead102f663d5c7a407964a13dc0cf9d977207/grh-compact/src/New/RealPlace/FloorKernels.lean\#L93-L95}{{\leanfont Level.\allowbreak{}IsPairingKernel.\allowbreak{}selfNorm\_eq\_add\_succ}}}
\expandafter\newcommand\csname LN@selfNormleselfNormzero\endcsname{\href{https://github.com/lsernau/grh-lean/blob/72fead102f663d5c7a407964a13dc0cf9d977207/grh-compact/src/New/RealPlace/FloorKernels.lean\#L97-L98}{{\leanfont Level.\allowbreak{}IsPairingKernel.\allowbreak{}selfNorm\_le\_selfNorm\_zero}}}
\expandafter\newcommand\csname LN@Level\endcsname{\href{https://github.com/lsernau/grh-lean/blob/72fead102f663d5c7a407964a13dc0cf9d977207/grh-compact/src/New/RealPlace/LevelGauge.lean\#L25-L35}{{\leanfont Level}}}
\expandafter\newcommand\csname LN@floorMeasure\endcsname{\href{https://github.com/lsernau/grh-lean/blob/72fead102f663d5c7a407964a13dc0cf9d977207/grh-compact/src/New/RealPlace/LevelGauge.lean\#L46-L47}{{\leanfont Level.\allowbreak{}floorMeasure}}}
\expandafter\newcommand\csname LN@stackMeasure\endcsname{\href{https://github.com/lsernau/grh-lean/blob/72fead102f663d5c7a407964a13dc0cf9d977207/grh-compact/src/New/RealPlace/LevelGauge.lean\#L53-L53}{{\leanfont Level.\allowbreak{}stackMeasure}}}
\expandafter\newcommand\csname LN@groupMeasure\endcsname{\href{https://github.com/lsernau/grh-lean/blob/72fead102f663d5c7a407964a13dc0cf9d977207/grh-compact/src/New/RealPlace/LevelGauge.lean\#L57-L57}{{\leanfont Level.\allowbreak{}groupMeasure}}}
\expandafter\newcommand\csname LN@upperMeasure\endcsname{\href{https://github.com/lsernau/grh-lean/blob/72fead102f663d5c7a407964a13dc0cf9d977207/grh-compact/src/New/RealPlace/LevelGauge.lean\#L61-L61}{{\leanfont Level.\allowbreak{}upperMeasure}}}
\expandafter\newcommand\csname LN@IsNear\endcsname{\href{https://github.com/lsernau/grh-lean/blob/72fead102f663d5c7a407964a13dc0cf9d977207/grh-compact/src/New/RealPlace/LevelGauge.lean\#L64-L64}{{\leanfont Level.\allowbreak{}IsNear}}}
\expandafter\newcommand\csname LN@orthonormalPart\endcsname{\href{https://github.com/lsernau/grh-lean/blob/72fead102f663d5c7a407964a13dc0cf9d977207/grh-compact/src/New/RealPlace/LevelGauge.lean\#L68-L69}{{\leanfont Level.\allowbreak{}orthonormalPart}}}
\expandafter\newcommand\csname LN@closePairs\endcsname{\href{https://github.com/lsernau/grh-lean/blob/72fead102f663d5c7a407964a13dc0cf9d977207/grh-compact/src/New/RealPlace/LevelGauge.lean\#L72-L72}{{\leanfont Level.\allowbreak{}closePairs}}}
\expandafter\newcommand\csname LN@floorValues\endcsname{\href{https://github.com/lsernau/grh-lean/blob/72fead102f663d5c7a407964a13dc0cf9d977207/grh-compact/src/New/RealPlace/LevelGauge.lean\#L76-L77}{{\leanfont Level.\allowbreak{}floorValues}}}
\expandafter\newcommand\csname LN@volumeExponent\endcsname{\href{https://github.com/lsernau/grh-lean/blob/72fead102f663d5c7a407964a13dc0cf9d977207/grh-compact/src/New/RealPlace/LevelGauge.lean\#L81-L81}{{\leanfont Level.\allowbreak{}volumeExponent}}}
\expandafter\newcommand\csname LN@floorWeight\endcsname{\href{https://github.com/lsernau/grh-lean/blob/72fead102f663d5c7a407964a13dc0cf9d977207/grh-compact/src/New/RealPlace/LevelGauge.lean\#L85-L85}{{\leanfont Level.\allowbreak{}floorWeight}}}
\expandafter\newcommand\csname LN@floorModulus\endcsname{\href{https://github.com/lsernau/grh-lean/blob/72fead102f663d5c7a407964a13dc0cf9d977207/grh-compact/src/New/RealPlace/LevelGauge.lean\#L89-L89}{{\leanfont Level.\allowbreak{}floorModulus}}}
\expandafter\newcommand\csname LN@measurePreservingorbitPoint\endcsname{\href{https://github.com/lsernau/grh-lean/blob/72fead102f663d5c7a407964a13dc0cf9d977207/grh-compact/src/New/RealPlace/LevelGauge.lean\#L95-L112}{{\leanfont Level.\allowbreak{}measurePreserving\_orbitPoint}}}
\expandafter\newcommand\csname LN@measurePreservingmulVecorbitPoint\endcsname{\href{https://github.com/lsernau/grh-lean/blob/72fead102f663d5c7a407964a13dc0cf9d977207/grh-compact/src/New/RealPlace/LevelGauge.lean\#L116-L120}{{\leanfont Level.\allowbreak{}measurePreserving\_mulVec\_orbitPoint}}}
\expandafter\newcommand\csname LN@aefloorMeasuremem\endcsname{\href{https://github.com/lsernau/grh-lean/blob/72fead102f663d5c7a407964a13dc0cf9d977207/grh-compact/src/New/RealPlace/LevelGauge.lean\#L122-L125}{{\leanfont Level.\allowbreak{}ae\_floorMeasure\_mem}}}
\expandafter\newcommand\csname LN@lintegralfloorMeasurecongr\endcsname{\href{https://github.com/lsernau/grh-lean/blob/72fead102f663d5c7a407964a13dc0cf9d977207/grh-compact/src/New/RealPlace/LevelGauge.lean\#L127-L129}{{\leanfont Level.\allowbreak{}lintegral\_floorMeasure\_congr}}}
\expandafter\newcommand\csname LN@floorMeasuretop\endcsname{\href{https://github.com/lsernau/grh-lean/blob/72fead102f663d5c7a407964a13dc0cf9d977207/grh-compact/src/New/RealPlace/LevelGauge.lean\#L131-L131}{{\leanfont Level.\allowbreak{}floorMeasure\_top}}}
\expandafter\newcommand\csname LN@upperMeasureeqsuccmconv\endcsname{\href{https://github.com/lsernau/grh-lean/blob/72fead102f663d5c7a407964a13dc0cf9d977207/grh-compact/src/New/RealPlace/LevelGauge.lean\#L133-L134}{{\leanfont Level.\allowbreak{}upperMeasure\_eq\_succ\_mconv}}}
\expandafter\newcommand\csname LN@descConvfloorMeasuresucc\endcsname{\href{https://github.com/lsernau/grh-lean/blob/72fead102f663d5c7a407964a13dc0cf9d977207/grh-compact/src/New/RealPlace/LevelGauge.lean\#L136-L140}{{\leanfont Level.\allowbreak{}descConv\_floorMeasure\_succ}}}
\expandafter\newcommand\csname LN@lintegralweylRatiomul\endcsname{\href{https://github.com/lsernau/grh-lean/blob/72fead102f663d5c7a407964a13dc0cf9d977207/grh-compact/src/New/RealPlace/LevelGauge.lean\#L144-L155}{{\leanfont Level.\allowbreak{}lintegral\_weylRatio\_mul}}}
\expandafter\newcommand\csname LN@lintegralfloorMeasuremulright\endcsname{\href{https://github.com/lsernau/grh-lean/blob/72fead102f663d5c7a407964a13dc0cf9d977207/grh-compact/src/New/RealPlace/LevelGauge.lean\#L159-L171}{{\leanfont Level.\allowbreak{}lintegral\_floorMeasure\_mul\_right}}}
\expandafter\newcommand\csname LN@aestackMeasuremem\endcsname{\href{https://github.com/lsernau/grh-lean/blob/72fead102f663d5c7a407964a13dc0cf9d977207/grh-compact/src/New/RealPlace/LevelGauge.lean\#L173-L174}{{\leanfont Level.\allowbreak{}ae\_stackMeasure\_mem}}}
\expandafter\newcommand\csname LN@stackMeasuremconvdirac\endcsname{\href{https://github.com/lsernau/grh-lean/blob/72fead102f663d5c7a407964a13dc0cf9d977207/grh-compact/src/New/RealPlace/LevelGauge.lean\#L178-L186}{{\leanfont Level.\allowbreak{}stackMeasure\_mconv\_dirac}}}
\expandafter\newcommand\csname LQ@stackMeasuremconvdirac\endcsname{\leanquoteline{0}{lemma stackMeasure\_mconv\_dirac (h : l + i = 𝔰.ϱ + 2) (hs : s ∈ GT l) :}\leanquoteline{2}{F.stackMeasure l i ∗ₘ dirac s = F.stackMeasure l i}}
\expandafter\newcommand\csname LN@isMulRightInvariantgroupMeasure\endcsname{\href{https://github.com/lsernau/grh-lean/blob/72fead102f663d5c7a407964a13dc0cf9d977207/grh-compact/src/New/RealPlace/LevelGauge.lean\#L188-L189}{{\leanfont Level.\allowbreak{}isMulRightInvariant\_groupMeasure}}}
\expandafter\newcommand\csname LN@stackMeasurezeroeq\endcsname{\href{https://github.com/lsernau/grh-lean/blob/72fead102f663d5c7a407964a13dc0cf9d977207/grh-compact/src/New/RealPlace/LevelGauge.lean\#L193-L196}{{\leanfont Level.\allowbreak{}stackMeasure\_zero\_eq}}}
\expandafter\newcommand\csname LN@stackMeasurenezero\endcsname{\href{https://github.com/lsernau/grh-lean/blob/72fead102f663d5c7a407964a13dc0cf9d977207/grh-compact/src/New/RealPlace/LevelGauge.lean\#L200-L204}{{\leanfont Level.\allowbreak{}stackMeasure\_ne\_zero}}}
\expandafter\newcommand\csname LN@lintegralmullefteqoffactor\endcsname{\href{https://github.com/lsernau/grh-lean/blob/72fead102f663d5c7a407964a13dc0cf9d977207/grh-compact/src/New/RealPlace/LevelGauge.lean\#L211-L224}{{\leanfont Level.\allowbreak{}lintegral\_mul\_left\_eq\_of\_factor}}}
\expandafter\newcommand\csname LN@lintegralstackMeasuremulleft\endcsname{\href{https://github.com/lsernau/grh-lean/blob/72fead102f663d5c7a407964a13dc0cf9d977207/grh-compact/src/New/RealPlace/LevelGauge.lean\#L228-L232}{{\leanfont Level.\allowbreak{}lintegral\_stackMeasure\_mul\_left}}}
\expandafter\newcommand\csname LN@FloorMassLe\endcsname{\href{https://github.com/lsernau/grh-lean/blob/72fead102f663d5c7a407964a13dc0cf9d977207/grh-compact/src/New/RealPlace/LevelGauge.lean\#L238-L239}{{\leanfont Level.\allowbreak{}FloorMassLe}}}
\expandafter\newcommand\csname LN@upperPairing\endcsname{\href{https://github.com/lsernau/grh-lean/blob/72fead102f663d5c7a407964a13dc0cf9d977207/grh-compact/src/New/RealPlace/LevelGauge.lean\#L243-L243}{{\leanfont Level.\allowbreak{}upperPairing}}}
\expandafter\newcommand\csname LN@IsPairingKernel\endcsname{\href{https://github.com/lsernau/grh-lean/blob/72fead102f663d5c7a407964a13dc0cf9d977207/grh-compact/src/New/RealPlace/LevelGauge.lean\#L248-L255}{{\leanfont Level.\allowbreak{}IsPairingKernel}}}
\expandafter\newcommand\csname LN@invWeight\endcsname{\href{https://github.com/lsernau/grh-lean/blob/72fead102f663d5c7a407964a13dc0cf9d977207/grh-compact/src/New/RealPlace/LevelGauge.lean\#L259-L259}{{\leanfont Level.\allowbreak{}invWeight}}}
\expandafter\newcommand\csname LN@pivotModulus\endcsname{\href{https://github.com/lsernau/grh-lean/blob/72fead102f663d5c7a407964a13dc0cf9d977207/grh-compact/src/New/RealPlace/LevelGauge.lean\#L263-L263}{{\leanfont Level.\allowbreak{}pivotModulus}}}
\expandafter\newcommand\csname LN@weightedSquare\endcsname{\href{https://github.com/lsernau/grh-lean/blob/72fead102f663d5c7a407964a13dc0cf9d977207/grh-compact/src/New/RealPlace/LevelGauge.lean\#L267-L267}{{\leanfont Level.\allowbreak{}weightedSquare}}}
\expandafter\newcommand\csname LN@weightedSquareSum\endcsname{\href{https://github.com/lsernau/grh-lean/blob/72fead102f663d5c7a407964a13dc0cf9d977207/grh-compact/src/New/RealPlace/LevelGauge.lean\#L271-L272}{{\leanfont Level.\allowbreak{}weightedSquareSum}}}
\expandafter\newcommand\csname LN@selfNorm\endcsname{\href{https://github.com/lsernau/grh-lean/blob/72fead102f663d5c7a407964a13dc0cf9d977207/grh-compact/src/New/RealPlace/LevelGauge.lean\#L276-L276}{{\leanfont Level.\allowbreak{}selfNorm}}}
\expandafter\newcommand\csname LN@Gauge\endcsname{\href{https://github.com/lsernau/grh-lean/blob/72fead102f663d5c7a407964a13dc0cf9d977207/grh-compact/src/New/RealPlace/LevelGauge.lean\#L283-L301}{{\leanfont Gauge}}}
\expandafter\newcommand\csname LN@oneleLambda\endcsname{\href{https://github.com/lsernau/grh-lean/blob/72fead102f663d5c7a407964a13dc0cf9d977207/grh-compact/src/New/RealPlace/LevelGauge.lean\#L303-L304}{{\leanfont Gauge.\allowbreak{}one\_le\_Lambda}}}
\expandafter\newcommand\csname LN@positionsOfSize\endcsname{\href{https://github.com/lsernau/grh-lean/blob/72fead102f663d5c7a407964a13dc0cf9d977207/grh-compact/src/New/RealPlace/LevelGauge.lean\#L307-L308}{{\leanfont Gauge.\allowbreak{}positionsOfSize}}}
\expandafter\newcommand\csname LN@Size\endcsname{\href{https://github.com/lsernau/grh-lean/blob/72fead102f663d5c7a407964a13dc0cf9d977207/grh-compact/src/New/RealPlace/Size.lean\#L15-L22}{{\leanfont Size}}}
\expandafter\newcommand\csname LN@weylRatio\endcsname{\href{https://github.com/lsernau/grh-lean/blob/72fead102f663d5c7a407964a13dc0cf9d977207/grh-compact/src/New/RealPlace/Size.lean\#L29-L29}{{\leanfont weylRatio}}}
\expandafter\newcommand\csname LN@coeweylRatio\endcsname{\href{https://github.com/lsernau/grh-lean/blob/72fead102f663d5c7a407964a13dc0cf9d977207/grh-compact/src/New/RealPlace/Size.lean\#L31-L34}{{\leanfont coe\_weylRatio}}}
\expandafter\newcommand\csname LN@weylRatioinv\endcsname{\href{https://github.com/lsernau/grh-lean/blob/72fead102f663d5c7a407964a13dc0cf9d977207/grh-compact/src/New/RealPlace/Size.lean\#L36-L36}{{\leanfont weylRatio\_inv}}}
\expandafter\newcommand\csname LN@cardsucclefloorOf\endcsname{\href{https://github.com/lsernau/grh-lean/blob/72fead102f663d5c7a407964a13dc0cf9d977207/grh-compact/src/New/RealPlace/Size.lean\#L41-L42}{{\leanfont card\_succ\_le\_floorOf}}}
\expandafter\newcommand\csname LN@vecsOfSize\endcsname{\href{https://github.com/lsernau/grh-lean/blob/72fead102f663d5c7a407964a13dc0cf9d977207/grh-compact/src/New/RealPlace/Size.lean\#L47-L47}{{\leanfont Size.\allowbreak{}vecsOfSize}}}
\expandafter\newcommand\csname LN@HasMatSize\endcsname{\href{https://github.com/lsernau/grh-lean/blob/72fead102f663d5c7a407964a13dc0cf9d977207/grh-compact/src/New/RealPlace/Size.lean\#L50-L50}{{\leanfont Size.\allowbreak{}HasMatSize}}}
\expandafter\newcommand\csname LN@IsHead\endcsname{\href{https://github.com/lsernau/grh-lean/blob/72fead102f663d5c7a407964a13dc0cf9d977207/grh-compact/src/New/RealPlace/Size.lean\#L54-L54}{{\leanfont Size.\allowbreak{}IsHead}}}
\expandafter\newcommand\csname LN@starhead\endcsname{\href{https://github.com/lsernau/grh-lean/blob/72fead102f663d5c7a407964a13dc0cf9d977207/grh-compact/src/New/RealPlace/Size.lean\#L59-L59}{{\leanfont Size.\allowbreak{}IsHead.\allowbreak{}star\_head}}}
\expandafter\newcommand\csname LN@memofSizemono\endcsname{\href{https://github.com/lsernau/grh-lean/blob/72fead102f663d5c7a407964a13dc0cf9d977207/grh-compact/src/New/RealPlace/Size.lean\#L61-L62}{{\leanfont Size.\allowbreak{}mem\_ofSize\_mono}}}
\expandafter\newcommand\csname LN@submemofSize\endcsname{\href{https://github.com/lsernau/grh-lean/blob/72fead102f663d5c7a407964a13dc0cf9d977207/grh-compact/src/New/RealPlace/Size.lean\#L64-L65}{{\leanfont Size.\allowbreak{}sub\_mem\_ofSize}}}
\expandafter\newcommand\csname LN@monoa\endcsname{\href{https://github.com/lsernau/grh-lean/blob/72fead102f663d5c7a407964a13dc0cf9d977207/grh-compact/src/New/RealPlace/Size.lean\#L67-L68}{{\leanfont Size.\allowbreak{}HasMatSize.\allowbreak{}mono}}}
\expandafter\newcommand\csname LN@one\endcsname{\href{https://github.com/lsernau/grh-lean/blob/72fead102f663d5c7a407964a13dc0cf9d977207/grh-compact/src/New/RealPlace/Size.lean\#L70-L70}{{\leanfont Size.\allowbreak{}HasMatSize.\allowbreak{}one}}}
\expandafter\newcommand\csname LN@sub\endcsname{\href{https://github.com/lsernau/grh-lean/blob/72fead102f663d5c7a407964a13dc0cf9d977207/grh-compact/src/New/RealPlace/Size.lean\#L72-L73}{{\leanfont Size.\allowbreak{}HasMatSize.\allowbreak{}sub}}}
\expandafter\newcommand\csname LN@dotProductmemofSize\endcsname{\href{https://github.com/lsernau/grh-lean/blob/72fead102f663d5c7a407964a13dc0cf9d977207/grh-compact/src/New/RealPlace/Size.lean\#L75-L77}{{\leanfont Size.\allowbreak{}dotProduct\_mem\_ofSize}}}
\expandafter\newcommand\csname LN@mulVec\endcsname{\href{https://github.com/lsernau/grh-lean/blob/72fead102f663d5c7a407964a13dc0cf9d977207/grh-compact/src/New/RealPlace/Size.lean\#L79-L80}{{\leanfont Size.\allowbreak{}HasMatSize.\allowbreak{}mulVec}}}
\expandafter\newcommand\csname LN@mul\endcsname{\href{https://github.com/lsernau/grh-lean/blob/72fead102f663d5c7a407964a13dc0cf9d977207/grh-compact/src/New/RealPlace/Size.lean\#L82-L83}{{\leanfont Size.\allowbreak{}HasMatSize.\allowbreak{}mul}}}
\expandafter\newcommand\csname LN@detmemofSize\endcsname{\href{https://github.com/lsernau/grh-lean/blob/72fead102f663d5c7a407964a13dc0cf9d977207/grh-compact/src/New/RealPlace/Size.lean\#L85-L88}{{\leanfont Size.\allowbreak{}det\_mem\_ofSize}}}
\expandafter\newcommand\csname LN@adjointapplymemofSize\endcsname{\href{https://github.com/lsernau/grh-lean/blob/72fead102f663d5c7a407964a13dc0cf9d977207/grh-compact/src/New/RealPlace/Size.lean\#L92-L97}{{\leanfont Size.\allowbreak{}adjoint\_apply\_mem\_ofSize}}}
\expandafter\newcommand\csname LN@inv\endcsname{\href{https://github.com/lsernau/grh-lean/blob/72fead102f663d5c7a407964a13dc0cf9d977207/grh-compact/src/New/RealPlace/Size.lean\#L99-L101}{{\leanfont Size.\allowbreak{}HasMatSize.\allowbreak{}inv}}}
\expandafter\newcommand\csname LN@rowmemofSizeofdualframe\endcsname{\href{https://github.com/lsernau/grh-lean/blob/72fead102f663d5c7a407964a13dc0cf9d977207/grh-compact/src/New/RealPlace/Size.lean\#L103-L110}{{\leanfont Size.\allowbreak{}row\_mem\_ofSize\_of\_dual\_frame}}}
\expandafter\newcommand\csname LN@hasMatSizeonesubinv\endcsname{\href{https://github.com/lsernau/grh-lean/blob/72fead102f663d5c7a407964a13dc0cf9d977207/grh-compact/src/New/RealPlace/Size.lean\#L115-L125}{{\leanfont Size.\allowbreak{}hasMatSize\_one\_sub\_inv}}}
\expandafter\newcommand\csname LN@minormuldualmemofSize\endcsname{\href{https://github.com/lsernau/grh-lean/blob/72fead102f663d5c7a407964a13dc0cf9d977207/grh-compact/src/New/RealPlace/Size.lean\#L133-L146}{{\leanfont Size.\allowbreak{}minor\_mul\_dual\_mem\_ofSize}}}
\expandafter\newcommand\csname LN@minorsmemofSize\endcsname{\href{https://github.com/lsernau/grh-lean/blob/72fead102f663d5c7a407964a13dc0cf9d977207/grh-compact/src/New/RealPlace/Size.lean\#L150-L154}{{\leanfont Size.\allowbreak{}minors\_mem\_ofSize}}}
\expandafter\newcommand\csname LN@headentriesmemofSize\endcsname{\href{https://github.com/lsernau/grh-lean/blob/72fead102f663d5c7a407964a13dc0cf9d977207/grh-compact/src/New/RealPlace/Size.lean\#L158-L164}{{\leanfont Size.\allowbreak{}head\_entries\_mem\_ofSize}}}
\expandafter\newcommand\csname LN@hasMatSizeofisHead\endcsname{\href{https://github.com/lsernau/grh-lean/blob/72fead102f663d5c7a407964a13dc0cf9d977207/grh-compact/src/New/RealPlace/Size.lean\#L168-L174}{{\leanfont Size.\allowbreak{}hasMatSize\_of\_isHead}}}
\expandafter\newcommand\csname LN@floorMassLeofchart\endcsname{\href{https://github.com/lsernau/grh-lean/blob/72fead102f663d5c7a407964a13dc0cf9d977207/grh-compact/src/New/RealPlace/WeightPlancherel.lean\#L20-L31}{{\leanfont Gauge.\allowbreak{}floorMassLe\_of\_chart}}}
\expandafter\newcommand\csname LN@floorMassLefloorWeightdiv\endcsname{\href{https://github.com/lsernau/grh-lean/blob/72fead102f663d5c7a407964a13dc0cf9d977207/grh-compact/src/New/RealPlace/WeightPlancherel.lean\#L35-L45}{{\leanfont Gauge.\allowbreak{}floorMassLe\_floorWeight\_div}}}
\expandafter\newcommand\csname LN@floorMassLeone\endcsname{\href{https://github.com/lsernau/grh-lean/blob/72fead102f663d5c7a407964a13dc0cf9d977207/grh-compact/src/New/RealPlace/WeightPlancherel.lean\#L49-L52}{{\leanfont Gauge.\allowbreak{}floorMassLe\_one}}}
\expandafter\newcommand\csname LN@floorConst\endcsname{\href{https://github.com/lsernau/grh-lean/blob/72fead102f663d5c7a407964a13dc0cf9d977207/grh-compact/src/New/RealPlace/WeightPlancherel.lean\#L56-L56}{{\leanfont Gauge.\allowbreak{}floorConst}}}
\expandafter\newcommand\csname LN@lintegralfloorWeightle\endcsname{\href{https://github.com/lsernau/grh-lean/blob/72fead102f663d5c7a407964a13dc0cf9d977207/grh-compact/src/New/RealPlace/WeightPlancherel.lean\#L62-L90}{{\leanfont Gauge.\allowbreak{}lintegral\_floorWeight\_le}}}
\expandafter\newcommand\csname LN@globalConst\endcsname{\href{https://github.com/lsernau/grh-lean/blob/72fead102f663d5c7a407964a13dc0cf9d977207/grh-compact/src/New/RealPlace/WeightPlancherel.lean\#L93-L93}{{\leanfont Gauge.\allowbreak{}globalConst}}}
\expandafter\newcommand\csname LN@lintegralmulinvWeightle\endcsname{\href{https://github.com/lsernau/grh-lean/blob/72fead102f663d5c7a407964a13dc0cf9d977207/grh-compact/src/New/RealPlace/WeightPlancherel.lean\#L98-L135}{{\leanfont Gauge.\allowbreak{}lintegral\_mul\_invWeight\_le}}}
\expandafter\newcommand\csname LN@amplitudesupportandplancherel\endcsname{\href{https://github.com/lsernau/grh-lean/blob/72fead102f663d5c7a407964a13dc0cf9d977207/grh-compact/src/New/RealPlace/WeightPlancherel.lean\#L141-L188}{{\leanfont Gauge.\allowbreak{}amplitude\_support\_and\_plancherel}}}
\expandafter\newcommand\csname LN@parsevalheisenbergMeasureofbox\endcsname{\href{https://github.com/lsernau/grh-lean/blob/72fead102f663d5c7a407964a13dc0cf9d977207/grh-compact/src/New/RealPlace/WeightPlancherel.lean\#L194-L220}{{\leanfont Gauge.\allowbreak{}parseval\_heisenbergMeasure\_of\_box}}}
\newcommand{\ConcordanceRows}{\crow{memvecsOfSizeofheisenbergElem}{\href{https://github.com/lsernau/grh-lean/blob/72fead102f663d5c7a407964a13dc0cf9d977207/grh-compact/src/New/RealPlace/FlatCharts.lean\#L21-L26}{{\leanfont Size.\allowbreak{}mem\_vecsOfSize\_of\_heisenbergElem}}}{FlatCharts.lean}{21--26}
\crow{isNearmuliff}{\href{https://github.com/lsernau/grh-lean/blob/72fead102f663d5c7a407964a13dc0cf9d977207/grh-compact/src/New/RealPlace/FlatCharts.lean\#L28-L29}{{\leanfont Level.\allowbreak{}isNear\_mul\_iff}}}{FlatCharts.lean}{28--29}
\crow{offloorGroupmul}{\href{https://github.com/lsernau/grh-lean/blob/72fead102f663d5c7a407964a13dc0cf9d977207/grh-compact/src/New/RealPlace/FlatCharts.lean\#L31-L33}{{\leanfont Level.\allowbreak{}IsNear.\allowbreak{}of\_floorGroup\_mul}}}{FlatCharts.lean}{31--33}
\crow{measurableSetisNear}{\href{https://github.com/lsernau/grh-lean/blob/72fead102f663d5c7a407964a13dc0cf9d977207/grh-compact/src/New/RealPlace/FlatCharts.lean\#L35-L37}{{\leanfont Level.\allowbreak{}measurableSet\_isNear}}}{FlatCharts.lean}{35--37}
\crow{upperMeasureisNearle}{\href{https://github.com/lsernau/grh-lean/blob/72fead102f663d5c7a407964a13dc0cf9d977207/grh-compact/src/New/RealPlace/FlatCharts.lean\#L41-L51}{{\leanfont Level.\allowbreak{}upperMeasure\_isNear\_le}}}{FlatCharts.lean}{41--51}
\crow{isOpenfloorValues}{\href{https://github.com/lsernau/grh-lean/blob/72fead102f663d5c7a407964a13dc0cf9d977207/grh-compact/src/New/RealPlace/FlatCharts.lean\#L53-L55}{{\leanfont Level.\allowbreak{}isOpen\_floorValues}}}{FlatCharts.lean}{53--55}
\crow{floorValuesmul}{\href{https://github.com/lsernau/grh-lean/blob/72fead102f663d5c7a407964a13dc0cf9d977207/grh-compact/src/New/RealPlace/FlatCharts.lean\#L57-L61}{{\leanfont Level.\allowbreak{}floorValues\_mul}}}{FlatCharts.lean}{57--61}
\crow{volumefloorValuesnezeronetop}{\href{https://github.com/lsernau/grh-lean/blob/72fead102f663d5c7a407964a13dc0cf9d977207/grh-compact/src/New/RealPlace/FlatCharts.lean\#L63-L66}{{\leanfont Level.\allowbreak{}volume\_floorValues\_ne\_zero\_ne\_top}}}{FlatCharts.lean}{63--66}
\crow{floorWeightsucc}{\href{https://github.com/lsernau/grh-lean/blob/72fead102f663d5c7a407964a13dc0cf9d977207/grh-compact/src/New/RealPlace/FlatCharts.lean\#L68-L70}{{\leanfont Level.\allowbreak{}floorWeight\_succ}}}{FlatCharts.lean}{68--70}
\crow{floorWeightnezeroandnetop}{\href{https://github.com/lsernau/grh-lean/blob/72fead102f663d5c7a407964a13dc0cf9d977207/grh-compact/src/New/RealPlace/FlatCharts.lean\#L72-L73}{{\leanfont Level.\allowbreak{}floorWeight\_ne\_zero\_and\_ne\_top}}}{FlatCharts.lean}{72--73}
\crow{floorWeightmulleft}{\href{https://github.com/lsernau/grh-lean/blob/72fead102f663d5c7a407964a13dc0cf9d977207/grh-compact/src/New/RealPlace/FlatCharts.lean\#L76-L77}{{\leanfont Level.\allowbreak{}floorWeight\_mul\_left}}}{FlatCharts.lean}{76--77}
\crow{invWeightmul}{\href{https://github.com/lsernau/grh-lean/blob/72fead102f663d5c7a407964a13dc0cf9d977207/grh-compact/src/New/RealPlace/FlatCharts.lean\#L80-L83}{{\leanfont Level.\allowbreak{}invWeight\_mul}}}{FlatCharts.lean}{80--83}
\crow{measurablefloorWeight}{\href{https://github.com/lsernau/grh-lean/blob/72fead102f663d5c7a407964a13dc0cf9d977207/grh-compact/src/New/RealPlace/FlatCharts.lean\#L84-L85}{{\leanfont measurable\_floorWeight}}}{FlatCharts.lean}{84--85}
\crow{measurableinvWeight}{\href{https://github.com/lsernau/grh-lean/blob/72fead102f663d5c7a407964a13dc0cf9d977207/grh-compact/src/New/RealPlace/FlatCharts.lean\#L87-L88}{{\leanfont Level.\allowbreak{}measurable\_invWeight}}}{FlatCharts.lean}{87--88}
\crow{measurableSetvecsOfSize}{\href{https://github.com/lsernau/grh-lean/blob/72fead102f663d5c7a407964a13dc0cf9d977207/grh-compact/src/New/RealPlace/FlatCharts.lean\#L92-L93}{{\leanfont Gauge.\allowbreak{}measurableSet\_vecsOfSize}}}{FlatCharts.lean}{92--93}
\crow{memvecsOfSizeoffloorBasismulVec}{\href{https://github.com/lsernau/grh-lean/blob/72fead102f663d5c7a407964a13dc0cf9d977207/grh-compact/src/New/RealPlace/FlatCharts.lean\#L95-L99}{{\leanfont Gauge.\allowbreak{}mem\_vecsOfSize\_of\_floorBasis\_mulVec}}}{FlatCharts.lean}{95--99}
\crow{HasFibreBound}{\href{https://github.com/lsernau/grh-lean/blob/72fead102f663d5c7a407964a13dc0cf9d977207/grh-compact/src/New/RealPlace/FlatCharts.lean\#L103-L106}{{\leanfont Gauge.\allowbreak{}HasFibreBound}}}{FlatCharts.lean}{103--106}
\crow{setLIntegralleofhasFibreBound}{\href{https://github.com/lsernau/grh-lean/blob/72fead102f663d5c7a407964a13dc0cf9d977207/grh-compact/src/New/RealPlace/FlatCharts.lean\#L108-L118}{{\leanfont Gauge.\allowbreak{}setLIntegral\_le\_of\_hasFibreBound}}}{FlatCharts.lean}{108--118}
\crow{hasFibreBoundfloorQuadricofle}{\href{https://github.com/lsernau/grh-lean/blob/72fead102f663d5c7a407964a13dc0cf9d977207/grh-compact/src/New/RealPlace/FlatCharts.lean\#L123-L131}{{\leanfont Gauge.\allowbreak{}hasFibreBound\_floorQuadric\_of\_le}}}{FlatCharts.lean}{123--131}
\crow{hasFibreBoundfloorQuadric}{\href{https://github.com/lsernau/grh-lean/blob/72fead102f663d5c7a407964a13dc0cf9d977207/grh-compact/src/New/RealPlace/FlatCharts.lean\#L135-L136}{{\leanfont Gauge.\allowbreak{}hasFibreBound\_floorQuadric}}}{FlatCharts.lean}{135--136}
\crow{hasFibreBoundfloorVolume}{\href{https://github.com/lsernau/grh-lean/blob/72fead102f663d5c7a407964a13dc0cf9d977207/grh-compact/src/New/RealPlace/FlatCharts.lean\#L141-L148}{{\leanfont Gauge.\allowbreak{}hasFibreBound\_floorVolume}}}{FlatCharts.lean}{141--148}
\crow{lintegralheisenbergMeasurele}{\href{https://github.com/lsernau/grh-lean/blob/72fead102f663d5c7a407964a13dc0cf9d977207/grh-compact/src/New/RealPlace/FlatCharts.lean\#L153-L173}{{\leanfont Gauge.\allowbreak{}lintegral\_heisenbergMeasure\_le}}}{FlatCharts.lean}{153--173}
\crow{setLIntegralorthonormalPartle}{\href{https://github.com/lsernau/grh-lean/blob/72fead102f663d5c7a407964a13dc0cf9d977207/grh-compact/src/New/RealPlace/FlatCharts.lean\#L175-L181}{{\leanfont Gauge.\allowbreak{}setLIntegral\_orthonormalPart\_le}}}{FlatCharts.lean}{175--181}
\crow{measurableSetpositionsOfSize}{\href{https://github.com/lsernau/grh-lean/blob/72fead102f663d5c7a407964a13dc0cf9d977207/grh-compact/src/New/RealPlace/FlatCharts.lean\#L183-L184}{{\leanfont Gauge.\allowbreak{}measurableSet\_positionsOfSize}}}{FlatCharts.lean}{183--184}
\crow{conjmulinvclosetoone}{\href{https://github.com/lsernau/grh-lean/blob/72fead102f663d5c7a407964a13dc0cf9d977207/grh-compact/src/New/RealPlace/FlatCharts.lean\#L188-L192}{{\leanfont Gauge.\allowbreak{}conj\_mul\_inv\_close\_to\_one}}}{FlatCharts.lean}{188--192}
\crow{orbitPointmemvecsOfSize}{\href{https://github.com/lsernau/grh-lean/blob/72fead102f663d5c7a407964a13dc0cf9d977207/grh-compact/src/New/RealPlace/FlatCharts.lean\#L196-L199}{{\leanfont Gauge.\allowbreak{}orbitPoint\_mem\_vecsOfSize}}}{FlatCharts.lean}{196--199}
\crow{existshasMatSizeofisNear}{\href{https://github.com/lsernau/grh-lean/blob/72fead102f663d5c7a407964a13dc0cf9d977207/grh-compact/src/New/RealPlace/FlatCharts.lean\#L203-L209}{{\leanfont Gauge.\allowbreak{}exists\_hasMatSize\_of\_isNear}}}{FlatCharts.lean}{203--209}
\crow{hasMatSizesuboneoffloorValues}{\href{https://github.com/lsernau/grh-lean/blob/72fead102f663d5c7a407964a13dc0cf9d977207/grh-compact/src/New/RealPlace/FlatCharts.lean\#L215-L233}{{\leanfont Gauge.\allowbreak{}hasMatSize\_sub\_one\_of\_floorValues}}}{FlatCharts.lean}{215--233}
\crow{existsnezeroofpivotModulus}{\href{https://github.com/lsernau/grh-lean/blob/72fead102f663d5c7a407964a13dc0cf9d977207/grh-compact/src/New/RealPlace/FloorKernels.lean\#L18-L21}{{\leanfont Level.\allowbreak{}exists\_ne\_zero\_of\_pivotModulus}}}{FloorKernels.lean}{18--21}
\crow{lintegralupperPairingsqle}{\href{https://github.com/lsernau/grh-lean/blob/72fead102f663d5c7a407964a13dc0cf9d977207/grh-compact/src/New/RealPlace/FloorKernels.lean\#L23-L26}{{\leanfont Level.\allowbreak{}lintegral\_upperPairing\_sq\_le}}}{FloorKernels.lean}{23--26}
\crow{stronglyMeasurableupperPairing}{\href{https://github.com/lsernau/grh-lean/blob/72fead102f663d5c7a407964a13dc0cf9d977207/grh-compact/src/New/RealPlace/FloorKernels.lean\#L32-L33}{{\leanfont Level.\allowbreak{}IsPairingKernel.\allowbreak{}stronglyMeasurable\_upperPairing}}}{FloorKernels.lean}{32--33}
\crow{mono}{\href{https://github.com/lsernau/grh-lean/blob/72fead102f663d5c7a407964a13dc0cf9d977207/grh-compact/src/New/RealPlace/FloorKernels.lean\#L35-L36}{{\leanfont Level.\allowbreak{}IsPairingKernel.\allowbreak{}mono}}}{FloorKernels.lean}{35--36}
\crow{integrableupperMeasure}{\href{https://github.com/lsernau/grh-lean/blob/72fead102f663d5c7a407964a13dc0cf9d977207/grh-compact/src/New/RealPlace/FloorKernels.lean\#L38-L40}{{\leanfont Level.\allowbreak{}IsPairingKernel.\allowbreak{}integrable\_upperMeasure}}}{FloorKernels.lean}{38--40}
\crow{upperPairingeqintegralsucc}{\href{https://github.com/lsernau/grh-lean/blob/72fead102f663d5c7a407964a13dc0cf9d977207/grh-compact/src/New/RealPlace/FloorKernels.lean\#L42-L48}{{\leanfont Level.\allowbreak{}IsPairingKernel.\allowbreak{}upperPairing\_eq\_integral\_succ}}}{FloorKernels.lean}{42--48}
\crow{enormupperPairingle}{\href{https://github.com/lsernau/grh-lean/blob/72fead102f663d5c7a407964a13dc0cf9d977207/grh-compact/src/New/RealPlace/FloorKernels.lean\#L50-L51}{{\leanfont Level.\allowbreak{}IsPairingKernel.\allowbreak{}enorm\_upperPairing\_le}}}{FloorKernels.lean}{50--51}
\crow{upperPairinginvariant}{\href{https://github.com/lsernau/grh-lean/blob/72fead102f663d5c7a407964a13dc0cf9d977207/grh-compact/src/New/RealPlace/FloorKernels.lean\#L53-L61}{{\leanfont Level.\allowbreak{}IsPairingKernel.\allowbreak{}upperPairing\_invariant}}}{FloorKernels.lean}{53--61}
\crow{measurableweightedSquareSum}{\href{https://github.com/lsernau/grh-lean/blob/72fead102f663d5c7a407964a13dc0cf9d977207/grh-compact/src/New/RealPlace/FloorKernels.lean\#L63-L64}{{\leanfont Level.\allowbreak{}IsPairingKernel.\allowbreak{}measurable\_weightedSquareSum}}}{FloorKernels.lean}{63--64}
\crow{weightedSquareinvariantsupport}{\href{https://github.com/lsernau/grh-lean/blob/72fead102f663d5c7a407964a13dc0cf9d977207/grh-compact/src/New/RealPlace/FloorKernels.lean\#L66-L71}{{\leanfont Level.\allowbreak{}IsPairingKernel.\allowbreak{}weightedSquare\_invariant\_support}}}{FloorKernels.lean}{66--71}
\crow{weightedSquareSuminvariantsupport}{\href{https://github.com/lsernau/grh-lean/blob/72fead102f663d5c7a407964a13dc0cf9d977207/grh-compact/src/New/RealPlace/FloorKernels.lean\#L73-L79}{{\leanfont Level.\allowbreak{}IsPairingKernel.\allowbreak{}weightedSquareSum\_invariant\_support}}}{FloorKernels.lean}{73--79}
\crow{pivotModulussqle}{\href{https://github.com/lsernau/grh-lean/blob/72fead102f663d5c7a407964a13dc0cf9d977207/grh-compact/src/New/RealPlace/FloorKernels.lean\#L84-L91}{{\leanfont Level.\allowbreak{}IsPairingKernel.\allowbreak{}pivotModulus\_sq\_le}}}{FloorKernels.lean}{84--91}
\crow{selfNormeqaddsucc}{\href{https://github.com/lsernau/grh-lean/blob/72fead102f663d5c7a407964a13dc0cf9d977207/grh-compact/src/New/RealPlace/FloorKernels.lean\#L93-L95}{{\leanfont Level.\allowbreak{}IsPairingKernel.\allowbreak{}selfNorm\_eq\_add\_succ}}}{FloorKernels.lean}{93--95}
\crow{selfNormleselfNormzero}{\href{https://github.com/lsernau/grh-lean/blob/72fead102f663d5c7a407964a13dc0cf9d977207/grh-compact/src/New/RealPlace/FloorKernels.lean\#L97-L98}{{\leanfont Level.\allowbreak{}IsPairingKernel.\allowbreak{}selfNorm\_le\_selfNorm\_zero}}}{FloorKernels.lean}{97--98}
\crow{Level}{\href{https://github.com/lsernau/grh-lean/blob/72fead102f663d5c7a407964a13dc0cf9d977207/grh-compact/src/New/RealPlace/LevelGauge.lean\#L25-L35}{{\leanfont Level}}}{LevelGauge.lean}{25--35}
\crow{floorMeasure}{\href{https://github.com/lsernau/grh-lean/blob/72fead102f663d5c7a407964a13dc0cf9d977207/grh-compact/src/New/RealPlace/LevelGauge.lean\#L46-L47}{{\leanfont Level.\allowbreak{}floorMeasure}}}{LevelGauge.lean}{46--47}
\crow{stackMeasure}{\href{https://github.com/lsernau/grh-lean/blob/72fead102f663d5c7a407964a13dc0cf9d977207/grh-compact/src/New/RealPlace/LevelGauge.lean\#L53-L53}{{\leanfont Level.\allowbreak{}stackMeasure}}}{LevelGauge.lean}{53--53}
\crow{groupMeasure}{\href{https://github.com/lsernau/grh-lean/blob/72fead102f663d5c7a407964a13dc0cf9d977207/grh-compact/src/New/RealPlace/LevelGauge.lean\#L57-L57}{{\leanfont Level.\allowbreak{}groupMeasure}}}{LevelGauge.lean}{57--57}
\crow{upperMeasure}{\href{https://github.com/lsernau/grh-lean/blob/72fead102f663d5c7a407964a13dc0cf9d977207/grh-compact/src/New/RealPlace/LevelGauge.lean\#L61-L61}{{\leanfont Level.\allowbreak{}upperMeasure}}}{LevelGauge.lean}{61--61}
\crow{IsNear}{\href{https://github.com/lsernau/grh-lean/blob/72fead102f663d5c7a407964a13dc0cf9d977207/grh-compact/src/New/RealPlace/LevelGauge.lean\#L64-L64}{{\leanfont Level.\allowbreak{}IsNear}}}{LevelGauge.lean}{64--64}
\crow{orthonormalPart}{\href{https://github.com/lsernau/grh-lean/blob/72fead102f663d5c7a407964a13dc0cf9d977207/grh-compact/src/New/RealPlace/LevelGauge.lean\#L68-L69}{{\leanfont Level.\allowbreak{}orthonormalPart}}}{LevelGauge.lean}{68--69}
\crow{closePairs}{\href{https://github.com/lsernau/grh-lean/blob/72fead102f663d5c7a407964a13dc0cf9d977207/grh-compact/src/New/RealPlace/LevelGauge.lean\#L72-L72}{{\leanfont Level.\allowbreak{}closePairs}}}{LevelGauge.lean}{72--72}
\crow{floorValues}{\href{https://github.com/lsernau/grh-lean/blob/72fead102f663d5c7a407964a13dc0cf9d977207/grh-compact/src/New/RealPlace/LevelGauge.lean\#L76-L77}{{\leanfont Level.\allowbreak{}floorValues}}}{LevelGauge.lean}{76--77}
\crow{volumeExponent}{\href{https://github.com/lsernau/grh-lean/blob/72fead102f663d5c7a407964a13dc0cf9d977207/grh-compact/src/New/RealPlace/LevelGauge.lean\#L81-L81}{{\leanfont Level.\allowbreak{}volumeExponent}}}{LevelGauge.lean}{81--81}
\crow{floorWeight}{\href{https://github.com/lsernau/grh-lean/blob/72fead102f663d5c7a407964a13dc0cf9d977207/grh-compact/src/New/RealPlace/LevelGauge.lean\#L85-L85}{{\leanfont Level.\allowbreak{}floorWeight}}}{LevelGauge.lean}{85--85}
\crow{floorModulus}{\href{https://github.com/lsernau/grh-lean/blob/72fead102f663d5c7a407964a13dc0cf9d977207/grh-compact/src/New/RealPlace/LevelGauge.lean\#L89-L89}{{\leanfont Level.\allowbreak{}floorModulus}}}{LevelGauge.lean}{89--89}
\crow{measurePreservingorbitPoint}{\href{https://github.com/lsernau/grh-lean/blob/72fead102f663d5c7a407964a13dc0cf9d977207/grh-compact/src/New/RealPlace/LevelGauge.lean\#L95-L112}{{\leanfont Level.\allowbreak{}measurePreserving\_orbitPoint}}}{LevelGauge.lean}{95--112}
\crow{measurePreservingmulVecorbitPoint}{\href{https://github.com/lsernau/grh-lean/blob/72fead102f663d5c7a407964a13dc0cf9d977207/grh-compact/src/New/RealPlace/LevelGauge.lean\#L116-L120}{{\leanfont Level.\allowbreak{}measurePreserving\_mulVec\_orbitPoint}}}{LevelGauge.lean}{116--120}
\crow{aefloorMeasuremem}{\href{https://github.com/lsernau/grh-lean/blob/72fead102f663d5c7a407964a13dc0cf9d977207/grh-compact/src/New/RealPlace/LevelGauge.lean\#L122-L125}{{\leanfont Level.\allowbreak{}ae\_floorMeasure\_mem}}}{LevelGauge.lean}{122--125}
\crow{lintegralfloorMeasurecongr}{\href{https://github.com/lsernau/grh-lean/blob/72fead102f663d5c7a407964a13dc0cf9d977207/grh-compact/src/New/RealPlace/LevelGauge.lean\#L127-L129}{{\leanfont Level.\allowbreak{}lintegral\_floorMeasure\_congr}}}{LevelGauge.lean}{127--129}
\crow{floorMeasuretop}{\href{https://github.com/lsernau/grh-lean/blob/72fead102f663d5c7a407964a13dc0cf9d977207/grh-compact/src/New/RealPlace/LevelGauge.lean\#L131-L131}{{\leanfont Level.\allowbreak{}floorMeasure\_top}}}{LevelGauge.lean}{131--131}
\crow{upperMeasureeqsuccmconv}{\href{https://github.com/lsernau/grh-lean/blob/72fead102f663d5c7a407964a13dc0cf9d977207/grh-compact/src/New/RealPlace/LevelGauge.lean\#L133-L134}{{\leanfont Level.\allowbreak{}upperMeasure\_eq\_succ\_mconv}}}{LevelGauge.lean}{133--134}
\crow{descConvfloorMeasuresucc}{\href{https://github.com/lsernau/grh-lean/blob/72fead102f663d5c7a407964a13dc0cf9d977207/grh-compact/src/New/RealPlace/LevelGauge.lean\#L136-L140}{{\leanfont Level.\allowbreak{}descConv\_floorMeasure\_succ}}}{LevelGauge.lean}{136--140}
\crow{lintegralweylRatiomul}{\href{https://github.com/lsernau/grh-lean/blob/72fead102f663d5c7a407964a13dc0cf9d977207/grh-compact/src/New/RealPlace/LevelGauge.lean\#L144-L155}{{\leanfont Level.\allowbreak{}lintegral\_weylRatio\_mul}}}{LevelGauge.lean}{144--155}
\crow{lintegralfloorMeasuremulright}{\href{https://github.com/lsernau/grh-lean/blob/72fead102f663d5c7a407964a13dc0cf9d977207/grh-compact/src/New/RealPlace/LevelGauge.lean\#L159-L171}{{\leanfont Level.\allowbreak{}lintegral\_floorMeasure\_mul\_right}}}{LevelGauge.lean}{159--171}
\crow{aestackMeasuremem}{\href{https://github.com/lsernau/grh-lean/blob/72fead102f663d5c7a407964a13dc0cf9d977207/grh-compact/src/New/RealPlace/LevelGauge.lean\#L173-L174}{{\leanfont Level.\allowbreak{}ae\_stackMeasure\_mem}}}{LevelGauge.lean}{173--174}
\crow{stackMeasuremconvdirac}{\href{https://github.com/lsernau/grh-lean/blob/72fead102f663d5c7a407964a13dc0cf9d977207/grh-compact/src/New/RealPlace/LevelGauge.lean\#L178-L186}{{\leanfont Level.\allowbreak{}stackMeasure\_mconv\_dirac}}}{LevelGauge.lean}{178--186}
\crow{isMulRightInvariantgroupMeasure}{\href{https://github.com/lsernau/grh-lean/blob/72fead102f663d5c7a407964a13dc0cf9d977207/grh-compact/src/New/RealPlace/LevelGauge.lean\#L188-L189}{{\leanfont Level.\allowbreak{}isMulRightInvariant\_groupMeasure}}}{LevelGauge.lean}{188--189}
\crow{stackMeasurezeroeq}{\href{https://github.com/lsernau/grh-lean/blob/72fead102f663d5c7a407964a13dc0cf9d977207/grh-compact/src/New/RealPlace/LevelGauge.lean\#L193-L196}{{\leanfont Level.\allowbreak{}stackMeasure\_zero\_eq}}}{LevelGauge.lean}{193--196}
\crow{stackMeasurenezero}{\href{https://github.com/lsernau/grh-lean/blob/72fead102f663d5c7a407964a13dc0cf9d977207/grh-compact/src/New/RealPlace/LevelGauge.lean\#L200-L204}{{\leanfont Level.\allowbreak{}stackMeasure\_ne\_zero}}}{LevelGauge.lean}{200--204}
\crow{lintegralmullefteqoffactor}{\href{https://github.com/lsernau/grh-lean/blob/72fead102f663d5c7a407964a13dc0cf9d977207/grh-compact/src/New/RealPlace/LevelGauge.lean\#L211-L224}{{\leanfont Level.\allowbreak{}lintegral\_mul\_left\_eq\_of\_factor}}}{LevelGauge.lean}{211--224}
\crow{lintegralstackMeasuremulleft}{\href{https://github.com/lsernau/grh-lean/blob/72fead102f663d5c7a407964a13dc0cf9d977207/grh-compact/src/New/RealPlace/LevelGauge.lean\#L228-L232}{{\leanfont Level.\allowbreak{}lintegral\_stackMeasure\_mul\_left}}}{LevelGauge.lean}{228--232}
\crow{FloorMassLe}{\href{https://github.com/lsernau/grh-lean/blob/72fead102f663d5c7a407964a13dc0cf9d977207/grh-compact/src/New/RealPlace/LevelGauge.lean\#L238-L239}{{\leanfont Level.\allowbreak{}FloorMassLe}}}{LevelGauge.lean}{238--239}
\crow{upperPairing}{\href{https://github.com/lsernau/grh-lean/blob/72fead102f663d5c7a407964a13dc0cf9d977207/grh-compact/src/New/RealPlace/LevelGauge.lean\#L243-L243}{{\leanfont Level.\allowbreak{}upperPairing}}}{LevelGauge.lean}{243--243}
\crow{IsPairingKernel}{\href{https://github.com/lsernau/grh-lean/blob/72fead102f663d5c7a407964a13dc0cf9d977207/grh-compact/src/New/RealPlace/LevelGauge.lean\#L248-L255}{{\leanfont Level.\allowbreak{}IsPairingKernel}}}{LevelGauge.lean}{248--255}
\crow{invWeight}{\href{https://github.com/lsernau/grh-lean/blob/72fead102f663d5c7a407964a13dc0cf9d977207/grh-compact/src/New/RealPlace/LevelGauge.lean\#L259-L259}{{\leanfont Level.\allowbreak{}invWeight}}}{LevelGauge.lean}{259--259}
\crow{pivotModulus}{\href{https://github.com/lsernau/grh-lean/blob/72fead102f663d5c7a407964a13dc0cf9d977207/grh-compact/src/New/RealPlace/LevelGauge.lean\#L263-L263}{{\leanfont Level.\allowbreak{}pivotModulus}}}{LevelGauge.lean}{263--263}
\crow{weightedSquare}{\href{https://github.com/lsernau/grh-lean/blob/72fead102f663d5c7a407964a13dc0cf9d977207/grh-compact/src/New/RealPlace/LevelGauge.lean\#L267-L267}{{\leanfont Level.\allowbreak{}weightedSquare}}}{LevelGauge.lean}{267--267}
\crow{weightedSquareSum}{\href{https://github.com/lsernau/grh-lean/blob/72fead102f663d5c7a407964a13dc0cf9d977207/grh-compact/src/New/RealPlace/LevelGauge.lean\#L271-L272}{{\leanfont Level.\allowbreak{}weightedSquareSum}}}{LevelGauge.lean}{271--272}
\crow{selfNorm}{\href{https://github.com/lsernau/grh-lean/blob/72fead102f663d5c7a407964a13dc0cf9d977207/grh-compact/src/New/RealPlace/LevelGauge.lean\#L276-L276}{{\leanfont Level.\allowbreak{}selfNorm}}}{LevelGauge.lean}{276--276}
\crow{Gauge}{\href{https://github.com/lsernau/grh-lean/blob/72fead102f663d5c7a407964a13dc0cf9d977207/grh-compact/src/New/RealPlace/LevelGauge.lean\#L283-L301}{{\leanfont Gauge}}}{LevelGauge.lean}{283--301}
\crow{oneleLambda}{\href{https://github.com/lsernau/grh-lean/blob/72fead102f663d5c7a407964a13dc0cf9d977207/grh-compact/src/New/RealPlace/LevelGauge.lean\#L303-L304}{{\leanfont Gauge.\allowbreak{}one\_le\_Lambda}}}{LevelGauge.lean}{303--304}
\crow{positionsOfSize}{\href{https://github.com/lsernau/grh-lean/blob/72fead102f663d5c7a407964a13dc0cf9d977207/grh-compact/src/New/RealPlace/LevelGauge.lean\#L307-L308}{{\leanfont Gauge.\allowbreak{}positionsOfSize}}}{LevelGauge.lean}{307--308}
\crow{Size}{\href{https://github.com/lsernau/grh-lean/blob/72fead102f663d5c7a407964a13dc0cf9d977207/grh-compact/src/New/RealPlace/Size.lean\#L15-L22}{{\leanfont Size}}}{Size.lean}{15--22}
\crow{weylRatio}{\href{https://github.com/lsernau/grh-lean/blob/72fead102f663d5c7a407964a13dc0cf9d977207/grh-compact/src/New/RealPlace/Size.lean\#L29-L29}{{\leanfont weylRatio}}}{Size.lean}{29--29}
\crow{coeweylRatio}{\href{https://github.com/lsernau/grh-lean/blob/72fead102f663d5c7a407964a13dc0cf9d977207/grh-compact/src/New/RealPlace/Size.lean\#L31-L34}{{\leanfont coe\_weylRatio}}}{Size.lean}{31--34}
\crow{weylRatioinv}{\href{https://github.com/lsernau/grh-lean/blob/72fead102f663d5c7a407964a13dc0cf9d977207/grh-compact/src/New/RealPlace/Size.lean\#L36-L36}{{\leanfont weylRatio\_inv}}}{Size.lean}{36--36}
\crow{cardsucclefloorOf}{\href{https://github.com/lsernau/grh-lean/blob/72fead102f663d5c7a407964a13dc0cf9d977207/grh-compact/src/New/RealPlace/Size.lean\#L41-L42}{{\leanfont card\_succ\_le\_floorOf}}}{Size.lean}{41--42}
\crow{vecsOfSize}{\href{https://github.com/lsernau/grh-lean/blob/72fead102f663d5c7a407964a13dc0cf9d977207/grh-compact/src/New/RealPlace/Size.lean\#L47-L47}{{\leanfont Size.\allowbreak{}vecsOfSize}}}{Size.lean}{47--47}
\crow{HasMatSize}{\href{https://github.com/lsernau/grh-lean/blob/72fead102f663d5c7a407964a13dc0cf9d977207/grh-compact/src/New/RealPlace/Size.lean\#L50-L50}{{\leanfont Size.\allowbreak{}HasMatSize}}}{Size.lean}{50--50}
\crow{IsHead}{\href{https://github.com/lsernau/grh-lean/blob/72fead102f663d5c7a407964a13dc0cf9d977207/grh-compact/src/New/RealPlace/Size.lean\#L54-L54}{{\leanfont Size.\allowbreak{}IsHead}}}{Size.lean}{54--54}
\crow{starhead}{\href{https://github.com/lsernau/grh-lean/blob/72fead102f663d5c7a407964a13dc0cf9d977207/grh-compact/src/New/RealPlace/Size.lean\#L59-L59}{{\leanfont Size.\allowbreak{}IsHead.\allowbreak{}star\_head}}}{Size.lean}{59--59}
\crow{memofSizemono}{\href{https://github.com/lsernau/grh-lean/blob/72fead102f663d5c7a407964a13dc0cf9d977207/grh-compact/src/New/RealPlace/Size.lean\#L61-L62}{{\leanfont Size.\allowbreak{}mem\_ofSize\_mono}}}{Size.lean}{61--62}
\crow{submemofSize}{\href{https://github.com/lsernau/grh-lean/blob/72fead102f663d5c7a407964a13dc0cf9d977207/grh-compact/src/New/RealPlace/Size.lean\#L64-L65}{{\leanfont Size.\allowbreak{}sub\_mem\_ofSize}}}{Size.lean}{64--65}
\crow{monoa}{\href{https://github.com/lsernau/grh-lean/blob/72fead102f663d5c7a407964a13dc0cf9d977207/grh-compact/src/New/RealPlace/Size.lean\#L67-L68}{{\leanfont Size.\allowbreak{}HasMatSize.\allowbreak{}mono}}}{Size.lean}{67--68}
\crow{one}{\href{https://github.com/lsernau/grh-lean/blob/72fead102f663d5c7a407964a13dc0cf9d977207/grh-compact/src/New/RealPlace/Size.lean\#L70-L70}{{\leanfont Size.\allowbreak{}HasMatSize.\allowbreak{}one}}}{Size.lean}{70--70}
\crow{sub}{\href{https://github.com/lsernau/grh-lean/blob/72fead102f663d5c7a407964a13dc0cf9d977207/grh-compact/src/New/RealPlace/Size.lean\#L72-L73}{{\leanfont Size.\allowbreak{}HasMatSize.\allowbreak{}sub}}}{Size.lean}{72--73}
\crow{dotProductmemofSize}{\href{https://github.com/lsernau/grh-lean/blob/72fead102f663d5c7a407964a13dc0cf9d977207/grh-compact/src/New/RealPlace/Size.lean\#L75-L77}{{\leanfont Size.\allowbreak{}dotProduct\_mem\_ofSize}}}{Size.lean}{75--77}
\crow{mulVec}{\href{https://github.com/lsernau/grh-lean/blob/72fead102f663d5c7a407964a13dc0cf9d977207/grh-compact/src/New/RealPlace/Size.lean\#L79-L80}{{\leanfont Size.\allowbreak{}HasMatSize.\allowbreak{}mulVec}}}{Size.lean}{79--80}
\crow{mul}{\href{https://github.com/lsernau/grh-lean/blob/72fead102f663d5c7a407964a13dc0cf9d977207/grh-compact/src/New/RealPlace/Size.lean\#L82-L83}{{\leanfont Size.\allowbreak{}HasMatSize.\allowbreak{}mul}}}{Size.lean}{82--83}
\crow{detmemofSize}{\href{https://github.com/lsernau/grh-lean/blob/72fead102f663d5c7a407964a13dc0cf9d977207/grh-compact/src/New/RealPlace/Size.lean\#L85-L88}{{\leanfont Size.\allowbreak{}det\_mem\_ofSize}}}{Size.lean}{85--88}
\crow{adjointapplymemofSize}{\href{https://github.com/lsernau/grh-lean/blob/72fead102f663d5c7a407964a13dc0cf9d977207/grh-compact/src/New/RealPlace/Size.lean\#L92-L97}{{\leanfont Size.\allowbreak{}adjoint\_apply\_mem\_ofSize}}}{Size.lean}{92--97}
\crow{inv}{\href{https://github.com/lsernau/grh-lean/blob/72fead102f663d5c7a407964a13dc0cf9d977207/grh-compact/src/New/RealPlace/Size.lean\#L99-L101}{{\leanfont Size.\allowbreak{}HasMatSize.\allowbreak{}inv}}}{Size.lean}{99--101}
\crow{rowmemofSizeofdualframe}{\href{https://github.com/lsernau/grh-lean/blob/72fead102f663d5c7a407964a13dc0cf9d977207/grh-compact/src/New/RealPlace/Size.lean\#L103-L110}{{\leanfont Size.\allowbreak{}row\_mem\_ofSize\_of\_dual\_frame}}}{Size.lean}{103--110}
\crow{hasMatSizeonesubinv}{\href{https://github.com/lsernau/grh-lean/blob/72fead102f663d5c7a407964a13dc0cf9d977207/grh-compact/src/New/RealPlace/Size.lean\#L115-L125}{{\leanfont Size.\allowbreak{}hasMatSize\_one\_sub\_inv}}}{Size.lean}{115--125}
\crow{minormuldualmemofSize}{\href{https://github.com/lsernau/grh-lean/blob/72fead102f663d5c7a407964a13dc0cf9d977207/grh-compact/src/New/RealPlace/Size.lean\#L133-L146}{{\leanfont Size.\allowbreak{}minor\_mul\_dual\_mem\_ofSize}}}{Size.lean}{133--146}
\crow{minorsmemofSize}{\href{https://github.com/lsernau/grh-lean/blob/72fead102f663d5c7a407964a13dc0cf9d977207/grh-compact/src/New/RealPlace/Size.lean\#L150-L154}{{\leanfont Size.\allowbreak{}minors\_mem\_ofSize}}}{Size.lean}{150--154}
\crow{headentriesmemofSize}{\href{https://github.com/lsernau/grh-lean/blob/72fead102f663d5c7a407964a13dc0cf9d977207/grh-compact/src/New/RealPlace/Size.lean\#L158-L164}{{\leanfont Size.\allowbreak{}head\_entries\_mem\_ofSize}}}{Size.lean}{158--164}
\crow{hasMatSizeofisHead}{\href{https://github.com/lsernau/grh-lean/blob/72fead102f663d5c7a407964a13dc0cf9d977207/grh-compact/src/New/RealPlace/Size.lean\#L168-L174}{{\leanfont Size.\allowbreak{}hasMatSize\_of\_isHead}}}{Size.lean}{168--174}
\crow{floorMassLeofchart}{\href{https://github.com/lsernau/grh-lean/blob/72fead102f663d5c7a407964a13dc0cf9d977207/grh-compact/src/New/RealPlace/WeightPlancherel.lean\#L20-L31}{{\leanfont Gauge.\allowbreak{}floorMassLe\_of\_chart}}}{WeightPlancherel.lean}{20--31}
\crow{floorMassLefloorWeightdiv}{\href{https://github.com/lsernau/grh-lean/blob/72fead102f663d5c7a407964a13dc0cf9d977207/grh-compact/src/New/RealPlace/WeightPlancherel.lean\#L35-L45}{{\leanfont Gauge.\allowbreak{}floorMassLe\_floorWeight\_div}}}{WeightPlancherel.lean}{35--45}
\crow{floorMassLeone}{\href{https://github.com/lsernau/grh-lean/blob/72fead102f663d5c7a407964a13dc0cf9d977207/grh-compact/src/New/RealPlace/WeightPlancherel.lean\#L49-L52}{{\leanfont Gauge.\allowbreak{}floorMassLe\_one}}}{WeightPlancherel.lean}{49--52}
\crow{floorConst}{\href{https://github.com/lsernau/grh-lean/blob/72fead102f663d5c7a407964a13dc0cf9d977207/grh-compact/src/New/RealPlace/WeightPlancherel.lean\#L56-L56}{{\leanfont Gauge.\allowbreak{}floorConst}}}{WeightPlancherel.lean}{56--56}
\crow{lintegralfloorWeightle}{\href{https://github.com/lsernau/grh-lean/blob/72fead102f663d5c7a407964a13dc0cf9d977207/grh-compact/src/New/RealPlace/WeightPlancherel.lean\#L62-L90}{{\leanfont Gauge.\allowbreak{}lintegral\_floorWeight\_le}}}{WeightPlancherel.lean}{62--90}
\crow{globalConst}{\href{https://github.com/lsernau/grh-lean/blob/72fead102f663d5c7a407964a13dc0cf9d977207/grh-compact/src/New/RealPlace/WeightPlancherel.lean\#L93-L93}{{\leanfont Gauge.\allowbreak{}globalConst}}}{WeightPlancherel.lean}{93--93}
\crow{lintegralmulinvWeightle}{\href{https://github.com/lsernau/grh-lean/blob/72fead102f663d5c7a407964a13dc0cf9d977207/grh-compact/src/New/RealPlace/WeightPlancherel.lean\#L98-L135}{{\leanfont Gauge.\allowbreak{}lintegral\_mul\_invWeight\_le}}}{WeightPlancherel.lean}{98--135}
\crow{amplitudesupportandplancherel}{\href{https://github.com/lsernau/grh-lean/blob/72fead102f663d5c7a407964a13dc0cf9d977207/grh-compact/src/New/RealPlace/WeightPlancherel.lean\#L141-L188}{{\leanfont Gauge.\allowbreak{}amplitude\_support\_and\_plancherel}}}{WeightPlancherel.lean}{141--188}
\crow{parsevalheisenbergMeasureofbox}{\href{https://github.com/lsernau/grh-lean/blob/72fead102f663d5c7a407964a13dc0cf9d977207/grh-compact/src/New/RealPlace/WeightPlancherel.lean\#L194-L220}{{\leanfont Gauge.\allowbreak{}parseval\_heisenbergMeasure\_of\_box}}}{WeightPlancherel.lean}{194--220}}
% the notations of the defined functions: one slot for each explicit argument of the Lean declaration
\newcommand{\HasFibreBound}[5]{\operatorname{Fib}_{#2,#4}(#1;#3)\le #5}
\newcommand{\Levels}[3]{\operatorname{Lev}(#1,#2,#3)}
\newcommand{\floorMeasure}[2]{\theta_{#2}(#1)}
\newcommand{\stackMeasure}[1]{\mu^{#1}}
\newcommand{\groupMeasure}[1]{\mathsf{M}(#1)}
\newcommand{\upperMeasure}[2]{\mathsf{M}_{\ge #2}(#1)}
\newcommand{\IsNear}[3]{#3\in\operatorname{Near}_{#2}(#1)}
\newcommand{\orthonormalPart}[1]{\mathsf{O}(#1)}
\newcommand{\closePairs}[1]{\operatorname{Cl}(#1)}
\newcommand{\floorValues}[4]{\operatorname{Val}_{#2}(#1;#3,#4)}
\newcommand{\volumeExponent}[1]{\mathsf{s}(#1)}
\newcommand{\floorWeight}[4]{\mathsf{w}_{#2}(#1;#3,#4)}
\newcommand{\floorModulus}[3]{\mathsf{d}_{#2}(#1;#3)}
\newcommand{\FloorMassLe}[4]{\operatorname{Mass}_{#2}(#1;#4)\le #3}
\newcommand{\upperPairing}[5]{\mathsf{Av}_{#2}(#1;#3)(#4,#5)}
\newcommand{\IsPairingKernel}[4]{(#3,#4)\in\operatorname{PK}_{#2}(#1)}
\newcommand{\invWeight}[3]{\mathsf{w}^{-1}_{#2}(#1;#3)}
\newcommand{\pivotModulus}[5]{\mathsf{Pv}_{#2}(#1;#3)(#4,#5)}
\newcommand{\weightedSquare}[5]{\mathsf{Sq}_{#2}(#1;#3,#4)(#5)}
\newcommand{\weightedSquareSum}[5]{\mathsf{Sm}_{#4}(#1;#2,#3)(#5)}
\newcommand{\selfNorm}[4]{\operatorname{SN}_{#2}(#1;#3,#4)}
\newcommand{\Gauges}[1]{\operatorname{Gau}(#1)}
\newcommand{\positionsOfSize}[2]{\operatorname{Pos}(#1;#2)}
\newcommand{\Sizes}[1]{\operatorname{Siz}(#1)}
\newcommand{\weylRatio}[3]{w_{#1}(#3)^{-1}w_{#1}(#2)}
\newcommand{\vecsOfSize}[1]{\mathrm{Vec}_{#1}}
\newcommand{\HasMatSize}[2]{\operatorname{sz}(#2)\le #1}
\newcommand{\IsHead}[2]{#2\in\operatorname{Hd}_{#1}}
\newcommand{\floorConst}[2]{C_{#2}(#1)}
\newcommand{\globalConst}[1]{C_{\mathrm{gl}}(#1)}
\newcommand{\UnitMembers}{239}
\newcommand{\UnitInNote}{116}
\newcommand{\UnitProofs}{108}
\newcommand{\UnitSteps}{639}
\newcommand{\UnitAllSteps}{665}


\linespread{1.08}                       % 11 pt on 14.7 pt
\setlength{\parindent}{0pt}
\setlength{\parskip}{5.5pt plus 1.5pt minus 0.5pt}
\setlength{\abovedisplayskip}{7.5pt plus 2pt minus 3pt}
\setlength{\belowdisplayskip}{7.5pt plus 2pt minus 3pt}
\setlength{\abovedisplayshortskip}{3pt plus 2pt}
\setlength{\belowdisplayshortskip}{6pt plus 2pt minus 2pt}
\emergencystretch=1.5em
\clubpenalty=10000 \widowpenalty=10000 \displaywidowpenalty=10000
\relpenalty=7000 \binoppenalty=9000
\raggedbottom
\pagestyle{fancy}\fancyhf{}\renewcommand{\headrulewidth}{0pt}
\fancyfoot[C]{\sffamily\fontsize{7.6}{9}\selectfont\color{inkfaint}\thepage\,/\,\pageref*{LastPage}}
\setlist{topsep=2pt,partopsep=0pt,parsep=0pt,itemsep=4pt,leftmargin=2.1em,labelsep=0.55em}

\newcommand{\Q}{\mathbb{Q}}
\newcommand{\Z}{\mathbb{Z}}
\newcommand{\F}{\mathbb{F}}
\newcommand{\len}{\operatorname{len}}
\newcommand{\cok}{\operatorname{cok}}
\newcommand{\End}{\operatorname{End}}
\newcommand{\Hom}{\operatorname{Hom}}
\newcommand{\Gal}{\operatorname{Gal}}
\newcommand{\diag}{\operatorname{diag}}
\newcommand{\kw}[1]{\emph{#1}}
\newcommand{\m}[1]{\mbox{$#1$}}        % a short formula that must not be broken across lines
\newcommand{\head}[1]{\par\addvspace{17pt plus 4pt minus 2pt}\Needspace{6\baselineskip}%
  \noindent{\fontsize{13.2}{16}\selectfont\bfseries\color{ink}#1}\par\nopagebreak\vspace{-1pt}}
\newcommand{\sect}[2]{\refstepcounter{section}\label{#2}\par\addvspace{19pt plus 4pt minus 2pt}\Needspace{7\baselineskip}%
  \noindent\llap{{\fontsize{13.2}{16}\selectfont\bfseries\color{inkfaint}\thesection}\hspace{9pt}}%
  {\fontsize{13.2}{16}\selectfont\bfseries\color{ink}#1}\par\nopagebreak\vspace{-1pt}}
\newcommand{\subhead}[1]{\par\addvspace{9pt plus 2pt}\Needspace{4\baselineskip}\noindent\textbf{#1}\enspace\ignorespaces}
% where a statement lives in the Lean archive
\newcommand{\leantag}{{\sffamily\bfseries\fontsize{6.3}{8}\selectfont\color{leancol}\addfontfeatures{LetterSpace=7}LEAN}}
\newcommand{\sep}{\unskip\nobreak\hspace{5pt}{\color{inkfaint}\textperiodcentered}\hspace{5pt plus 2pt}\ignorespaces}
\newcommand{\leanline}[1]{\par\nopagebreak\vspace{-2.5pt}{\fontsize{8.9}{12.4}\selectfont\color{inksoft}\raggedright
  \noindent\llap{\leantag\hspace{9pt}}#1\par}\vspace{1pt}}
% a displayed statement in a tinted box: #1 colour, #2 text
\newcommand{\statement}[2]{\par\addvspace{7pt}\begingroup\setlength{\fboxsep}{8.5pt}\setlength{\fboxrule}{0.6pt}%
  \noindent\fcolorbox{#1!30!white}{#1!3!white}{\parbox{\dimexpr\linewidth-2\fboxsep-2\fboxrule\relax}{#2}}%
  \endgroup\par\addvspace{7pt}}
\newcommand{\nt}[2]{\par\noindent\hangindent=1.1em\hangafter=1 #1\enspace{\color{inksoft}#2}\par\vspace{1pt}}
\newcommand{\legdot}[1]{\tikz[baseline=-0.62ex]\fill[#1] (0,0) circle (2.5pt);\hspace{3pt}}
\newcommand{\crow}[4]{\csname CS@#1\endcsname & #2 & \textcolor{inksoft}{#3} & \textcolor{inksoft}{#4}\\}
\newcommand{\leanquoteline}[2]{\par\noindent\hangindent=\dimexpr#1\fontdimen2\font+2.2em\relax\hangafter=1
  \hspace*{#1\fontdimen2\font}#2}
\newcommand{\quoted}[1]{\par\vspace{3pt}{\ttfamily\fontsize{7.5}{10.4}\selectfont\color{leancol}\raggedright
  \setlength{\parskip}{0pt}\leanquote{#1}\par}\vspace{4pt}}
\newcommand{\bx}[3]{{\sffamily\bfseries\fontsize{7.3}{9.2}\selectfont\color{#1!82!black}#2}\\[1.5pt]#3}
\newcommand{\cs}[2]{\expandafter\newcommand\csname CS@#1\endcsname{#2}}

% theorems
\newtheoremstyle{plainx}{7pt plus 2pt}{7pt plus 2pt}{\itshape}{}{\bfseries}{.}{0.5em}{}
\newtheoremstyle{uprightx}{7pt plus 2pt}{7pt plus 2pt}{}{}{\bfseries}{.}{0.5em}{}
\theoremstyle{plainx}
\newtheorem{theorem}{Theorem}[section]
\newtheorem{lemma}[theorem]{Lemma}
\newtheorem{proposition}[theorem]{Proposition}
\newtheorem{corollary}[theorem]{Corollary}
\theoremstyle{uprightx}
\newtheorem{definition}[theorem]{Definition}
\newtheorem{setting}[theorem]{Setting}          % hypotheses shared by several statements; each of them cites it by number
\newtheorem{remark}[theorem]{Remark}
\newtheorem{example}[theorem]{Example}
\renewcommand{\qedsymbol}{{\setlength{\fboxsep}{0pt}\setlength{\fboxrule}{0.5pt}\fbox{\rule{0pt}{0.46em}\hspace{0.46em}}}}
% what holds throughout the note: the ambient objects only (implicit and instance arguments), never a hypothesis
\newenvironment{conventions}{\par\addvspace{5pt}\noindent\textbf{Conventions.}\enspace\ignorespaces}{\par\addvspace{3pt}}
\makeatletter
\renewenvironment{proof}[1][Proof]{\par\pushQED{\qed}\normalfont\topsep0pt\partopsep0pt
  \trivlist\item[\hskip\labelsep\itshape #1.]\ignorespaces}{\popQED\endtrivlist\@endpefalse}
\makeatother

% notations of definitions of unit 132 that this note uses and another note of the unit states
% (made by shared/mkspec.py from shared/notation-132.json; the same TeX as in the note that states them)
\newcommand{\sieveVector}[2]{\Xi(#1;#2)}
\newcommand{\RealPlaceSizes}[4]{\operatorname{RPS}(#1,#2,#3,#4)}
\newcommand{\haarRatio}[1]{\lambda(#1)}
\newcommand{\gauge}[2]{\Gamma(#1;#2)}
\newcommand{\IsCutBessel}[7]{\operatorname{CB}(#1,#2,#3,#4,#5,#6,#7)}
\newcommand{\formConst}[1]{c_{\mathrm{f}}(#1)}
\newcommand{\gramSection}[1]{\gamma(#1)}
\newcommand{\invCorner}[1]{\operatorname{ic}(#1)}
\newcommand{\cornerDefect}[1]{\operatorname{cd}(#1)}
\newcommand{\IsBoundedInCell}[2]{#1\in\operatorname{Cell}_{\le #2}}
\newcommand{\diagLast}[1]{\operatorname{d}_{\varrho}(#1)}
\newcommand{\sectionCoeff}[2]{\mathsf{c}(#1,#2)}
\newcommand{\sectionFrame}[3]{\mathsf{f}(#1,#2,#3)}
\newcommand{\sectionMatrix}[3]{\mathsf{m}(#1,#2,#3)}
\newcommand{\headDensity}[4]{\mathsf{h}_{#1}(#4;#2,#3)}
\newcommand{\translatePairing}[5]{\langle #1(#4),#2(#3#5)\rangle}
\newcommand{\IsBessel}[2]{#1\in\operatorname{Bes}(#2)}
\newcommand{\IsCut}[1]{#1\in\operatorname{Cut}}
\newcommand{\pairingNorm}[2]{\operatorname{PN}(#1,#2)}
\newcommand{\formIntegrandOne}[6]{\mathsf{I}^{1}_{#1}(#2;#3,#4,#5)(#6)}
\newcommand{\formIntegrand}[7]{\mathsf{I}_{#1}(#2,#3;#4,#5,#6)(#7)}
\newcommand{\floorForm}[5]{\mathcal{I}_{#1}(#2,#3;#4,#5)}
\newcommand{\formModulus}[6]{\mathcal{I}^{\mathrm{mod}}_{#1}(#2,#3;#4,#5,#6)}
\newcommand{\formMajorant}[6]{\mathcal{I}^{\mathrm{maj}}_{#1}(#2,#3;#4,#5,#6)}
\newcommand{\selfNormProduct}[7]{\operatorname{SNP}_{#1}(#7;#2,#3;#4,#5,#6)}
\newcommand{\frameVolume}[1]{\Delta(#1)}
\newcommand{\frameXPart}[1]{\mathsf{x}_0(#1)}

% hand macros of the three notes of unit 132 (objects of the background library; letters of shared/PLAN-132.md)
\newcommand{\cL}{\mathcal{L}}
\newcommand{\CC}{\mathbb{C}}
\newcommand{\RR}{\mathbb{R}}
\newcommand{\NN}{\mathbb{N}}
\newcommand{\dd}{\mathrm{d}}
\newcommand{\GT}[1]{G(#1)}
\newcommand{\JT}[1]{J(#1)}
\newcommand{\Hl}[2]{\mathcal{H}_{#1}(#2)}
\newcommand{\Ql}[2]{\mathcal{Q}_{#1}(#2)}
\newcommand{\Nb}[2]{\mathcal{N}_{<#1}(#2)}
\newcommand{\Nf}[2]{\mathcal{N}_{#1}(#2)}
\newcommand{\UT}[1]{\mathcal{U}(#1)}
\newcommand{\modA}{\operatorname{mod}_A}
\newcommand{\ZT}[2]{Z_{#1}(#2)}
\newcommand{\sd}{\operatorname{sd}}
\newcommand{\wl}[2]{w_{#1}(#2)}
\newcommand{\wo}[2]{w^{\circ}_{#1}(#2)}
\newcommand{\ip}[2]{\langle #1,#2\rangle}
\newcommand{\gL}[1]{g_{#1}}
\newcommand{\BoxL}[1]{\operatorname{Box}(#1)}
\newcommand{\Posn}[1]{\operatorname{Posn}(#1)}
\newcommand{\Ball}[1]{\mathfrak{B}(#1)}
\newcommand{\etab}[1]{\eta_{<#1}}
\newcommand{\vol}{\operatorname{vol}}
\newcommand{\cLC}[1]{\mathcal{L}_{\CC}(#1)}
\newcommand{\Gone}[1]{G_1(#1)}
\newcommand{\Jone}[1]{J_1(#1)}
\newcommand{\volU}[1]{\vol\mathrm{U}(#1)}
\newcommand{\hw}[3]{\operatorname{hw}_{#1,#2}(#3)}
\newcommand{\damp}[2]{\ell_{#1}(#2)}
\newcommand{\bandc}[1]{\Theta(#1)}
\newcommand{\cellc}[1]{c_G(#1)}
\newcommand{\corr}[3]{J^{#1}_{#2}(#3)}
\newcommand{\boost}[2]{\chi_{#1}(#2)}
\newcommand{\quadric}[2]{\delta^{(#2)}(#1)}
\newcommand{\Inp}[1]{\textup{(#1)}}
% where each declaration of the concordance is treated
\cs{Size}{Definition~\ref{d:size}}
\cs{weylRatio}{Definition~\ref{d:wr}}
\cs{coeweylRatio}{Definition~\ref{d:wr}}
\cs{weylRatioinv}{Definition~\ref{d:wr}}
\cs{vecsOfSize}{Definition~\ref{d:sz}}
\cs{HasMatSize}{Definition~\ref{d:sz}}
\cs{IsHead}{Definition~\ref{d:sz}}
\cs{starhead}{Lemma~\ref{l:rules}}
\cs{memofSizemono}{Lemma~\ref{l:rules}}
\cs{submemofSize}{Lemma~\ref{l:rules}}
\cs{monoa}{Lemma~\ref{l:rules}}
\cs{one}{Lemma~\ref{l:rules}}
\cs{sub}{Lemma~\ref{l:rules}}
\cs{dotProductmemofSize}{Lemma~\ref{l:rules}}
\cs{mulVec}{Lemma~\ref{l:rules}}
\cs{mul}{Lemma~\ref{l:rules}}
\cs{detmemofSize}{Lemma~\ref{l:det}}
\cs{adjointapplymemofSize}{Lemma~\ref{l:adj}}
\cs{inv}{Lemma~\ref{l:adj}}
\cs{rowmemofSizeofdualframe}{Lemma~\ref{l:row}}
\cs{hasMatSizeonesubinv}{Lemma~\ref{l:backsub}}
\cs{minormuldualmemofSize}{Lemma~\ref{l:minor}}
\cs{minorsmemofSize}{Lemma~\ref{l:minor}}
\cs{headentriesmemofSize}{Lemma~\ref{l:minor}}
\cs{hasMatSizeofisHead}{Lemma~\ref{l:minor}}
\cs{cardsucclefloorOf}{Lemma~\ref{l:minor}}
\cs{Level}{Definition~\ref{d:level}}
\cs{floorMeasure}{Definition~\ref{d:fm}}
\cs{floorMeasuretop}{Definition~\ref{d:fm}}
\cs{stackMeasure}{Definition~\ref{d:tower}}
\cs{groupMeasure}{Definition~\ref{d:tower}}
\cs{upperMeasure}{Definition~\ref{d:tower}}
\cs{upperMeasureeqsuccmconv}{Definition~\ref{d:tower}}
\cs{aestackMeasuremem}{Definition~\ref{d:tower}}
\cs{IsNear}{Definition~\ref{d:near}}
\cs{isNearmuliff}{Definition~\ref{d:near}}
\cs{offloorGroupmul}{Definition~\ref{d:near}}
\cs{orthonormalPart}{Definition~\ref{d:val}}
\cs{closePairs}{Definition~\ref{d:val}}
\cs{floorValues}{Definition~\ref{d:val}}
\cs{volumeExponent}{Definition~\ref{d:weight}}
\cs{floorWeight}{Definition~\ref{d:weight}}
\cs{invWeight}{Definition~\ref{d:weight}}
\cs{floorModulus}{Definition~\ref{d:weight}}
\cs{floorWeightsucc}{Definition~\ref{d:weight}}
\cs{floorWeightnezeroandnetop}{Definition~\ref{d:weight}}
\cs{floorWeightmulleft}{Definition~\ref{d:weight}}
\cs{measurablefloorWeight}{Definition~\ref{d:weight}}
\cs{measurableinvWeight}{Definition~\ref{d:weight}}
\cs{FloorMassLe}{Definition~\ref{d:mass}}
\cs{upperPairing}{Definition~\ref{d:av}}
\cs{pivotModulus}{Definition~\ref{d:av}}
\cs{IsPairingKernel}{Definition~\ref{d:pk}}
\cs{weightedSquare}{Definition~\ref{d:sn}}
\cs{weightedSquareSum}{Definition~\ref{d:sn}}
\cs{selfNorm}{Definition~\ref{d:sn}}
\cs{Gauge}{Definition~\ref{d:gauge}}
\cs{oneleLambda}{Definition~\ref{d:gauge}}
\cs{positionsOfSize}{Definition~\ref{d:consts}}
\cs{floorConst}{Definition~\ref{d:consts}}
\cs{globalConst}{Definition~\ref{d:consts}}
\cs{measurableSetpositionsOfSize}{Definition~\ref{d:consts}}
\cs{measurableSetvecsOfSize}{Definition~\ref{d:consts}}
\cs{HasFibreBound}{Definition~\ref{d:fib}}
\cs{measurePreservingorbitPoint}{Lemma~\ref{l:orbit}}
\cs{measurePreservingmulVecorbitPoint}{Lemma~\ref{l:orbit}}
\cs{aefloorMeasuremem}{Lemma~\ref{l:carried}}
\cs{lintegralfloorMeasurecongr}{Lemma~\ref{l:carried}}
\cs{descConvfloorMeasuresucc}{Lemma~\ref{l:desc}}
\cs{lintegralweylRatiomul}{Lemma~\ref{l:weylratio}}
\cs{lintegralfloorMeasuremulright}{Lemma~\ref{l:right}}
\cs{stackMeasurenezero}{Lemma~\ref{l:nz}}
\cs{lintegralmullefteqoffactor}{Lemma~\ref{l:factor}}
\cs{memvecsOfSizeofheisenbergElem}{Lemma~\ref{l:heis}}
\cs{measurableSetisNear}{Lemma~\ref{l:meas}}
\cs{isOpenfloorValues}{Lemma~\ref{l:meas}}
\cs{volumefloorValuesnezeronetop}{Lemma~\ref{l:meas}}
\cs{floorValuesmul}{Lemma~\ref{l:valinv}}
\cs{invWeightmul}{Lemma~\ref{l:valinv}}
\cs{upperMeasureisNearle}{Lemma~\ref{l:upper}}
\cs{memvecsOfSizeoffloorBasismulVec}{Lemma~\ref{l:basis}}
\cs{setLIntegralleofhasFibreBound}{Lemma~\ref{l:fibint}}
\cs{setLIntegralorthonormalPartle}{Lemma~\ref{l:fibint}}
\cs{hasFibreBoundfloorQuadricofle}{Lemma~\ref{l:fibq}}
\cs{hasFibreBoundfloorQuadric}{Lemma~\ref{l:fibq}}
\cs{hasFibreBoundfloorVolume}{Lemma~\ref{l:fibq}}
\cs{lintegralheisenbergMeasurele}{Lemma~\ref{l:heisint}}
\cs{conjmulinvclosetoone}{Lemma~\ref{l:three}}
\cs{orbitPointmemvecsOfSize}{Lemma~\ref{l:three}}
\cs{existshasMatSizeofisNear}{Lemma~\ref{l:three}}
\cs{hasMatSizesuboneoffloorValues}{Lemma~\ref{l:subone}}
\cs{existsnezeroofpivotModulus}{Lemma~\ref{l:nokernel}}
\cs{lintegralupperPairingsqle}{Lemma~\ref{l:nokernel}}
\cs{integrableupperMeasure}{Lemma~\ref{l:succ}}
\cs{upperPairingeqintegralsucc}{Lemma~\ref{l:succ}}
\cs{enormupperPairingle}{Lemma~\ref{l:succ}}
\cs{stronglyMeasurableupperPairing}{Lemma~\ref{l:succ}}
\cs{mono}{Lemma~\ref{l:succ}}
\cs{upperPairinginvariant}{Lemma~\ref{l:kinv}}
\cs{weightedSquareinvariantsupport}{Lemma~\ref{l:sqinv}}
\cs{weightedSquareSuminvariantsupport}{Lemma~\ref{l:sqinv}}
\cs{measurableweightedSquareSum}{Lemma~\ref{l:sqinv}}
\cs{pivotModulussqle}{Lemma~\ref{l:cs}}
\cs{selfNormeqaddsucc}{Lemma~\ref{l:sn}}
\cs{selfNormleselfNormzero}{Lemma~\ref{l:sn}}
\cs{floorMassLeofchart}{Lemma~\ref{l:chart}}
\cs{floorMassLefloorWeightdiv}{Lemma~\ref{l:mass}}
\cs{floorMassLeone}{Lemma~\ref{l:mass}}
\cs{stackMeasuremconvdirac}{Theorem A (i)}
\cs{isMulRightInvariantgroupMeasure}{Theorem A (i)}
\cs{stackMeasurezeroeq}{Theorem A (ii)}
\cs{lintegralstackMeasuremulleft}{Theorem A (iii)}
\cs{lintegralfloorWeightle}{Theorem B (i)}
\cs{lintegralmulinvWeightle}{Theorem B (ii)}
\cs{amplitudesupportandplancherel}{Theorem C (i)}
\cs{parsevalheisenbergMeasureofbox}{Theorem C (ii)}

\begin{document}
\color{ink}

{\raggedright
{\sffamily\bfseries\fontsize{7.2}{9}\selectfont\color{leancol}\addfontfeatures{LetterSpace=8}READBACK NOTE 28\par}
\vspace{5pt}
{\fontsize{21.5}{25.5}\selectfont\bfseries The tower method at a place: floor measures, gauges and the bounds on each floor\par}
\vspace{5pt}
{\fontsize{12.4}{16}\selectfont\itshape\color{inksoft} One unit of the Lean artifact \ArtifactName, read back into
mathematics\par}}
\vspace{8pt}
{\color{rulegrey}\hrule height 0.6pt}
\vspace{7pt}
\thispagestyle{empty}

\textbf{Summary.} Let $\GT{T}$ be the unitary group of a hermitian space with $\varrho$ hyperbolic planes over a local
ring $A$. The stabilisers $\Hl{l}{T}$ of the first $l$ planes form a tower, and the quotient of two consecutive
ones is a quadric. A \emph{level} puts a measure on each floor, the image of an explicit measure on a big cell,
and their convolution $\groupMeasure{\cL}$ is a right invariant measure on $\GT{T}$: the method works when it is a
Haar measure. A kernel on $\GT{T}\times\GT{T}$ is averaged floor by floor, and on each floor three tools bound the
averages: a Cauchy--Schwarz inequality against the mass of the floor, integrals of weights made of volumes, and
Plancherel's inequality on the quadric and on a Heisenberg group. All sizes are kept by a \emph{gauge}, a filtration
of~$A$ with one constant~$\Lambda$. This note sets up these objects and proves the tools; the companion notes
apply them.

\statement{leancol}{\fontsize{9.6}{13.2}\selectfont\textbf{What this note is.} The artifact \ArtifactName{} is a
body of Lean code whose final theorems are the Riemann hypothesis and the generalized Riemann hypothesis for Dirichlet
$L$-functions. Its authors report that it compiles and that Lean's kernel accepts these theorems from the three
standard axioms; nothing was built for this note. That is a statement about a formal text: no mathematician has read
the argument through, and this note does not say that the Riemann hypothesis is proved. The note treats one part of
one unit of the artifact (the folder of the analysis at the real places): \UnitInNote{} of the unit's \UnitMembers{}
declarations, those of five files; two companion notes treat the others (Note~L, \emph{The local theorem for floor
forms and the band limits at the complex place}; Note~C, \emph{The complex place: the open-cell chart, the Haar measure
of the tower, the Bessel cut and the volume weight}). Every result of the five files is stated here, compared
argument by argument with its Lean text, and proved from the inputs of Section~\ref{s:inputs}, which belong to the
background library; the last section says which proofs follow the formal text step by step and which were only
traced through its terms. The Lean text read is \PinText; Appendix~A is the concordance of \ConcDecls{} declarations.}

\sect{The statements}{s:intro}
The objects are defined in Section~\ref{s:setting}. $A$ is a place, $T$ an $\mathfrak{s}$-form over $A$, $\GT{T}$
the unitary group of its split Gram matrix, $\cL$ a level, $\Gamma$ a gauge on~$\cL$.

A word on the names. A \emph{real place} of the write-up is an infinite place of the maximal totally real subfield
of the CM field $E(D)$ of the overview; the completion of $E(D)$ above it is $\CC$, and the unitary group of the
rows there is $\GT{T}$ for $A=\CC$ and the form $T$ that is the identity on the $x_i$ and $-1$ at $z$, a unitary group
of signature $(\varrho+k,\varrho+1)$, written $G_1(\mathfrak{s})$ in the companion notes. In the formalism of places
of this note that is the place $A=\CC$ with complex conjugation, which the three notes call \emph{the complex
place}; so ``the real places'' (the name of the unit and of the overview) and ``the complex place'' (Notes L and~C)
are the same local situation. This note works at a general place.

\statement{arith}{\phantomsection\label{t:A}\textbf{Theorem A (the tower measure).}\enspace Let $A$,
$\mathfrak{s}$, $T$, $\Omega$, $\mathfrak{m}$ be as in the conventions and $\cL$ a level in
$\Levels{\Omega}{T}{\mathfrak{m}}$.

\textbf{(i)} Let $l$, $i$ be natural numbers with $l+i=\varrho+2$ and $s\in\Hl{l}{T}$. Then
$\stackMeasure{\cL}_{l,i}*\delta_s=\stackMeasure{\cL}_{l,i}$. In particular $\groupMeasure{\cL}$ is invariant under
right translations.

\textbf{(ii)} For every natural number $j$:
$\stackMeasure{\cL}_{0,j}=\etab{j}*(\floorMeasure{\cL}{j-1}*\dots*\floorMeasure{\cL}{0})$.

\textbf{(iii)} Suppose that $\groupMeasure{\cL}$ is a Haar measure. Let $j$ be a natural number with
$j\le\varrho+2$, let $p\colon\NN\to A^{\times}$, and let $\Phi\colon\GT{T}\to[0,\infty]$ be measurable and invariant
under left translations by $\Hl{j}{T}$. Then
\[\int\Phi\bigl(\wl{j}{p}^{-1}g\bigr)\,\dd\stackMeasure{\cL}_{0,j}(g)=\int\Phi(g)\,\dd\stackMeasure{\cL}_{0,j}(g).\]}
\leanline{\lean{stackMeasuremconvdirac}\sep\lean{isMulRightInvariantgroupMeasure}\sep\lean{stackMeasurezeroeq}\sep\lean{lintegralstackMeasuremulleft}}

\begin{setting}[a level with a gauge]\label{set:gauge}
Let $A$, $\mathfrak{s}$, $T$, $\Omega$, $\mathfrak{m}$ be as in the conventions, $\cL$ a level in
$\Levels{\Omega}{T}{\mathfrak{m}}$ and $\Gamma$ a gauge in $\Gauges{\cL}$. Sizes ($\operatorname{sz}$,
$\mathrm{Vec}$, heads) are taken for $\Gamma$.
\end{setting}

\statement{arith}{\phantomsection\label{t:B}\textbf{Theorem B (two integrals of weights).}\enspace In
Setting~\ref{set:gauge}, let $j$ be a natural number with $j\le\varrho$.

\textbf{(i)} Let $g,y\in\GT{T}$ with $g$ near at $j$: \m{\IsNear{\cL}{j}{g}}. Then
\[\int\mathbf{1}_{\positionsOfSize{\cL}{\Gamma}}(g^{-1}\nu y)\;\floorWeight{\cL}{j}{g}{g^{-1}\nu y}\;\dd\etab{j}(\nu)\ \le\ \prod_{l<j}\floorConst{\Gamma}{l}.\]

\textbf{(ii)} Let $p\colon\NN\to A^{\times}$ be a head of length $j$, \m{\IsHead{j}{p}}, let $x\in\Posn{\cL}$, let
$b$ be a natural number with $j\le b$, let $\rho$ be an $s$-finite measure on $\GT{T}$, and let
$\Phi\colon\GT{T}\to[0,\infty]$ be measurable, invariant under left translations by $\Ql{b}{T}$, and zero at every
point that is not near at~$b$. Then
\[\int\Phi(g)\;\invWeight{\cL}{j}{x}\bigl(\wl{j}{p}g\bigr)\;\dd(\etab{j}*\rho)(g)\ \le\
\prod_{a<j}\floorModulus{\cL}{a}{p_a^{-1}}\cdot\globalConst{\Gamma}\cdot\int\Phi\,\dd\rho.\]}
\leanline{\lean{lintegralfloorWeightle}\sep\lean{lintegralmulinvWeightle}}

\statement{arith}{\phantomsection\label{t:C}\textbf{Theorem C (Plancherel on a floor and on a Heisenberg group).}\enspace
In Setting~\ref{set:gauge}:

\textbf{(i)} Let $b$ be a natural number with $b\le\varrho$, let $t\in\Hl{b}{T}$ and $q,x,y\in\GT{T}$ with
$t^{-1}q\in\Nb{b}{T}$, let $j$ be a natural number with \m{\HasMatSize{j}{qx}}, and let
\m{\IsPairingKernel{\cL}{b}{k}{d}}. Put $F(\sigma)=\upperPairing{\cL}{b+1}{k}{\sigma x}{\sigma y}$. Then for
$\floorMeasure{\cL}{b}$-almost every $\sigma$ with $F(\sigma)\ne0$ the vector $t\sigma^{-1}f_b$ lies in
$\vecsOfSize{j+6}$; and for every $\xi\in A^{\times}$, every $y'\in V$, every measure $\mathsf{m}$ on $V$, natural number $i$
and $M\in[0,\infty]$ with \m{\HasFibreBound{\Gamma}{b}{\mathsf{m}}{i}{M}},
\[\begin{aligned}\int_{\vecsOfSize{i}}\Bigl|\int\psi_A\bigl(\xi\,h_T(v+y',t\sigma^{-1}f_b)\bigr)F(\sigma)\,\dd\floorMeasure{\cL}{b}(\sigma)\Bigr|^{2}&\dd\mathsf{m}(v)\\
\le M\,\Lambda(\Gamma)^{99n(j+12)}\Bigl(\frac{c_A}{\modA(\xi)}\Bigr)^{|\Omega|}&\int|F|^{2}\,\dd\floorMeasure{\cL}{b}.\end{aligned}\]

\textbf{(ii)} Let $a$ be a natural number with $a<\varrho$, $\xi,u\in A^{\times}$, $\varepsilon\in\NN$,
$C\in[0,\infty]$, \m{\IsPairingKernel{\cL}{a+1}{k}{d}}, $x,y\in\GT{T}$, $S\subseteq\GT{T}$,
$\Phi\colon\GT{T}\to[0,\infty]$ and $i,j\in\NN$. Suppose: \emph{(e)} for all $s\in\Nf{a}{T}$,
$\gL{\cL}^{-1}\bigl(\sd(dsd^{-1})-\sd(s^{\varepsilon})\bigr)=\xi\,\sd(s)$; \emph{(o)} for all $\nu\in\Nf{a}{T}$ with
$\Phi(\nu)\ne0$: $\nu\in S$ and $\Phi(\nu)\le C\,\bigl|\upperPairing{\cL}{a+1}{k}{\nu^{\varepsilon}x}{\nu y}\bigr|^{2}$;
\emph{(h)} for every $m\in S$ there are $t\in\Ql{a+1}{T}\cap\Hl{a}{T}$ and $q\in\Nb{a+1}{T}$ with
\m{\HasMatSize{i}{tqx}} such that for every $\nu\in S$ there is $\nu'\in\Nb{a}{T}$ for which the column at $e^*_a$
of the matrix $u^{-1}\bigl(\nu'\,tm(t\nu)^{-1}-1\bigr)$ lies in $\vecsOfSize{j}$. Then
\[\begin{aligned}\int\Phi\,\dd\eta_a\le{}& C\Bigl(\frac{c_A}{\modA(\xi u)}\Bigr)^{|\Omega|}\floorModulus{\cL}{a}{u}\,\Lambda(\Gamma)^{100n(i+j+31)}\\
&\cdot\int\bigl|\upperPairing{\cL}{a+2}{k}{\sigma x}{\sigma y}\bigr|^{2}\,\dd\floorMeasure{\cL}{a+1}(\sigma).\end{aligned}\]}
\leanline{\lean{amplitudesupportandplancherel}\sep\lean{parsevalheisenbergMeasureofbox}}

\sect{The setting}{s:setting}
\begin{conventions}\phantomsection\label{c:conventions}
$A$ is a commutative ring with an involution $x\mapsto\bar x$; from Definition~\ref{d:level} on it carries a topology
and a Haar measure $\dd x$ and is a \emph{place} (Section~\ref{s:space}). $\mathfrak{s}=(\varrho,\kappa,k_0)$ is a
shape, $n=\varrho+1$, $k=|\kappa|$, $r=2\varrho+k+1$. $T$ is an $\mathfrak{s}$-form over $A$; from
Lemma~\ref{l:adj} on, its split Gram matrix $\JT{T}$ is a hermitian unit. $\Omega$ is a finite set~$\Omega$,
$\mathfrak{m}$ a natural number, $\cL$ a level in $\Levels{\Omega}{T}{\mathfrak{m}}$. $\mathcal{S}$ is a size in
$\Sizes{T}$; the predicates $\operatorname{sz}$, $\operatorname{Hd}$ and the sets $\mathrm{Vec}_i$ refer to a size
that the Lean text takes as an implicit argument, and every statement names it. Index sets of vectors and matrices
are finite sets. Integrals of functions with values in $[0,\infty]$ are Lebesgue integrals; a product $0\cdot\infty$
is $0$; subtraction of natural numbers is truncated at $0$. ``Invariant'' means invariant under left translations
unless said otherwise; a measure is \emph{carried} by a set if the complement is a null set.
\end{conventions}

\subhead{The space and its groups (background library, units 112 and 119).}\phantomsection\label{s:space} A
\emph{shape} is a triple $\mathfrak{s}=(\varrho,\kappa,k_0)$ of a natural number, a finite set and an element of it.
The free module $V=A^{r}$ has the basis $e_0,\dots,e_{\varrho-1}$, $x_i$ ($i\in\kappa$), $z$,
$e^*_{\varrho-1},\dots,e^*_0$. An \emph{$\mathfrak{s}$-form} $T$ is a square matrix over $A$ on the indices of the
$x_i$ and of $z$; its \emph{split Gram matrix} $\JT{T}$ has the entries $1$ at $(e_a,e^*_a)$ and $(e^*_a,e_a)$, the
block $T$ on $V_\varrho=\langle x_i,z\rangle$ and $0$ elsewhere. It is a \emph{hermitian unit} if
$\JT{T}^{*}=\JT{T}$ and $\det\JT{T}$ is a unit; $X^{*}$ is the conjugate transpose. We write $h_T(v,w)=v^{*}\JT{T}w$,
and $\GT{T}$ is the group of the invertible matrices $g$ with $g^{*}\JT{T}g=\JT{T}$. Each basis vector has a
\emph{floor}: $a$ for $e_a$ and $e^*_a$, $\varrho$ for $z$, $n$ for the $x_i$; and a \emph{rank in the flag}: $a$ for
$e_a$, $\varrho$ for $z$ and the $x_i$, $2\varrho-a$ for $e^*_a$. $V_l$ is the span of the basis vectors of floor at
least $l$; it has rank $k+1+2(\varrho-l)$ for $l\le\varrho$. The \emph{frame} is $f_0,\dots,f_\varrho$, $f_a=e_a$
($a<\varrho$), $f_\varrho=z$; $\ZT{T}{g}=(h_T(f_a,gf_b))_{a,b}$ is the frame Gram matrix. The subgroups of $\GT{T}$:
\begin{itemize}
\item $\Hl{l}{T}$: the $g$ whose matrix agrees with the identity in every row and column of floor less than $l$;
\item $\Ql{l}{T}$: the $g$ with $g_{\alpha\beta}=\delta_{\alpha\beta}$ whenever the rank of $\beta$ is at most that of
  $\alpha$ and the larger of the floor and the rank of $\beta$ is less than $l$; for $l\le n$ this says
  $gf_a-f_a\in\langle e_0,\dots,e_{a-1}\rangle$ for $a<l$;
\item $\UT{T}$: the $g$ with $g_{\alpha\beta}=\delta_{\alpha\beta}$ whenever the rank of $\beta$ is at most that of
  $\alpha$; and $\Nb{l}{T}$: the same condition for the ranks in which every basis vector of floor at least $l$ gets
  the rank $\varrho$; so $\Nb{l}{T}=\UT{T}$ for $l\ge\varrho$;
\item $\Nf{l}{T}=\Hl{l}{T}\cap\Nb{l+1}{T}$.
\end{itemize}
For $u\in\GT{T}$, $\sd(u)=\sum_{a<\varrho}u_{f_a,f_{a+1}}$. For $p\colon\NN\to A^{\times}$ the \emph{Weyl element}
$\wl{l}{p}$ maps $e_a\mapsto p_ae^*_a$, $e^*_a\mapsto\bar p_a^{-1}e_a$ for $a<\min(l,\varrho)$ and fixes the other
basis vectors; $\wo{a}{c}$ ($a<\varrho$, $c\in A^{\times}$) does this for the one pair $(e_a,e^*_a)$. If
$\tau_A\in A$ has $\tau_A+\bar\tau_A=1$, then for $a<\varrho$, a vector $v\in V_{a+1}$ and a skew $\beta\in A$
($\bar\beta=-\beta$) the \emph{Heisenberg element} $\mathrm{n}_a(v,\beta)\in\Nf{a}{T}$ is the isometry that maps
$e^*_a\mapsto e^*_a+v+(\beta-\tau_Ah_T(v,v))e_a$, $w\mapsto w-h_T(v,w)e_a$ for $w\in V_{a+1}$, and fixes $e_a$ and the
basis vectors of floor less than~$a$.

\emph{Facts used} (all from unit 119): $\Hl{l}{T}$ and $\Ql{l}{T}$ decrease with $l$;
$\Hl{l}{T}\subseteq\Ql{l}{T}$; $\Nb{l}{T}\subseteq\UT{T}\subseteq\Ql{l'}{T}$; $\Hl{l}{T}=1$ for $l\ge\varrho+2$
and $\Ql{\varrho+2}{T}=\UT{T}$; $\Hl{l}{T}$ normalises $\Nb{l}{T}$ and $\Hl{l+1}{T}$ normalises $\Nf{l}{T}$;
\begin{enumerate}[label=\textup{(F\arabic*)}]
\item every $s\in\Ql{l}{T}$ is $\gamma\nu$ with $\gamma\in\Hl{l}{T}$, $\nu\in\Nb{l}{T}$; every
  $s\in\Ql{l+1}{T}\cap\Hl{l}{T}$ is $g\nu$ with $g\in\Hl{l+1}{T}$, $\nu\in\Nf{l}{T}$; every element of $\Nf{a}{T}$,
  $a<\varrho$, is a Heisenberg element;
\item for $b\le\varrho$ and $s\in\Hl{b}{T}$: $s\in\Ql{b+1}{T}$ if and only if $s^{-1}f_b=f_b$;
\item for $s\in\Ql{j}{T}$, $a<j$ and $v$ with $v_{e^*_b}=0$ for $b<a$: $(sv)_{e^*_a}=v_{e^*_a}$; the row of $s$ at
  $e^*_a$ agrees with that of the identity outside the columns $e^*_b$, $b<a$;
\item for $g\in\Nb{j}{T}$ and $v$ with $v_{e^*_b}=0$ for $b<j$: $(gv)_\alpha=v_\alpha$ at every $\alpha$ of floor at
  least $j$ or of rank more than $\varrho$;
\item $\wl{l}{p}$ commutes with $\Hl{l}{T}$; $\wl{l+1}{p}=\wo{l}{p_l}\wl{l}{p}$; $\wl{l}{\bar p}=\wl{l}{p}^{-1}$;
\item for $s\in\Hl{a+1}{T}$: $s\,\mathrm{n}_a(v,\beta)s^{-1}=\mathrm{n}_a(sv,\beta)$; the product and the inverse of
  Heisenberg elements are Heisenberg elements, with
  $\mathrm{n}_a(v,\beta)^{-1}=\mathrm{n}_a(-v,\beta')$ for a skew $\beta'$; and
  $\sd(s\,\mathrm{n}_a(v,\beta)s^{-1})=-h_T(v,s^{-1}f_{a+1})$;
\item the elements of $\GT{T}$ preserve $h_T$, and $\JT{T}g^{-1}=g^{*}\JT{T}$.
\end{enumerate}

\subhead{A place (unit 113).} A \emph{place} is a commutative ring $A$ with involution that is a locally compact
Polish topological ring, with a Haar measure $\dd x$ on $A$ and: an element $\tau_A$ with $\tau_A+\bar\tau_A=1$; a
unit $\delta_A$ with $\bar\delta_A=-\delta_A$; a Haar measure $\nu_A$ on the group $A^{-}$ of the skew elements such
that $(\beta,\beta')\mapsto\tau_A\delta_A^{-1}\beta+\beta'$ carries $\nu_A\times\nu_A$ to $\dd x$; a measurable
multiplicative $\modA\colon A\to[0,\infty]$ with $\modA(\bar x)=\modA(-x)=\modA(x)$ such that $v\mapsto Mv$ carries
$\dd v$ on $A^{N}$ to $\modA(\det M)^{-1}\dd v$ whenever $\det M$ is a unit; a continuous additive character
$\psi_A$ with $|\psi_A|=1$; a constant $c_A\in[0,\infty]$ bounding the $L^1$ norm of the Fourier transform of some
continuous integrable $G\ge0$ with $G(0)=1$; almost every element of $A$ is a unit; and $\modA(u)^{-1}\le|g(u)|^2$ on
units for some locally integrable~$g$.

\subhead{Quadric measures and hyperbolic bases (units 113 and 119).} For a hermitian matrix $T'$ on a finite index
set $\iota$ and a skew $\beta$, the \emph{flat chart} of the quadric $h=\delta_A^{-1}\beta$ in the space
$A^2\oplus A^{\iota}$ with the form (hyperbolic plane) $\oplus\,T'$ is
$(t,w,\alpha)\mapsto\bigl(\bar t^{-1}\tau_A(\delta_A^{-1}\beta-w^{*}T'w)+t\alpha,\ t,\ w\bigr)$, for $t\in A$,
$w\in A^{\iota}$, $\alpha\in A^{-}$, and the \emph{quadric measure} is the image of $\dd t\,\dd w\,\dd\nu_A(\alpha)$
under it. A \emph{hyperbolic basis} of $V_\varrho$, with parameters $(\Omega,T,\mathfrak{m})$, is a matrix $E$ whose
columns are $\mathfrak{m}+1$ hyperbolic pairs and vectors indexed by $\Omega$ with $E^{*}TE$ the Gram matrix of
$\mathfrak{m}+1$ hyperbolic planes plus the identity on $\Omega$, with $k+1=|\Omega|+2\mathfrak{m}+2$ and
$\modA(\det T)=1$. For $b\le\varrho$ its \emph{floor basis} $E_b$ is the basis of $V_b$ made of $E$ and of the pairs
$(e_a,e^*_a)$, $b\le a<\varrho$; its coordinates are: the first hyperbolic pair, a first half of the other pairs
(indexed by a set $\pi_b$ of $\mathfrak{m}+\varrho-b$ elements), the orthonormal part $\Omega$, and the second half.
The \emph{floor quadric measure} $\quadric{\cL}{b}$ of a level is the image under $E_b$ of the quadric measure with
$\beta=\delta_Ah_T(f_b,f_b)$ in these coordinates, and $\vol_b$ is the Haar measure of $V_b$, the image of the volume
of the coordinates of floor at least $b$.

\subhead{Sizes.}
\begin{definition}[size]\label{d:size}
Let $A$ be a commutative ring with involution, $\mathfrak{s}$ a shape and $T$ an $\mathfrak{s}$-form over $A$. A
\emph{size} for $T$, an element of $\Sizes{T}$, is a family $\mathcal{S}=(\mathcal{S}_i)_{i\in\NN}$ of subsets of
$A$ such that: $\mathcal{S}_i\mathcal{S}_{i'}\subseteq\mathcal{S}_{i+i'}$; the sum of a family of elements of
$\mathcal{S}_i$ indexed by a finite set of at most $r+1$ elements lies in $\mathcal{S}_{i+1}$; $-x\in\mathcal{S}_i$
if and only if $x\in\mathcal{S}_i$, and $\bar x\in\mathcal{S}_i$ if and only if $x\in\mathcal{S}_i$;
$0,1\in\mathcal{S}_i$ for all $i$; and all entries of $\JT{T}$ and of $\JT{T}^{-1}$ lie in $\mathcal{S}_1$.
\end{definition}
\leanline{\lean{Size}}
\begin{definition}[vectors and matrices of a size; heads]\label{d:sz}
Let $\mathcal{S}$ be a size in $\Sizes{T}$ (an implicit argument of the three notions). For a natural number $i$,
$\vecsOfSize{i}$ is the set of the vectors $v$, on a given finite index set, with all coordinates in $\mathcal{S}_i$.
A matrix $X$ over $A$ \emph{has size at most $i$}, \m{\HasMatSize{i}{X}}, if all its entries lie in
$\mathcal{S}_i$; for $g\in\GT{T}$ this refers to the matrix of $g$. A sequence $p\colon\NN\to A^{\times}$ is a
\emph{head of length $j$}, \m{\IsHead{j}{p}}, if $(p_0p_1\cdots p_{l-1})^{-1}\in\mathcal{S}_1$ for every $l\le j$.
\end{definition}
\leanline{\lean{vecsOfSize}\sep\lean{HasMatSize}\sep\lean{IsHead}}
\begin{definition}[Weyl ratio]\label{d:wr}
For a natural number $j$ and $p,q\colon\NN\to A^{\times}$, the element $\weylRatio{j}{p}{q}$ of $\GT{T}$ is the
product of the inverse of the Weyl element of $q$ and the Weyl element of $p$, both of length~$j$.
\end{definition}
\leanline{\lean{weylRatio}\sep\lean{coeweylRatio}\sep\lean{weylRatioinv}}
Its matrix is diagonal: the entry is $p_a/q_a$ at $e_a$ and $\overline{q_a/p_a}$ at $e^*_a$ for
$a<\min(j,\varrho)$, and $1$ elsewhere; and $(\weylRatio{j}{p}{q})^{-1}=\weylRatio{j}{q}{p}$.

\subhead{Levels and the tower measure.}
\begin{definition}[level]\label{d:level}
Let $A$ be a place, $T$ an $\mathfrak{s}$-form whose split Gram matrix is a hermitian unit, $\Omega$ a finite set
and $\mathfrak{m}$ a natural number. A \emph{level}, an element $\cL$ of $\Levels{\Omega}{T}{\mathfrak{m}}$, consists
of: a hyperbolic basis $E(\cL)$ of $V_\varrho$ with parameters $(\Omega,T,\mathfrak{m})$; two $s$-finite measures
$\theta$ and $\mu_H$ on $\GT{T}$ with $\mu_H\ne0$; a unit $\gL{\cL}$ of $A$ with $\overline{\gL{\cL}}=\gL{\cL}$; a set
$\Ball{\cL}\subseteq V$, the \emph{ball}; and two sets $\BoxL{\cL},\Posn{\cL}\subseteq\GT{T}$, the \emph{box} and the
\emph{positions}; such that
\begin{enumerate}[label=\textup{(L\arabic*)}]
\item $\theta$ is carried by $\Hl{\varrho}{T}$, and its image under $\sigma\mapsto\sigma^{-1}z$ is the floor quadric
  measure $\quadric{\cL}{\varrho}$;
\item $\mu_H$ is carried by $\Hl{n}{T}$; $\mu_H*\delta_s=\mu_H$ for $s\in\Hl{n}{T}$; and
  $\mu_H\{h:u\,h\,g\in\BoxL{\cL}\text{ for some }u\in\UT{T}\}\le1$ for every $g\in\GT{T}$;
\item $\Ball{\cL}$ is open, contains $0$ and has compact closure; $\BoxL{\cL}$ is compact; $\Posn{\cL}$ is
  measurable.
\end{enumerate}
\end{definition}
\leanline{\lean{Level}}
\begin{definition}[floor measures]\label{d:fm}
For a level $\cL$ in $\Levels{\Omega}{T}{\mathfrak{m}}$ and a natural number $l$, the measure
$\floorMeasure{\cL}{l}$ on $\GT{T}$ is: $\varpi_l*\eta_l$ for $l<\varrho$; $\theta$ for $l=\varrho$; $\mu_H$ for
$l=n$; and the unit mass $\delta_1$ for $l>n$. Here $*$ is the convolution of measures, $\varpi_l$ (the \emph{Weyl
law}) is the image of $\modA(c)^{k+1+2(\varrho-(l+1))}\dd c$ under $c\mapsto\wo{l}{c}^{-1}$, and $\eta_l$ (the
\emph{Heisenberg measure}) is the image of $\dd v\,\dd\nu_A(\beta)$ on (coordinates of $V_{l+1})\times A^{-}$ under
$(v,\beta)\mapsto\mathrm{n}_l(v,\beta)$; we put $\eta_l=\delta_1$ for $l\ge\varrho$ and
$\etab{j}=\eta_0*\eta_1*\dots*\eta_{j-1}$.
\end{definition}
\leanline{\lean{floorMeasure}\sep\lean{floorMeasuretop}}
\begin{definition}[stack, group and upper measures]\label{d:tower}
For a level $\cL$ in $\Levels{\Omega}{T}{\mathfrak{m}}$, $\stackMeasure{\cL}$ is the family of measures
\[\stackMeasure{\cL}_{a,i}=(\eta_{a+i-1}*\floorMeasure{\cL}{a+i-1})*\dots*(\eta_{a+1}*\floorMeasure{\cL}{a+1})*(\eta_a*\floorMeasure{\cL}{a})\]
($i$ factors; $\delta_1$ for $i=0$), indexed by natural numbers $a,i$. The \emph{tower measure} is
$\groupMeasure{\cL}=\stackMeasure{\cL}_{0,\varrho+2}$, and for a natural number $l$,
$\upperMeasure{\cL}{l}=\floorMeasure{\cL}{\varrho+1}*\dots*\floorMeasure{\cL}{l}$ ($\varrho+2-l$ factors; $\delta_1$
for $l\ge\varrho+2$).
\end{definition}
\leanline{\lean{stackMeasure}\sep\lean{groupMeasure}\sep\lean{upperMeasure}\sep\lean{upperMeasureeqsuccmconv}\sep\lean{aestackMeasuremem}}
So $\upperMeasure{\cL}{l}=\upperMeasure{\cL}{l+1}*\floorMeasure{\cL}{l}$, and $\stackMeasure{\cL}_{l,i}$ is carried
by $\Hl{l}{T}$ (Lemma~\ref{l:carried} and \Inp{I2}).

\subhead{Nearness, floor values, weights.}
\begin{definition}[near points]\label{d:near}
For a level $\cL$ in $\Levels{\Omega}{T}{\mathfrak{m}}$, a natural number $j$ and $g\in\GT{T}$: $g$ is \emph{near at
$j$}, \m{\IsNear{\cL}{j}{g}}, if $\xi g\in\BoxL{\cL}$ for some $\xi\in\Ql{j}{T}$.
\end{definition}
\leanline{\lean{IsNear}\sep\lean{isNearmuliff}\sep\lean{offloorGroupmul}}
Hence nearness at $j$ is invariant under $\Ql{j}{T}$; and if $xg$ is near at $i$, $x\in\Hl{j}{T}$ and $j\le i$, then
$g$ is near at $j$, because $\Ql{i}{T}\,x\subseteq\Ql{j}{T}$.
\begin{definition}[orthonormal part, close pairs, floor values]\label{d:val}
Let $\cL$ be a level in $\Levels{\Omega}{T}{\mathfrak{m}}$. Its \emph{orthonormal part} $\orthonormalPart{\cL}$ is
the matrix whose columns are the vectors of $E(\cL)$ indexed by $\Omega$, as vectors of $V$. The set
$\closePairs{\cL}$ consists of the pairs $(u_1,u_2)$ of vectors of $V$ with $u_1\in\Ball{\cL}$ and
$\gL{\cL}^{-1}(u_1-u_2)\in\Ball{\cL}$. For a natural number $a<\varrho$ and $x,y\in\GT{T}$, the set of \emph{floor
values} is the subset of $A$
\[\floorValues{\cL}{a}{x}{y}=\bigl\{w_{e^*_a}:\ w\in V,\ (x^{-1}w,(xy)^{-1}w)\in\closePairs{\cL},\ w_{e^*_b}=0\text{ for all }b<a\bigr\}.\]
\end{definition}
\leanline{\lean{orthonormalPart}\sep\lean{closePairs}\sep\lean{floorValues}}
\begin{definition}[volume exponent, weights, floor module]\label{d:weight}
For a level $\cL$ in $\Levels{\Omega}{T}{\mathfrak{m}}$: the \emph{volume exponent} is the natural number
$\volumeExponent{\cL}=|\Omega|-(3n+1+\mathfrak{m})$; for a natural number $j$ and $x,y\in\GT{T}$ the \emph{weight} is
\[\floorWeight{\cL}{j}{x}{y}=\prod_{a<\min(j,\varrho)}\vol\bigl(\floorValues{\cL}{a}{x}{y}\bigr)^{\volumeExponent{\cL}}\in[0,\infty];\]
for $x\in\GT{T}$, $\invWeight{\cL}{j}{x}$ is the function on $\GT{T}$ that is $z\mapsto\floorWeight{\cL}{j}{z}{x}^{-1}$
at the points $z$ near at $j$ and $0$ elsewhere; and for a natural number $l$ and $g\in A^{\times}$ the \emph{floor
module} is $\floorModulus{\cL}{l}{g}=\modA(g)^{k+1+2(\varrho-(l+1))}\cdot\modA(g)^{1/2}$. In the Lean text the last
argument of the floor module has the default value $\gL{\cL}$. (The census lists, among the arguments of the
inverse weight, the hypothesis that the tower measure is a Haar measure; the definition does not use it.)
\end{definition}
\leanline{\lean{volumeExponent}\sep\lean{floorWeight}\sep\lean{invWeight}\sep\lean{floorModulus}\sep\lean{floorWeightsucc}\sep\lean{floorWeightnezeroandnetop}\sep\lean{floorWeightmulleft}\sep\lean{measurablefloorWeight}\sep\lean{measurableinvWeight}}
\begin{definition}[mass of a floor]\label{d:mass}
For a level $\cL$ in $\Levels{\Omega}{T}{\mathfrak{m}}$, a natural number $j$, $K\in[0,\infty]$ and a function
$c\colon\GT{T}\to[0,\infty]$: \m{\FloorMassLe{\cL}{j}{K}{c}} means that for every $g\in\GT{T}$
\[\int_{\{\sigma:\ \sigma g\text{ is near at }j+1\}}c(\sigma g)\,\dd\floorMeasure{\cL}{j}(\sigma)\le K.\]
\end{definition}
\leanline{\lean{FloorMassLe}}

\subhead{Kernels averaged over floors.}
\begin{definition}[averaged kernel, pivot]\label{d:av}
For a level $\cL$ in $\Levels{\Omega}{T}{\mathfrak{m}}$, a natural number $j$, a function
$k\colon\GT{T}\times\GT{T}\to\CC$ and $g,y\in\GT{T}$,
\[\begin{gathered}\upperPairing{\cL}{j}{k}{g}{y}=\int k(\sigma g,\sigma y)\,\dd\upperMeasure{\cL}{j}(\sigma),\\
\pivotModulus{\cL}{j}{k}{g}{y}=\int\bigl|\upperPairing{\cL}{j+1}{k}{\sigma g}{\sigma y}\bigr|\,\dd\floorMeasure{\cL}{j}(\sigma).\end{gathered}\]
The first is a Bochner integral, $0$ if the integrand is not integrable; the second lies in $[0,\infty]$. (The
census lists, among the arguments of the pivot, the hypothesis that the tower measure is a Haar measure; the
definition does not use it.)
\end{definition}
\leanline{\lean{upperPairing}\sep\lean{pivotModulus}}
\begin{definition}[pairing kernel]\label{d:pk}
For a level $\cL$ in $\Levels{\Omega}{T}{\mathfrak{m}}$, a natural number $j$, $k\colon\GT{T}\times\GT{T}\to\CC$ and
$w\in\GT{T}$: the pair is a \emph{pairing kernel at $j$}, \m{\IsPairingKernel{\cL}{j}{k}{w}}, if
\begin{enumerate}[label=\textup{(K\arabic*)}]
\item $k$ is measurable;
\item for all $l$, all $\nu,\nu'\in\Nb{l}{T}$ and all $g,y$:
  \[\upperPairing{\cL}{l+1}{k}{\nu g}{\nu'y}=\psi_A\bigl(\gL{\cL}^{-1}(\sd(w\nu'w^{-1})-\sd(\nu))\bigr)\,\upperPairing{\cL}{l+1}{k}{g}{y};\]
\item $wsw^{-1}=s$ for all $s\in\Hl{j}{T}$;
\item for every $l$ there is $c<\infty$ with $\int|k(\sigma g,\sigma y)|\,\dd\upperMeasure{\cL}{l}(\sigma)\le c$ for
  all $g,y$;
\item for all $l,g,y$: if $\upperPairing{\cL}{l}{k}{g}{y}\ne0$ then $g$ is near at $l$.
\end{enumerate}
\end{definition}
\leanline{\lean{IsPairingKernel}}
\begin{definition}[weighted squares, self-norm]\label{d:sn}
For a level $\cL$ in $\Levels{\Omega}{T}{\mathfrak{m}}$, $k\colon\GT{T}\times\GT{T}\to\CC$, natural numbers $j$ and
$x,g\in\GT{T}$:
\[\weightedSquare{\cL}{j}{k}{x}{g}=\floorWeight{\cL}{j}{g}{x}\cdot\bigl|\upperPairing{\cL}{j+1}{k}{g}{gx}\bigr|^{2};\]
the \emph{weighted square sum} is defined by recursion on $n-j$:
\[\weightedSquareSum{\cL}{k}{x}{j}{g}=\int\Bigl(\weightedSquare{\cL}{j}{k}{x}{\sigma g}+\weightedSquareSum{\cL}{k}{x}{j+1}{\sigma g}\Bigr)\dd\floorMeasure{\cL}{j}(\sigma)\]
for $j<n$ and $\weightedSquareSum{\cL}{k}{x}{j}{g}=0$ for $j\ge n$; and the \emph{self-norm} is
\[\selfNorm{\cL}{j}{k}{x}=\int\weightedSquareSum{\cL}{k}{x}{j}{g}\,\dd(\floorMeasure{\cL}{j-1}*\dots*\floorMeasure{\cL}{0})(g).\]
(The census lists, among the arguments of the first two notions, the hypothesis that the tower measure is a Haar
measure; the definitions do not use it.)
\end{definition}
\leanline{\lean{weightedSquare}\sep\lean{weightedSquareSum}\sep\lean{selfNorm}}

\subhead{Gauges.}
\begin{definition}[gauge]\label{d:gauge}
Let $\cL$ be a level in $\Levels{\Omega}{T}{\mathfrak{m}}$. A \emph{gauge} on $\cL$, an element $\Gamma$ of
$\Gauges{\cL}$, is a size $(\Gamma_i)_{i\in\NN}$ for $T$ together with a constant $\Lambda(\Gamma)\in[0,\infty]$
such that, with sizes taken for $\Gamma$ and $\Lambda=\Lambda(\Gamma)$:
\begin{enumerate}[label=\textup{(G\arabic*)}]
\item every $\Gamma_j$ is compact; $2\le\Lambda<\infty$; $\mathfrak{m}\le9n$; $c_A\ge1$;
\item $\gL{\cL}\in\Gamma_1$; \m{\HasMatSize{1}{E(\cL)}}; $\tau_A\delta_A^{-1}\in\Gamma_1$;
\item \m{\HasMatSize{1}{x}} for every $x\in\BoxL{\cL}$; $\Ball{\cL}\subseteq\vecsOfSize{1}$;
\item $\modA(u)\le\Lambda^{j}$ for every unit $u\in\Gamma_j$; $\vol(\Gamma_j)\le\Lambda^{j+1}$;
\item \emph{(cone)} for every finite set $\pi$ with $\mathfrak{m}\le|\pi|\le\mathfrak{m}+\varrho$ and every $j$, the
  quadric measure with $\beta=0$ of the sum of $|\pi|+1$ hyperbolic planes gives $\vecsOfSize{j}$ mass at most
  $\Lambda^{99n(j+1)}$;
\item \emph{(volumes)} for every $x\in\BoxL{\cL}$, $y\in\Posn{\cL}$ and every matrix $v$ over $A$ with $\varrho$
  rows and columns indexed by $\Omega$, let
  \[W=\{w\in A^{\Omega}:vw=0,\ (x^{-1}\orthonormalPart{\cL}w,(xy)^{-1}\orthonormalPart{\cL}w)\in\closePairs{\cL}\};\]
  then there is a measurable $\varphi\colon A^{\Omega}\to[0,\infty]$ such that for every $u\in A^{\Omega}$
  \[\vol\{u^{*}w:w\in W\}^{-1}\le\varphi(u),\qquad\vol\,\mathcal{C}_u\le\varphi(u),\]
  where $\mathcal{C}_u$ is the set of the $c\in A$ such that for every $w\in W$ there are
  $z,z'\in\closePairs{\cL}$ with $(cu)^{*}w=\gL{\cL}^{-1}\bigl(h_T(z_1,z'_1)-h_T(z_2,z'_2)\bigr)$; and
  $\int_{\vecsOfSize{j}}\varphi(u)^{\volumeExponent{\cL}}\dd u\le\Lambda^{99n(j+1)}$ for every $j$.
\end{enumerate}
\end{definition}
\leanline{\lean{Gauge}\sep\lean{oneleLambda}}
In particular $\Lambda(\Gamma)\ge1$.
\begin{definition}[positions of size $1$; constants]\label{d:consts}
For a level $\cL$ in $\Levels{\Omega}{T}{\mathfrak{m}}$ and a gauge $\Gamma$ in $\Gauges{\cL}$,
$\positionsOfSize{\cL}{\Gamma}$ is the set of the $x\in\Posn{\cL}$ with \m{\HasMatSize{1}{\gL{\cL}^{-1}(x-1)}} for
$\Gamma$; for a natural number $l$,
$\floorConst{\Gamma}{l}=\floorModulus{\cL}{l}{\gL{\cL}}\cdot c_A^{|\Omega|}\cdot\Lambda(\Gamma)^{4600n}$; and
$\globalConst{\Gamma}=\Lambda(\Gamma)^{10^{6}n^{3}}$.
\end{definition}
\leanline{\lean{positionsOfSize}\sep\lean{floorConst}\sep\lean{globalConst}\sep\lean{measurableSetpositionsOfSize}\sep\lean{measurableSetvecsOfSize}}
The sets $\vecsOfSize{j}$ (for a gauge) and $\positionsOfSize{\cL}{\Gamma}$ are measurable, by compactness of the
$\Gamma_j$.
\begin{definition}[fibre bound]\label{d:fib}
Let $\cL$ be a level in $\Levels{\Omega}{T}{\mathfrak{m}}$, $\Gamma$ a gauge in $\Gauges{\cL}$, $b$ a natural number
with $b\le\varrho$, $q$ a measure on $V$, $j$ a natural number and $M\in[0,\infty]$. Write a coordinate vector for
the floor basis $E_b$ as $(t',t,\pi',w,\pi'')$: the first hyperbolic pair $(t',t)$, the first half $\pi'$ of the
other pairs, the orthonormal part $w\in A^{\Omega}$, the second half $\pi''$; call $P=(t,\pi')$ the \emph{fixed}
coordinates. Then \m{\HasFibreBound{\Gamma}{b}{q}{j}{M}} means: there are a measurable space $Z$ with an $s$-finite
measure $\zeta$, a measure $\rho'$ on coordinate vectors with $E_b\rho'=q$, a map $\mathsf{k}$ from triples $(P,w,z)$,
$z\in Z$, to coordinate vectors that carries $\dd P\,\dd w\,\dd\zeta(z)$ to $\rho'$ and whose values have the fixed
coordinates $P$ and the orthonormal coordinates $w$, and a measurable $\mu$ with
\[\zeta\{z:\mathsf{k}(P,w,z)\in\vecsOfSize{j+4}\}\le\mu(P)\ \text{ for all }P,w,\qquad\int\mu(P)\,\dd P\le M.\]
\end{definition}
\leanline{\lean{HasFibreBound}}
\sect{What is taken for granted}{s:inputs}
Besides the facts \textup{(F1)}--\textup{(F7)} about the groups, the proofs use the following results of the
background library. None is a member of this unit; the unit that owns each is given.

\Inp{I1} \emph{Convolution} (unit 119). For $s$-finite measures $\mu,\nu$ on $\GT{T}$ and measurable $f\ge0$:
$\int f\,\dd(\mu*\nu)=\iint f(xy)\,\dd\mu(x)\,\dd\nu(y)$; $\mu*\delta_s$ and $\delta_s*\mu$ are the right and left
translates of $\mu$; if $\delta_s*\nu=\nu*\delta_s$ for $\mu$-almost every $s$, then $\mu*\nu=\nu*\mu$; if
$\delta_s*\mu=\mu*\delta_t$, then $\int f\,\dd\mu=\int f(s^{-1}xt)\,\dd\mu(x)$.

\Inp{I2} (unit 119) A convolution of measures carried by a measurable subgroup is carried by it; the groups
$\Hl{l}{T}$, $\Ql{l}{T}$, $\Nb{l}{T}$ are closed (unit 112).

\Inp{I3} \emph{Heisenberg measures} (unit 119). $\eta_l$ is carried by $\Nf{l}{T}$;
$\eta_l*\delta_s=\eta_l=\delta_s*\eta_l$ for $s\in\Nf{l}{T}$; $\delta_s*\eta_l=\eta_l*\delta_s$ for
$s\in\Hl{l+1}{T}$. $\etab{j}$ is carried by $\Nb{j}{T}$; $\etab{j}*\delta_s=\etab{j}=\delta_s*\etab{j}$ for
$s\in\Nb{j}{T}$; $\delta_s*\etab{j}=\etab{j}*\delta_s$ for $s\in\Hl{j}{T}$; and $\etab{j+1}=\eta_j*\etab{j}$.

\Inp{I4} \emph{Disintegration} (units 113, 119). Call a family $(\lambda_\beta)$ of measures on a space with a
hermitian form $h$, indexed by the skew $\beta$, a \emph{disintegration} of a measure $\mu$ if it is weakly
continuous in $\beta$, finite on compact sets, and
$\int f(h(x,x),x)\,\dd\mu=\iint f(\delta_A^{-1}\beta,x)\,\dd\lambda_\beta(x)\,\dd\nu_A(\beta)$ for measurable $f\ge0$.
Then: (a) the quadric measures of the flat charts disintegrate the volume of $A^2\oplus A^{\iota}$; (b) the image of
a disintegration under a linear map with a left inverse that respects the forms is a disintegration of the image
measure; (c) a disintegration is unique; (d) the floor quadric measures of a hyperbolic basis disintegrate $\vol_b$
along $h_T$, the floor basis $E_b$ carries the volume of its coordinates to $\vol_b$, and $E_b^{*}\JT{T}E_b$ is the
Gram matrix of hyperbolic planes plus the identity on $\Omega$; (e) for $s\in\Hl{l}{T}$ the action of $s$ on $V_l$
preserves $\vol_l$. The catalogue of the write-up classes (a) as a standard local computation (Leray measures on
quadrics); reference to be confirmed.

\Inp{I5} \emph{Functions on a floor} (unit 119). A measurable function on $\GT{T}$ with values in a standard Borel
space that is invariant under $\Ql{b+1}{T}\cap\Hl{b}{T}$, $b\le\varrho$, is on $\Hl{b}{T}$ a measurable function of
$\sigma^{-1}f_b$.

\Inp{I6} \emph{Plancherel's inequality at a place} (unit 119). For $\xi\in A^{\times}$, a finite set $\iota$ and
$f\colon A^{\iota}\to\CC$ measurable and integrable:
$\int|\int\psi_A(\xi\,v^{*}w)f(w)\,\dd w|^2\dd v\le(c_A/\modA(\xi))^{|\iota|}\int|f|^2$. The second reader of the
artifact relates it to J.~Tate, \emph{Fourier analysis in number fields and Hecke's zeta-functions} (1950), \S2.2, as
a weaker form: a one-sided consequence of abelian Plancherel. Tate's \S2.2 covers local fields; for a general place
(a locally compact ring) the second reader's record points to the abelian Plancherel theorem (H.~Cartan,
R.~Godement, \emph{Th\'eorie de la dualit\'e et analyse harmonique dans les groupes ab\'eliens localement compacts},
Ann.\ Sci.\ \'Ec.\ Norm.\ Sup.\ 64 (1947)) together with an argument on the normalisation; the place of a reference
for the general case is not confirmed, and neither reference was opened for this note.

\Inp{I7} \emph{Skew elements} (unit 113). If $S\subseteq A^{-}$, $u\in A^{\times}$ and $C\subseteq A$ are such that
$u(\tau_A\delta_A^{-1}\beta+\beta')\in C$ for all differences $\beta,\beta'$ of elements of $S$, and
$\vol(C)/\modA(u)\le t^2$, then $\nu_A(S)\le t$.

\Inp{I8} The \emph{Cauchy--Schwarz inequality} for integrals of functions with values in $[0,\infty]$ (unit 117):
if $h^2\le f_1f_2$ pointwise then $(\int h)^2\le\int f_1\int f_2$.

\Inp{I9} Laplace expansion of a determinant along a row; the determinant of a triangular matrix (unit 112).

\sect{Sizes}{s:sizes}
In this section $\mathcal{S}$ is a size in $\Sizes{T}$, and sizes are taken for $\mathcal{S}$; index sets of rows
may be any finite sets, and vectors and columns are indexed by the basis of $V$.

\begin{lemma}[rules]\label{l:rules}
Let $\mathcal{S}$ be a size in $\Sizes{T}$, $i,j$ natural numbers, $\alpha,\beta\in A$, $v,w\in V$, and $X,Z$
matrices with columns indexed by the basis of $V$, $Y$ a matrix with rows indexed by it. Then:
\textup{(a)} if $\alpha\in\mathcal{S}_i$ and $i\le j$ then $\alpha\in\mathcal{S}_j$; if \m{\HasMatSize{i}{X}} and
$i\le j$ then \m{\HasMatSize{j}{X}};
\textup{(b)} if $\alpha,\beta\in\mathcal{S}_i$ then $\alpha-\beta\in\mathcal{S}_{i+1}$; if \m{\HasMatSize{i}{X}} and
\m{\HasMatSize{i}{Z}} then \m{\HasMatSize{i+1}{X-Z}}; the identity matrix has \m{\HasMatSize{i}{1}};
\textup{(c)} if $w_\alpha\in\mathcal{S}_i$ at every index where $v_\alpha\ne0$, and $v\in\vecsOfSize{j}$, then
$\sum_\alpha w_\alpha v_\alpha\in\mathcal{S}_{i+j+1}$;
\textup{(d)} if \m{\HasMatSize{i}{X}} and $v\in\vecsOfSize{j}$ then $Xv\in\vecsOfSize{i+j+1}$; if moreover
\m{\HasMatSize{j}{Y}} then \m{\HasMatSize{i+j+1}{XY}};
\textup{(e)} if $p$ is a head of length $j$, \m{\IsHead{j}{p}}, then so is $\bar p$: \m{\IsHead{j}{\bar p}}.
\end{lemma}
\begin{proof}
(a) $\alpha=\alpha\cdot1$ with $1\in\mathcal{S}_{j-i}$. (b) $\alpha-\beta$ is a sum of two elements of
$\mathcal{S}_i$, and $2\le r+1$. (c) Each term $w_\alpha v_\alpha$ is $0$ or lies in $\mathcal{S}_{i+j}$, and there
are $r$ terms. (d) follows from (c), entry by entry. (e) The sets $\mathcal{S}_1$ are stable under the involution.
\end{proof}
\leanline{\lean{memofSizemono}\sep\lean{monoa}\sep\lean{submemofSize}\sep\lean{sub}\sep\lean{one}\sep\lean{dotProductmemofSize}\sep\lean{mulVec}\sep\lean{mul}\sep\lean{starhead}}

\begin{lemma}[determinants]\label{l:det}
Let $\mathcal{S}$ be a size in $\Sizes{T}$, $j$ a natural number with $j\le\varrho$ and $Z$ a $j\times j$ matrix
over $A$ with \m{\HasMatSize{7}{Z}}. Then $\det Z\in\mathcal{S}_{8j}$.
\end{lemma}
\begin{proof}
Induction on $j$; for $j=0$ the determinant is $1$. For $j+1$, expand along the first row \Inp{I9}: each of the
$j+1$ terms is $\pm$ an entry of $Z$ times the determinant of a $j\times j$ minor, so it lies in
$\mathcal{S}_{7+8j}$; and $j+1\le r+1$, so the sum lies in $\mathcal{S}_{8j+8}$.
\end{proof}
\leanline{\lean{detmemofSize}}

\begin{lemma}[adjoints and inverses]\label{l:adj}
Let $\mathcal{S}$ be a size in $\Sizes{T}$ and suppose that $\JT{T}$ is a hermitian unit. \textup{(i)} Let $X,Y$ be
square matrices on the basis of $V$ with $\JT{T}X=Y^{*}\JT{T}$, let $a,b$ be two indices and $i$ a natural number
such that for every index $\nu$ with $\JT{T}_{\nu b}\ne0$ the row $\nu$ of $Y$ lies in $\vecsOfSize{i}$. Then
$X_{ab}\in\mathcal{S}_{i+4}$. \textup{(ii)} If $x\in\GT{T}$ has \m{\HasMatSize{i}{x}}, then
\m{\HasMatSize{i+4}{x^{-1}}}.
\end{lemma}
\begin{proof}
(i) $X=\JT{T}^{-1}(Y^{*}\JT{T})$. The entry $(m,b)$ of $Y^{*}\JT{T}$ is $\sum_\nu\overline{Y_{\nu m}}\JT{T}_{\nu b}$,
in $\mathcal{S}_{i+2}$ by rule (c), because the column $b$ of $\JT{T}$ has size $1$; and the row $a$ of
$\JT{T}^{-1}$ has size $1$, so $X_{ab}\in\mathcal{S}_{1+(i+2)+1}$. (ii) Apply (i) to $X=x^{-1}$, $Y=x$, by (F7).
\end{proof}
\leanline{\lean{adjointapplymemofSize}\sep\lean{inv}}

\begin{lemma}[back substitution in a row]\label{l:row}
Let $\mathcal{S}$ be a size in $\Sizes{T}$, $X$ a matrix with columns indexed by the basis of $V$, $m$ one of its
rows, and $c,d,j$ natural numbers. Suppose \textup{(h)} for every $a<\min(j,\varrho)$ there is $w\in V$ with
$w_{e^*_b}=0$ for $b<a$, $w_{e^*_a}=1$, $w\in\vecsOfSize{c}$ and $Xw\in\vecsOfSize{d}$; and \textup{(z)}
$X_{mk}\in\mathcal{S}_d$ for every column $k$ of floor at least $j$ or of rank at most $\varrho$. Then for every
natural number $i$ and every column $k$ whose floor is at least $j-i$ or whose rank is at most $\varrho$,
$X_{mk}\in\mathcal{S}_{d+(c+3)i}$.
\end{lemma}
\begin{proof}
Induction on $i$; $i=0$ is (z). A column new at step $i+1$ is $k=e^*_a$ with $a=j-i-1$. Take $w$ from (h) for this
$a$ and let $w'=w-e^*_a$; then $X_{mk}=(Xw)_m-\sum_{k'}X_{mk'}w'_{k'}$. The non-zero coordinates of $w'$ are at
columns of rank at most $\varrho$ or at $e^*_{b}$ with $b>a$, where the induction hypothesis gives
$X_{mk'}\in\mathcal{S}_{d+(c+3)i}$; and $w'\in\vecsOfSize{c+1}$. By rules (c) and (b),
$X_{mk}\in\mathcal{S}_{d+(c+3)i+(c+1)+1+1}$.
\end{proof}
\leanline{\lean{rowmemofSizeofdualframe}}

\begin{lemma}[back substitution for an inverse]\label{l:backsub}
Let $\mathcal{S}$ be a size in $\Sizes{T}$, $\JT{T}$ a hermitian unit, $j$ and $c$ natural numbers,
$u\in\Nb{j}{T}$, and $r\in A^{\times}$ with $\bar r=r$ and $r\in\mathcal{S}_1$. Suppose that for every
$a<\min(j,\varrho)$ there is $w\in V$ with $w_{e^*_b}=0$ for $b<a$, $w_{e^*_a}=1$, $w\in\vecsOfSize{c}$ and
$r^{-1}(1-u)w\in\vecsOfSize{c}$. Then \m{\HasMatSize{c+4+(2c+8)j}{r^{-1}(1-u^{-1})}}.
\end{lemma}
\begin{proof}
Let $Y=r^{-1}(1-u)$ and $X=r^{-1}(1-u^{-1})$; then $\JT{T}X=Y^{*}\JT{T}$ by (F7), as $\bar r=r$. \emph{Rows of $Y$.}
If the row $\nu$ of $\JT{T}$ is not zero at a column $k$ of floor at least $j$ or of rank at most $\varrho$, then
the entries of row $\nu$ of $u$ at such columns are those of the identity (definition of $\Nb{j}{T}$), so (z) of
Lemma~\ref{l:row} holds for $Y$ with any $d$, and (h) holds by hypothesis with $d=c$; taking $i=j$, the whole row
$\nu$ of $Y$ lies in $\vecsOfSize{c+(c+3)j}$. \emph{Entries of $X$ at the columns of (z).} By
Lemma~\ref{l:adj}(i) they lie in $\mathcal{S}_{d'}$, $d'=c+(c+3)j+4$. \emph{All entries of $X$.} For $w$ as in the
hypothesis, $uw$ has the same coordinates at $e^*_b$, $b\le a$ (F3), $uw=w-r\cdot(Yw)\in\vecsOfSize{c+2}$, and
$X(uw)=-Yw\in\vecsOfSize{c}$. So Lemma~\ref{l:row} applies to $X$ with $c+2$, $d'$ and $i=j$, and gives the size
$d'+(c+5)j$.
\end{proof}
\leanline{\lean{hasMatSizeonesubinv}}

\begin{setting}[two frames near a set of size $1$]\label{set:two}
Let $\mathcal{S}$ be a size in $\Sizes{T}$ and $\JT{T}$ a hermitian unit. Let $\mathsf{K}\subseteq\GT{T}$ be a set with
\m{\HasMatSize{1}{x}} for all $x\in\mathsf{K}$, let $j$ be a natural number with $j\le\varrho$, $p\colon\NN\to A^{\times}$,
and $g,s\in\GT{T}$ such that $\eta\,\wl{j}{p}^{-1}g\in\mathsf{K}$ for some $\eta\in\Ql{j}{T}$, and such that
$s\in\Ql{j}{T}$ and $sg\in\mathsf{K}$.
\end{setting}

\begin{lemma}[minors and heads]\label{l:minor}
In Setting~\ref{set:two}: \textup{(i)} for every $a<j$ and every index $i$ of the basis,
$(p_0p_1\cdots p_a)\cdot s_{i,e^*_a}\in\mathcal{S}_{9n}$; \textup{(ii)} for every $l\le j$,
$p_0\cdots p_{l-1}\in\mathcal{S}_{9n}$; \textup{(iii)} if $p$ is a head of length $j$, \m{\IsHead{j}{p}}, and
$a<j$, then $p_a$, $p_a^{-1}$ and all entries $s_{i,e^*_a}$ lie in $\mathcal{S}_{10n}$; \textup{(iv)} if
\m{\IsHead{j}{p}} and $s\in\Nb{j}{T}$, then \m{\HasMatSize{60n^2}{g}} and \m{\HasMatSize{99n^2}{\wl{j}{p}^{-1}g}}.
\end{lemma}
\begin{proof}
(i) The element $M=(sg)(\eta\wl{j}{p}^{-1}g)^{-1}=s\,\wl{j}{p}\,\eta^{-1}$ has size at most $1+(1+4)+1=7$. Take its
minor with the rows $e^*_0,\dots,e^*_{a-1},i$ and the columns $e_0,\dots,e_a$; by Lemma~\ref{l:det} its
determinant lies in $\mathcal{S}_{8(a+1)}\subseteq\mathcal{S}_{9n}$. Since $\eta^{-1}\in\Ql{j}{T}$ acts on
$\langle e_0,\dots,e_a\rangle$ by a unitriangular matrix, the determinant is that of the same minor of
$s\,\wl{j}{p}$; $\wl{j}{p}$ maps $e_b$ to $p_be^*_b$; and by (F3) the matrix of $s$ with the rows
$e^*_0,\dots,e^*_{a-1},i$ and the columns $e^*_0,\dots,e^*_a$ is lower triangular with diagonal
$1,\dots,1,s_{i,e^*_a}$ \Inp{I9}. So the determinant is $p_0\cdots p_a\cdot s_{i,e^*_a}$. (ii) Take $i=e^*_a$,
$a=l-1$, where $s_{e^*_a,e^*_a}=1$. (iii) $p_a=(p_0\cdots p_{a-1})^{-1}(p_0\cdots p_a)$,
$p_a^{-1}=(p_0\cdots p_a)^{-1}(p_0\cdots p_{a-1})$ and $s_{i,e^*_a}=(p_0\cdots p_a)^{-1}\cdot(p_0\cdots
p_a\,s_{i,e^*_a})$ lie in $\mathcal{S}_{1+9n}$. (iv) Apply Lemma~\ref{l:backsub} to $u=s$, $r=1$ and the vectors
$w=e^*_a$: $(1-s)e^*_a$ has size $10n+1$ by (iii), so $1-s^{-1}$ has size $10n+5+(20n+10)j$,
$g=s^{-1}(sg)$ has size at most $10n+8+(20n+10)(n-1)\le60n^2$, and $\wl{j}{p}^{-1}$, whose entries are $0$, $1$,
$p_a^{-1}$ or $\bar p_a$, has size $10n$.
\end{proof}
\leanline{\lean{minormuldualmemofSize}\sep\lean{minorsmemofSize}\sep\lean{headentriesmemofSize}\sep\lean{hasMatSizeofisHead}\sep\lean{cardsucclefloorOf}}

\sect{The tower measure}{s:tower}
\begin{lemma}[a floor measure and its quadric]\label{l:orbit}
Let $\cL$ be a level in $\Levels{\Omega}{T}{\mathfrak{m}}$ and $b$ a natural number with $b\le\varrho$.
\textup{(i)} The map $\sigma\mapsto\sigma^{-1}f_b$ is measurable and carries $\floorMeasure{\cL}{b}$ to the floor
quadric measure $\quadric{\cL}{b}$. \textup{(ii)} For $s\in\Hl{b}{T}$ the map $\sigma\mapsto s\sigma^{-1}f_b$ carries
$\floorMeasure{\cL}{b}$ to $\quadric{\cL}{b}$.
\end{lemma}
\begin{proof}
(i) For $b=\varrho$ this is (L1). Let $b=a<\varrho$ and $\sigma=\wo{a}{c}^{-1}\mathrm{n}_a(v,\beta)$. Then
$\sigma^{-1}e_a=\mathrm{n}_a(v,\beta)^{-1}(ce^*_a)=c\,e^*_a-cv+c(\beta'-\tau_Ah_T(v,v))e_a$ with $\beta'$ skew, by
(F6). In the decomposition $V_a=\langle e_a,e^*_a\rangle\oplus V_{a+1}$ this is the point of the flat chart of the
cone $h_T=0$ with the parameters $t=c$, $w=-cv$ and the skew parameter $\alpha=\beta'$: with these $t$ and $w$ the
chart has the $e_a$-coordinate $\bar t^{-1}\tau_A(-w^{*}T'w)+t\alpha=-c\,\tau_Ah_T(v,v)+c\,\alpha$. By (F6), $\beta'$ is
$-\beta$ plus a skew element that depends on $v$ only. So $(c,v,\beta)\mapsto(c,-cv,\beta')$ carries the measure
$\modA(c)^{N}\dd c\,\dd v\,\dd\nu_A(\beta)$, $N=k+1+2(\varrho-a-1)$ the rank of $V_{a+1}$, to
$\dd t\,\dd w\,\dd\nu_A(\alpha)$, the measure of the flat chart: the place axiom for $v\mapsto-cv$ removes the factor
$\modA(c)^{N}$, and the Haar measure $\nu_A$ is invariant under translations and under $\beta\mapsto-\beta$. So the
image of
$\floorMeasure{\cL}{a}$ is the quadric measure of $V_a$ for the standard pair $(e_a,e^*_a)$. It and
$\quadric{\cL}{a}$ are both the member at $\beta=0$ of a disintegration of $\vol_a$ along $h_T$ (\Inp{I4}(a),(b),(d)),
and a disintegration is unique (\Inp{I4}(c)). (ii) By \Inp{I4}(e) and (b), $s$ carries the disintegration of
$\vol_b$ to a disintegration of $\vol_b$, so $s$ preserves $\quadric{\cL}{b}$ by uniqueness; compose with (i).
\end{proof}
\leanline{\lean{measurePreservingorbitPoint}\sep\lean{measurePreservingmulVecorbitPoint}}

\begin{lemma}[floor measures are carried by their floors]\label{l:carried}
Let $\cL$ be a level in $\Levels{\Omega}{T}{\mathfrak{m}}$ and $l$ a natural number. Then $\floorMeasure{\cL}{l}$ is
carried by $\Hl{l}{T}$. Hence, if a measure $\mu$ equals $\floorMeasure{\cL}{l}$ and two functions $f,\Phi$ agree
on $\Hl{l}{T}$, then $\int f\,\dd\floorMeasure{\cL}{l}=\int\Phi\,\dd\mu$.
\end{lemma}
\begin{proof}
For $l<\varrho$: $\wo{l}{c}\in\Hl{l}{T}$, $\eta_l$ is carried by $\Nf{l}{T}\subseteq\Hl{l}{T}$ \Inp{I3}, and
\Inp{I2}. For $l=\varrho$, $l=n$ these are (L1), (L2); for $l>n$ the measure is $\delta_1$.
\end{proof}
\leanline{\lean{aefloorMeasuremem}\sep\lean{lintegralfloorMeasurecongr}}

\begin{lemma}[moving a Weyl law to the left]\label{l:desc}
Let $\cL$ be a level in $\Levels{\Omega}{T}{\mathfrak{m}}$, $l,i$ natural numbers with $l<\varrho$. Then
\[\floorMeasure{\cL}{l+i}*\dots*\floorMeasure{\cL}{l+1}*\floorMeasure{\cL}{l}=
\bigl(\varpi_l*(\floorMeasure{\cL}{l+i}*\dots*\floorMeasure{\cL}{l+1})\bigr)*\eta_l.\]
\end{lemma}
\begin{proof}
$\floorMeasure{\cL}{l}=\varpi_l*\eta_l$, and $\varpi_l$ commutes with every measure carried by $\Hl{l+1}{T}$, because
$\wo{l}{c}$ commutes with $\Hl{l+1}{T}$ (\Inp{I1}, Lemma~\ref{l:carried}, \Inp{I2}).
\end{proof}
\leanline{\lean{descConvfloorMeasuresucc}}

\begin{lemma}[translation by a Weyl ratio]\label{l:weylratio}
Let $\cL$ be a level in $\Levels{\Omega}{T}{\mathfrak{m}}$, $j$ a natural number with $j\le\varrho$,
$p,q\colon\NN\to A^{\times}$ and $\Phi\colon\GT{T}\to[0,\infty]$ measurable. Then
\[\int\Phi\bigl(\weylRatio{j}{p}{q}\,x\bigr)\,\dd\upperMeasure{\cL}{0}(x)=\prod_{a<j}\modA(p_a/q_a)^{k+2(\varrho-a)}\int\Phi\,\dd\upperMeasure{\cL}{0}.\]
\end{lemma}
\begin{proof}
Induction on $j$. Write $\upperMeasure{\cL}{0}=\upperMeasure{\cL}{j}*(\floorMeasure{\cL}{j-1}*\dots*\floorMeasure{\cL}{0})$
and use Lemma~\ref{l:desc} at $l=j$: $\upperMeasure{\cL}{0}=(\varpi_j*M')*\eta_j*M''$ with $M'$ carried by
$\Hl{j+1}{T}$. The Weyl ratio of length $j+1$ is that of length $j$ times
$\wo{j}{q_j}^{-1}\wo{j}{p_j}$, and $\wo{j}{q_j}^{-1}\wo{j}{p_j}\wo{j}{c}^{-1}=\wo{j}{cq_j/p_j}^{-1}$ up to an element
that fixes the measure; the change of variable $c\mapsto cp_j/q_j$ in
$\modA(c)^{k+1+2(\varrho-j-1)}\dd c$ produces the factor $\modA(p_j/q_j)^{k+1+2(\varrho-j-1)+1}$. The remaining
translation, by the ratio of length $j$, commutes with the factors carried by $\Hl{j}{T}$ and is handled by the
induction hypothesis.
\end{proof}
\leanline{\lean{lintegralweylRatiomul}}

\begin{lemma}[right translations on a floor]\label{l:right}
Let $\cL$ be a level in $\Levels{\Omega}{T}{\mathfrak{m}}$, $l$ a natural number, $S$ a subgroup of $\GT{T}$ that
contains $\Hl{l}{T}$, $f\colon\GT{T}\to[0,\infty]$ a measurable function invariant under left translations by
$\Ql{l+1}{T}\cap S$, and $s\in\Ql{l}{T}\cap S$. Then
$\int f(xs)\,\dd\floorMeasure{\cL}{l}(x)=\int f(x)\,\dd\floorMeasure{\cL}{l}(x)$.
\end{lemma}
\begin{proof}
Write $s=\gamma\nu$ with $\gamma\in\Hl{l}{T}$, $\nu\in\Nb{l}{T}$ (F1). For $x\in\Hl{l}{T}$, $x\gamma\nu(x\gamma)^{-1}$
lies in $\Nb{l}{T}\subseteq\Ql{l+1}{T}$ and in $S$, so $f(x\gamma\nu)=f(x\gamma)$; by Lemma~\ref{l:carried} we may
assume $s=\gamma\in\Hl{l}{T}$. If $l\le\varrho$, then on $\Hl{l}{T}$ we have $f(\sigma)=F(\sigma^{-1}f_l)$ with $F$
measurable \Inp{I5}, and by Lemma~\ref{l:orbit}(ii) and (i)
$\int F(\gamma^{-1}x^{-1}f_l)\,\dd\floorMeasure{\cL}{l}(x)=\int F\,\dd\quadric{\cL}{l}=\int F(x^{-1}f_l)\,\dd\floorMeasure{\cL}{l}(x)$.
If $l=n$ this is the right invariance in (L2); if $l>n$ then $\Hl{l}{T}=1$.
\end{proof}
\leanline{\lean{lintegralfloorMeasuremulright}}

\begin{proof}[Proof of Theorem A\textup{(i)}, \textup{(ii)}]
(i) Induction on $i$. For $i=0$, $l=\varrho+2$ and $s=1$. For $i+1$:
$\stackMeasure{\cL}_{l,i+1}=\stackMeasure{\cL}_{l+1,i}*\eta_l*\floorMeasure{\cL}{l}$. For measurable $f\ge0$ let
$F(x)=\int f(mx)\,\dd(\stackMeasure{\cL}_{l+1,i}*\eta_l)(m)$. An element of $\Ql{l+1}{T}\cap\Hl{l}{T}$ is $g\nu$ with
$g\in\Hl{l+1}{T}$, $\nu\in\Nf{l}{T}$ (F1); $\eta_l*\delta_\nu=\eta_l$,
$\eta_l*\delta_g=\delta_g*\eta_l$ \Inp{I3}, and $\stackMeasure{\cL}_{l+1,i}*\delta_g=\stackMeasure{\cL}_{l+1,i}$ by
induction; so $F$ is invariant under $\Ql{l+1}{T}\cap\Hl{l}{T}$. By Lemma~\ref{l:right} with $S=\Hl{l}{T}$,
$\int F(xs)\,\dd\floorMeasure{\cL}{l}(x)=\int F\,\dd\floorMeasure{\cL}{l}$, which is the claim. For $l=0$,
$i=\varrho+2$ every $s\in\GT{T}=\Hl{0}{T}$ is allowed.

(ii) Induction on $j$: $\stackMeasure{\cL}_{0,j+1}=\eta_j*\floorMeasure{\cL}{j}*\etab{j}*(\floorMeasure{\cL}{j-1}*\dots)$;
$\floorMeasure{\cL}{j}$ is carried by $\Hl{j}{T}$, whose unit masses commute with $\etab{j}$, so
$\floorMeasure{\cL}{j}*\etab{j}=\etab{j}*\floorMeasure{\cL}{j}$ \Inp{I1}; and $\eta_j*\etab{j}=\etab{j+1}$ \Inp{I3}.
\end{proof}

\begin{lemma}[the stack measures are not zero]\label{l:nz}
Let $\cL$ be a level in $\Levels{\Omega}{T}{\mathfrak{m}}$ and $l,i$ natural numbers. Then
$\stackMeasure{\cL}_{l,i}\ne0$.
\end{lemma}
\begin{proof}
The total mass of a convolution is the product of the total masses. $\eta_l$ and $\delta_1$ are not zero,
$\floorMeasure{\cL}{n}\ne0$ by definition of a level, and for $l\le\varrho$ the image $\quadric{\cL}{l}$ of
$\floorMeasure{\cL}{l}$ (Lemma~\ref{l:orbit}) is not zero.
\end{proof}
\leanline{\lean{stackMeasurenezero}}

\begin{lemma}[left translations of a factor of a Haar measure]\label{l:factor}
Let $\cL$ be a level in $\Levels{\Omega}{T}{\mathfrak{m}}$ whose tower measure $\groupMeasure{\cL}$ is a Haar
measure. Let $\rho,\nu$ be $s$-finite measures on $\GT{T}$ with $\groupMeasure{\cL}=\rho*\nu$, let $S$ be a subgroup
and $s\in\GT{T}$, and suppose that $\rho$ is carried by $S$, that $\rho*\delta_{s'}=\rho$ for all $s'\in S$, and that
$\delta_s*\rho=\rho*\delta_s$. Then for every measurable $\Phi\colon\GT{T}\to[0,\infty]$ that is invariant under
left translations by $S$,
\[\int\Phi(sg)\,\dd\nu(g)=\int\Phi\,\dd\nu.\]
\end{lemma}
\begin{proof}
Choose a measurable $f>0$ with $\int f\,\dd\groupMeasure{\cL}<\infty$ and put $D(x)=\int f(yx)\,\dd\rho(y)$. We claim
that for \emph{every} $s$-finite $\nu'$ with $\groupMeasure{\cL}=\rho*\nu'$
\[\int\Phi(x)D(x)^{-1}f(x)\,\dd\groupMeasure{\cL}(x)=\int\Phi\,\dd\nu'.\]
Indeed, by \Inp{I1} the left side is
\[\iint\Phi(yx)D(yx)^{-1}f(yx)\,\dd\rho(y)\,\dd\nu'(x);\]
for $y\in S$, $\Phi(yx)=\Phi(x)$ and
$D(yx)=D(x)$ by the right invariance of $\rho$; $D(x)<\infty$ for $\nu'$-almost every $x$ and $D(x)>0$ because
$\rho\ne0$ (Lemma~\ref{l:nz} excludes $\rho=0$, as $\groupMeasure{\cL}\ne0$); so the inner integral is $\Phi(x)$.
Apply the claim to $\nu$ and to $\delta_s*\nu$: the latter is allowed since
$\rho*(\delta_s*\nu)=\delta_s*\rho*\nu=\delta_s*\groupMeasure{\cL}=\groupMeasure{\cL}$ by left invariance.
\end{proof}
\leanline{\lean{lintegralmullefteqoffactor}}

\begin{proof}[Proof of Theorem A\textup{(iii)}]
$\groupMeasure{\cL}=\stackMeasure{\cL}_{j,\varrho+2-j}*\stackMeasure{\cL}_{0,j}$. The first factor is carried by
$\Hl{j}{T}$, is invariant under right translations by $\Hl{j}{T}$ by (i), and commutes with the unit mass at
$\wl{j}{p}^{-1}$ by (F5) and \Inp{I1}. Apply Lemma~\ref{l:factor} with $S=\Hl{j}{T}$ and $s=\wl{j}{p}^{-1}$.
\end{proof}

\sect{Nearness, floor values and weights}{s:near}
\begin{lemma}[measurability and openness]\label{l:meas}
Let $\cL$ be a level in $\Levels{\Omega}{T}{\mathfrak{m}}$, $j$ a natural number and $a$ a natural number with
$a<\varrho$. \textup{(i)} The set of the points near at $j$ is measurable. \textup{(ii)} The set of the triples
$(x,y,c)\in\GT{T}\times\GT{T}\times A$ with $c\in\floorValues{\cL}{a}{x}{y}$ is open. \textup{(iii)} For all
$x,y\in\GT{T}$: $0<\vol(\floorValues{\cL}{a}{x}{y})<\infty$.
\end{lemma}
\begin{proof}
(i) It is the product $\Ql{j}{T}^{-1}\BoxL{\cL}$ of a closed set \Inp{I2} and a compact set, hence closed. (ii) If
$c=w_{e^*_a}$ for an admissible $w$, then for $(x',y',c')$ near $(x,y,c)$ the vector $w$ with the coordinate at
$e^*_a$ replaced by $c'$ is admissible, since $\Ball{\cL}$ is open and the conditions are continuous. (iii) The
set is open by (ii) and contains $0$ (take $w=0$; $0\in\Ball{\cL}$), so its volume is positive; it is contained in
the image of the compact set $x\,\overline{\Ball{\cL}}$ under a coordinate, so its volume is finite.
\end{proof}
\leanline{\lean{measurableSetisNear}\sep\lean{isOpenfloorValues}\sep\lean{volumefloorValuesnezeronetop}}

\begin{lemma}[invariance of floor values and weights]\label{l:valinv}
Let $\cL$ be a level in $\Levels{\Omega}{T}{\mathfrak{m}}$, $x,y,g\in\GT{T}$. \textup{(i)} For $a<\varrho$ and
$\xi\in\Ql{a+1}{T}$: $\floorValues{\cL}{a}{\xi x}{y}=\floorValues{\cL}{a}{x}{y}$. \textup{(ii)} For a natural number
$j$ and $\xi\in\Ql{j}{T}$: $\floorWeight{\cL}{j}{\xi x}{y}=\floorWeight{\cL}{j}{x}{y}$ and
$\invWeight{\cL}{j}{x}(\xi g)=\invWeight{\cL}{j}{x}(g)$.
\end{lemma}
\begin{proof}
(i) If $w$ is admissible for $(x,y)$ then $\xi w$ is admissible for $(\xi x,y)$, because
$(\xi x)^{-1}\xi w=x^{-1}w$ and $(\xi xy)^{-1}\xi w=(xy)^{-1}w$, and by (F3) $\xi w$ has the same coordinates at
$e^*_b$, $b\le a$, as $w$; conversely with $\xi^{-1}$. (ii) The weight is a product over $a<\min(j,\varrho)$, and
$\Ql{j}{T}\subseteq\Ql{a+1}{T}$; nearness at $j$ is $\Ql{j}{T}$-invariant.
\end{proof}
\leanline{\lean{floorValuesmul}\sep\lean{invWeightmul}}
Also $\floorWeight{\cL}{a+1}{x}{y}=\floorWeight{\cL}{a}{x}{y}\cdot\vol(\floorValues{\cL}{a}{x}{y})^{\volumeExponent{\cL}}$
for $a<\varrho$, the weights are neither $0$ nor $\infty$ (Lemma~\ref{l:meas}(iii)), and they are measurable in
$(x,y)$ because the set of Lemma~\ref{l:meas}(ii) is open.

\begin{lemma}[the upper measure of the near points]\label{l:upper}
Let $\cL$ be a level in $\Levels{\Omega}{T}{\mathfrak{m}}$ and $V\in[0,\infty]$ with \m{\FloorMassLe{\cL}{a}{V}{1}}
for every $a\le\varrho$. Then for every natural number $l$ and every $c\in\GT{T}$,
\[\upperMeasure{\cL}{l}\{g:\ gc\text{ is near at }\varrho+2\}\le V^{\,n-l}.\]
\end{lemma}
\begin{proof}
Induction on $\varrho+2-l$. For $l\ge n$ the bound is $1$: for $l=n$ it is the last part of (L2), since
$\Ql{\varrho+2}{T}=\UT{T}$. For $l<n$, $\upperMeasure{\cL}{l}=\upperMeasure{\cL}{l+1}*\floorMeasure{\cL}{l}$, so the
measure is $\int\upperMeasure{\cL}{l+1}\{g:g\sigma c\text{ near at }\varrho+2\}\,\dd\floorMeasure{\cL}{l}(\sigma)$
\Inp{I1}. The integrand is at most $V^{n-l-1}$ by induction, and it is $0$ unless $\sigma c$ is near at $l+1$,
because $\upperMeasure{\cL}{l+1}$ is carried by $\Hl{l+1}{T}$. The $\floorMeasure{\cL}{l}$-measure of these $\sigma$
is at most $V$.
\end{proof}
\leanline{\lean{upperMeasureisNearle}}

\sect{Kernels}{s:kernels}
\begin{lemma}[without the kernel axioms]\label{l:nokernel}
Let $\cL$ be a level in $\Levels{\Omega}{T}{\mathfrak{m}}$, $k\colon\GT{T}\times\GT{T}\to\CC$, $l$ a natural number
and $g,y,x\in\GT{T}$. \textup{(i)} If $\pivotModulus{\cL}{l}{k}{g}{y}\ne0$, then $k(\sigma g,\sigma y)\ne0$ for some
$\sigma\in\Hl{l}{T}$. \textup{(ii)} If $l\le\varrho$, then
\[\int\bigl|\upperPairing{\cL}{l+1}{k}{\sigma g}{\sigma gx}\bigr|^{2}\,\dd\floorMeasure{\cL}{l}(\sigma)\le
\frac{\weightedSquareSum{\cL}{k}{x}{l}{g}}{\floorWeight{\cL}{l}{g}{x}}.\]
\end{lemma}
\begin{proof}
(i) The measure $\floorMeasure{\cL}{l}$ is carried by $\Hl{l}{T}$ and $\upperMeasure{\cL}{l+1}$ by $\Hl{l+1}{T}$.
So there is $\sigma_1\in\Hl{l}{T}$ with
\[\upperPairing{\cL}{l+1}{k}{\sigma_1g}{\sigma_1y}\ne0,\]
and then $k(\sigma_2\sigma_1g,\sigma_2\sigma_1y)\ne0$ for some $\sigma_2\in\Hl{l+1}{T}$. (ii) By definition the weighted square sum is at least
$\int\weightedSquare{\cL}{l}{k}{x}{\sigma g}\,\dd\floorMeasure{\cL}{l}(\sigma)$, and
$\floorWeight{\cL}{l}{\sigma g}{x}=\floorWeight{\cL}{l}{g}{x}$ for $\sigma\in\Hl{l}{T}$
(Lemma~\ref{l:valinv}).
\end{proof}
\leanline{\lean{existsnezeroofpivotModulus}\sep\lean{lintegralupperPairingsqle}}

\begin{setting}[a pairing kernel]\label{set:pk}
Let $\cL$ be a level in $\Levels{\Omega}{T}{\mathfrak{m}}$, $j$ a natural number,
$k\colon\GT{T}\times\GT{T}\to\CC$ and $d\in\GT{T}$ with \m{\IsPairingKernel{\cL}{j}{k}{d}}.
\end{setting}

\begin{lemma}[averages floor by floor]\label{l:succ}
In Setting~\ref{set:pk}, let $l$ be a natural number and $g,y\in\GT{T}$. Then
\textup{(i)} $x\mapsto k(xg,xy)$ is integrable for $\upperMeasure{\cL}{l}$;
\textup{(ii)} $\sigma\mapsto\upperPairing{\cL}{l+1}{k}{\sigma g}{\sigma y}$ is integrable for $\floorMeasure{\cL}{l}$
and
\[\upperPairing{\cL}{l}{k}{g}{y}=\int\upperPairing{\cL}{l+1}{k}{\sigma g}{\sigma y}\,\dd\floorMeasure{\cL}{l}(\sigma);\]
hence $|\upperPairing{\cL}{l}{k}{g}{y}|\le\pivotModulus{\cL}{l}{k}{g}{y}$, and $\upperPairing{\cL}{l}{k}{\cdot}{\cdot}$
is measurable;
\textup{(iii)} if $j\le l$, then \m{\IsPairingKernel{\cL}{l}{k}{d}}.
\end{lemma}
\begin{proof}
(i) is (K1) and (K4). (ii) is Fubini's theorem for
$\upperMeasure{\cL}{l}=\upperMeasure{\cL}{l+1}*\floorMeasure{\cL}{l}$, with (i). (iii) Only (K3) depends on $j$, and
$\Hl{l}{T}\subseteq\Hl{j}{T}$.
\end{proof}
\leanline{\lean{integrableupperMeasure}\sep\lean{upperPairingeqintegralsucc}\sep\lean{enormupperPairingle}\sep\lean{stronglyMeasurableupperPairing}\sep\lean{mono}}

\begin{lemma}[invariance of the averages]\label{l:kinv}
In Setting~\ref{set:pk}, let $l$ be a natural number with $j\le l$, let $s\in\Ql{l}{T}\cap\Hl{j}{T}$ and
$g,y\in\GT{T}$. Then $\upperPairing{\cL}{l}{k}{sg}{sy}=\upperPairing{\cL}{l}{k}{g}{y}$.
\end{lemma}
\begin{proof}
Induction on $\varrho+2-l$. If $l\ge\varrho+2$, then $s\in\UT{T}=\Nb{\varrho+1}{T}$, the average is $k$ itself, and
(K2) with $\nu=\nu'=s$ gives the factor $\psi_A(\gL{\cL}^{-1}(\sd(dsd^{-1})-\sd(s)))=1$ by (K3). Otherwise write
both sides by Lemma~\ref{l:succ}(ii). The function
$\sigma\mapsto\upperPairing{\cL}{l+1}{k}{\sigma g}{\sigma y}$ is invariant under $\Ql{l+1}{T}\cap\Hl{j}{T}$ by
induction, so Lemma~\ref{l:right} with $S=\Hl{j}{T}\supseteq\Hl{l}{T}$ allows to replace $\sigma$ by $\sigma s$.
\end{proof}
\leanline{\lean{upperPairinginvariant}}

\begin{lemma}[invariance and support of the weighted squares]\label{l:sqinv}
In Setting~\ref{set:pk}, let $x\in\GT{T}$. \textup{(i)} The function $\weightedSquare{\cL}{j}{k}{x}{\cdot}$ is
invariant under $\Ql{j+1}{T}$ and vanishes at every $g$ that is not near at $j+1$. \textup{(ii)} The function
$\weightedSquareSum{\cL}{k}{x}{j}{\cdot}$ is invariant under $\Ql{j}{T}$ and vanishes at every $g$ that is not near
at $j$; it is measurable in $(x,g)$.
\end{lemma}
\begin{proof}
(i) Let $s\in\Ql{j+1}{T}$, and write $s=\gamma\nu$ with $\gamma\in\Hl{j}{T}$, $\nu\in\Nb{j}{T}$ (F1); then
$\gamma=s\nu^{-1}\in\Ql{j+1}{T}\cap\Hl{j}{T}$. The weight is invariant (Lemma~\ref{l:valinv}). For the average at
$j+1$: Lemma~\ref{l:kinv} removes $\gamma$, and (K2) at $l=j$ with $\nu'=\nu$ shows that
$\upperPairing{\cL}{j+1}{k}{\nu g}{\nu gx}$ and $\upperPairing{\cL}{j+1}{k}{g}{gx}$ differ by a factor of absolute
value $1$. The support is
(K5). (ii) Downward induction on $j$ along the recursion: by (i), (iii) of Lemma~\ref{l:succ} and the induction
hypothesis the integrand $\sigma\mapsto(\dots)(\sigma g)$ is $\Ql{j+1}{T}$-invariant in $\sigma g$, and
Lemma~\ref{l:right} with $S=\GT{T}$ gives the invariance under $s\in\Ql{j}{T}$. If the sum at $g$ is not $0$, some
$\sigma\in\Hl{j}{T}$ has $\sigma g$ near at $j+1$, so $g$ is near at $j$.
\end{proof}
\leanline{\lean{weightedSquareinvariantsupport}\sep\lean{weightedSquareSuminvariantsupport}\sep\lean{measurableweightedSquareSum}}

\begin{lemma}[Cauchy--Schwarz on a floor]\label{l:cs}
In Setting~\ref{set:pk}, let $l$ be a natural number, $g,y\in\GT{T}$, $K\in[0,\infty]$ and
$u\colon\GT{T}\to[0,\infty]$, $v\colon\GT{T}\to[0,\infty]$ measurable with \m{\FloorMassLe{\cL}{l}{K}{u}} and
\[\bigl|\upperPairing{\cL}{l+1}{k}{\sigma g}{\sigma y}\bigr|^{2}\le u(\sigma g)\,v(\sigma)\qquad\text{for all }\sigma\in\Hl{l}{T}.\]
Then $\pivotModulus{\cL}{l}{k}{g}{y}^{2}\le K\int v\,\dd\floorMeasure{\cL}{l}$.
\end{lemma}
\begin{proof}
By (K5) the integrand of the pivot vanishes unless $\sigma g$ is near at $l+1$. Restrict the integral to these
$\sigma$ and to $\Hl{l}{T}$, apply \Inp{I8} with $f_1=u(\sigma g)$ and $f_2=v$, and bound the integral of $f_1$ by
$K$.
\end{proof}
\leanline{\lean{pivotModulussqle}}

\begin{lemma}[the self-norm, floor by floor]\label{l:sn}
In Setting~\ref{set:pk}, let $x\in\GT{T}$ and $l$ a natural number. \textup{(i)} If $l\le\varrho$, then
\[\selfNorm{\cL}{l}{k}{x}=\int\weightedSquare{\cL}{l}{k}{x}{g}\,\dd(\floorMeasure{\cL}{l}*\dots*\floorMeasure{\cL}{0})(g)+\selfNorm{\cL}{l+1}{k}{x}.\]
\textup{(ii)} If $l\le n$, then $\selfNorm{\cL}{l}{k}{x}\le\selfNorm{\cL}{0}{k}{x}$.
\end{lemma}
\begin{proof}
(i) Insert the recursion for the weighted square sum and use \Inp{I1} for
$\floorMeasure{\cL}{l}*(\floorMeasure{\cL}{l-1}*\dots*\floorMeasure{\cL}{0})$. (ii) follows from (i) by induction.
\end{proof}
\leanline{\lean{selfNormeqaddsucc}\sep\lean{selfNormleselfNormzero}}
\sect{Sizes in a gauge and fibre bounds}{s:gsizes}
\begin{lemma}[the vector of a Heisenberg element]\label{l:heis}
Let $A$ be a place, $\mathcal{S}$ a size in $\Sizes{T}$, $a$ a natural number with $a<\varrho$, $u\in A^{\times}$,
$i$ a natural number, $v\in V_{a+1}$ and $\beta$ a skew element of $A$. If there is $\nu\in\Nb{a}{T}$ such that the
column at $e^*_a$ of $u^{-1}(\nu\,\mathrm{n}_a(v,\beta)-1)$ lies in $\vecsOfSize{i}$, then $u^{-1}v\in\vecsOfSize{i}$
and $u^{-1}(\beta-\tau_Ah_T(v,v))\in\mathcal{S}_i$.
\end{lemma}
\begin{proof}
The column at $e^*_a$ of $\mathrm{n}_a(v,\beta)$ is $e^*_a+v+(\beta-\tau_Ah_T(v,v))e_a$; its coordinates at
$e^*_b$, $b<a$, vanish, so by (F4) multiplication by $\nu$ does not change its coordinates of floor at least $a$:
those on $V_{a+1}$, which are $v$, and the one at $e_a$.
\end{proof}
\leanline{\lean{memvecsOfSizeofheisenbergElem}}

\begin{lemma}[coordinates in a floor basis]\label{l:basis}
In Setting~\ref{set:gauge}, let $b\le\varrho$, $j$ a natural number and $x$ a coordinate vector for the floor basis
$E_b$ with $E_bx\in\vecsOfSize{j}$. Then $x\in\vecsOfSize{j+4}$.
\end{lemma}
\begin{proof}
By \Inp{I4}(d), $E_b^{*}\JT{T}E_bx$ is $x$ with its coordinates permuted. The entries of $E_b$ are those of
$E(\cL)$ or $0$, $1$, so $E_b^{*}$ has size $1$ by (G2), $E_b^{*}\JT{T}$ has size $3$, and
$E_b^{*}\JT{T}(E_bx)\in\vecsOfSize{3+j+1}$.
\end{proof}
\leanline{\lean{memvecsOfSizeoffloorBasismulVec}}

\begin{lemma}[integrating against a fibre bound]\label{l:fibint}
In Setting~\ref{set:gauge}, let $b\le\varrho$, $q$ a measure on $V$, $j$ a natural number and $M,K\in[0,\infty]$
with \m{\HasFibreBound{\Gamma}{b}{q}{j}{M}}. \textup{(i)} Let $\Phi\colon V\to[0,\infty]$ and let $G$ be a function
of the fixed and the orthonormal coordinates with values in $[0,\infty]$ such that, for every $P$, $G(P,\cdot)$ is
measurable with $\int G(P,w)\,\dd w\le K$, and such that $\Phi(E_bx)\le G(P(x),w(x))$ for every coordinate vector
$x\in\vecsOfSize{j+4}$ with fixed coordinates $P(x)$ and orthonormal coordinates $w(x)$. Then
$\int_{\vecsOfSize{j}}\Phi\,\dd q\le M\cdot K$. \textup{(ii)} Let $\varphi\colon A^{\Omega}\to[0,\infty]$ be
measurable with $\int_{\vecsOfSize{j+4}}\varphi(w)\,\dd w\le K$. Then
\[\int_{\vecsOfSize{j}}\varphi\bigl(\orthonormalPart{\cL}^{*}\JT{T}x\bigr)\,\dd q(x)\le M\cdot K.\]
\end{lemma}
\begin{proof}
(i) With the notation of Definition~\ref{d:fib}, $\int_{\vecsOfSize{j}}\Phi\,\dd q$ is the integral of
$\mathbf{1}[E_b\mathsf{k}\in\vecsOfSize{j}]\,\Phi(E_b\mathsf{k})$ over $(P,w,z)$. By Lemma~\ref{l:basis} the
integrand is at most $\mathbf{1}[\mathsf{k}(P,w,z)\in\vecsOfSize{j+4}]\,G(P,w)$; integrate over $z$ (at most
$\mu(P)G(P,w)$), then over $w$ (at most $\mu(P)K$), then over $P$. (ii) By \Inp{I4}(d),
$\orthonormalPart{\cL}^{*}\JT{T}E_bx=w(x)$; apply (i) with $G(P,w)=\mathbf{1}[w\in\vecsOfSize{j+4}]\varphi(w)$.
\end{proof}
\leanline{\lean{setLIntegralleofhasFibreBound}\sep\lean{setLIntegralorthonormalPartle}}

\begin{lemma}[fibre bounds for the quadric and for the volume]\label{l:fibq}
In Setting~\ref{set:gauge}, let $b\le\varrho$ and $j$ a natural number. \textup{(i)} Let $M\in[0,\infty]$ be such
that the quadric measure with $\beta=0$ of the sum of $|\pi_b|+1$ hyperbolic planes gives $\vecsOfSize{j+5}$ mass at
most $M$. Then \m{\HasFibreBound{\Gamma}{b}{\quadric{\cL}{b}}{j}{M}}. \textup{(ii)} In particular
\m{\HasFibreBound{\Gamma}{b}{\quadric{\cL}{b}}{j}{\Lambda(\Gamma)^{99n(j+6)}}}. \textup{(iii)} Also
\m{\HasFibreBound{\Gamma}{b}{\vol_b}{j}{\Lambda(\Gamma)^{99n(j+6)}}}.
\end{lemma}
\begin{proof}
(i) In the coordinates of $E_b$, $\quadric{\cL}{b}$ is the image of $\dd t\,\dd(\pi',w,\pi'')\,\dd\nu_A(\alpha)$
under the flat chart, which fills in the coordinate $t'$. Take $Z$ the space of the pairs $(\pi'',\alpha)$ with
its product measure and $\mathsf{k}$ the flat chart after regrouping the variables as $((t,\pi'),w,(\pi'',\alpha))$,
which preserves the product measure. Fix $P=(t,\pi')$ and $w$, and a point $z_0$ of the fibre set
$\{z:\mathsf{k}(P,w,z)\in\vecsOfSize{j+4}\}$. For another point $z$ of it, the difference of the two values of
$\mathsf{k}$ has $w$-coordinate $0$, sizes at most $j+5$ by rule (b), and is the value at $(P,0,z-z_0)$ of the flat
chart of the cone of the hyperbolic planes alone (the chart is affine in $(\pi'',\alpha)$ for fixed $t,\pi',w$). So
the fibre set, translated by $z_0$, lies in the corresponding set for the cone and size $j+5$, whose measure we
call $\mu(P)$; and $\int\mu(P)\,\dd P$ is the cone measure of $\vecsOfSize{j+5}$. (ii) is (i) with (G5), since
$\mathfrak{m}\le|\pi_b|\le\mathfrak{m}+\varrho$. (iii) Take $Z$ the space of the coordinates $(t',\pi'')$ with its
volume, $\mathsf{k}$ the regrouping, and $\mu(P)=\vol(\Gamma_{j+4})^{1+|\pi_b|}$ for $P\in\vecsOfSize{j+4}$, $0$
otherwise; then $\int\mu\le\Lambda(\Gamma)^{2(j+5)(1+\mathfrak{m}+\varrho)}$ by (G4), and
$2(1+\mathfrak{m}+\varrho)\le20n$ by (G1).
\end{proof}
\leanline{\lean{hasFibreBoundfloorQuadricofle}\sep\lean{hasFibreBoundfloorQuadric}\sep\lean{hasFibreBoundfloorVolume}}

\begin{lemma}[an integral over a Heisenberg group]\label{l:heisint}
In Setting~\ref{set:gauge}, let $a<\varrho$, $u\in A^{\times}$, $j$ a natural number, $C,K,M\in[0,\infty]$,
$\Phi\colon\GT{T}\to[0,\infty]$ and $\Psi\colon V\to[0,\infty]$ measurable. Suppose
$\int_{\vecsOfSize{j}}\Psi\,\dd\vol_{a+1}\le M$ and $\floorModulus{\cL}{a}{u}\cdot\Lambda(\Gamma)^{j+2}\cdot M\le K$,
and that for every $v\in V_{a+1}$ and skew $\beta$ with $\Phi(\mathrm{n})\ne0$, $\mathrm{n}=\mathrm{n}_a(v,\beta)$:
there is $\nu\in\Nb{a}{T}$ such that the column at $e^*_a$ of $u^{-1}(\nu\mathrm{n}^{-1}-1)$ lies in
$\vecsOfSize{j}$, and $\Phi(\mathrm{n})\le C\cdot\Psi(u^{-1}v)$. Then $\int\Phi\,\dd\eta_a\le C\cdot K$.
\end{lemma}
\begin{proof}
$\eta_a$ is the image of $\dd v\,\dd\nu_A(\beta)$. Let $\Phi(\mathrm{n}_a(v,\beta))\ne0$. As
$\mathrm{n}^{-1}=\mathrm{n}_a(-v,\beta')$ with $\beta'$ a translate of $-\beta$ depending on $v$ only (F6),
Lemma~\ref{l:heis} gives $u^{-1}v\in\vecsOfSize{j}$ and $u^{-1}(\beta'-\tau_Ah_T(v,v))\in\Gamma_j$. So for fixed
$v$ the set $S_v$ of the $\beta$ with $\Phi\ne0$ has $u^{-1}(\beta_1-\beta_2)\in\Gamma_{j+1}$ for
$\beta_1,\beta_2\in S_v$; by (G2) and the rules, $u^{-1}(\tau_A\delta_A^{-1}\beta+\beta'')\in\Gamma_{j+3}$ for
differences $\beta,\beta''$, and by \Inp{I7} and (G4)
$\nu_A(S_v)\le(\Lambda(\Gamma)^{j+4}\modA(u))^{1/2}\le\Lambda(\Gamma)^{j+2}\modA(u)^{1/2}$. Hence
\[\begin{aligned}\int\Phi\,\dd\eta_a&\le C\,\Lambda(\Gamma)^{j+2}\modA(u)^{1/2}\int\mathbf{1}[u^{-1}v\in\vecsOfSize{j}]\,\Psi(u^{-1}v)\,\dd v\\
&=C\,\Lambda(\Gamma)^{j+2}\modA(u)^{1/2}\modA(u)^{N}\int_{\vecsOfSize{j}}\Psi\,\dd\vol_{a+1},\end{aligned}\]
where $N=k+1+2(\varrho-(a+1))$ is the rank of $V_{a+1}$. The factor in front is
$C\,\floorModulus{\cL}{a}{u}\,\Lambda(\Gamma)^{j+2}$.
\end{proof}
\leanline{\lean{lintegralheisenbergMeasurele}}

\begin{lemma}[three size bounds]\label{l:three}
In Setting~\ref{set:gauge}: \textup{(i)} if $g\in\BoxL{\cL}$, $x,y\in\positionsOfSize{\cL}{\Gamma}$ and
$\sigma=gx(gy)^{-1}$, then \m{\HasMatSize{14}{\gL{\cL}^{-1}(\sigma-1)}}; \textup{(ii)} if $b\le\varrho$,
$\sigma\in\Hl{b}{T}$, $g\in\GT{T}$ with \m{\HasMatSize{i}{g}} for a natural number $i$, and $\sigma g$ is near at
$b+1$, \m{\IsNear{\cL}{b+1}{\sigma g}}, then $\sigma^{-1}f_b\in\vecsOfSize{i+6}$; \textup{(iii)} if $j\le\varrho$,
$p\colon\NN\to A^{\times}$ is a head of length $j$, \m{\IsHead{j}{p}}, and $g$ and $\wl{j}{p}^{-1}g$ are near at
$j$, \m{\IsNear{\cL}{j}{g}} and \m{\IsNear{\cL}{j}{\wl{j}{p}^{-1}g}}, then there is $\sigma\in\Hl{j}{T}$ with
\m{\HasMatSize{60n^2}{\sigma g}} and \m{\HasMatSize{99n^2}{\sigma\wl{j}{p}^{-1}g}}.
\end{lemma}
\begin{proof}
(i) $\sigma-1=g(x-y)(gy)^{-1}$. Here $\gL{\cL}^{-1}(x-y)$ has size $2$, $y=1+\gL{\cL}\cdot\gL{\cL}^{-1}(y-1)$ has
size $3$, $gy$ size $5$, its inverse size $9$ (Lemma~\ref{l:adj}), and the product size $(1+2+1)+9+1$. (ii) Let
$\xi\in\Ql{b+1}{T}$ with $\xi\sigma g\in\BoxL{\cL}$. By (F2) and (F3) the coordinates of $\sigma^{-1}f_b$ of floor
less than $b$ vanish and the others are those of $\sigma^{-1}\xi^{-1}f_b$, a column of
$g(\xi\sigma g)^{-1}$, of size $i+5+1$. (iii) Let $\xi g\in\BoxL{\cL}$, $\xi=\sigma\nu$ with $\sigma\in\Hl{j}{T}$,
$\nu\in\Nb{j}{T}$ (F1). Then $(\sigma\nu\sigma^{-1})(\sigma g)\in\BoxL{\cL}$ with
$\sigma\nu\sigma^{-1}\in\Nb{j}{T}$, and $\wl{j}{p}^{-1}\sigma g=\sigma\wl{j}{p}^{-1}g$ is near at $j$; apply
Lemma~\ref{l:minor}(iv) to $\mathsf{K}=\BoxL{\cL}$, $\sigma g$ and $s=\sigma\nu\sigma^{-1}$.
\end{proof}
\leanline{\lean{conjmulinvclosetoone}\sep\lean{orbitPointmemvecsOfSize}\sep\lean{existshasMatSizeofisNear}}

\begin{lemma}[a position between two near points is small]\label{l:subone}
In Setting~\ref{set:gauge}, let $\Gamma_S$ be a second gauge in $\Gauges{\cL}$, $j$ a natural number with
$j\le\varrho$, and $p,q\colon\NN\to A^{\times}$ heads of length $j$ for $\Gamma_S$: \m{\IsHead{j}{p}} and
\m{\IsHead{j}{q}} for $\Gamma_S$. Suppose that, if $j\ne0$, then $\gL{\cL}^{-1}c\in\Gamma_1$ for all
$c\in(\Gamma_S)_{10n}$. Let $x,y\in\GT{T}$ be such that: there are $u,v\in\UT{T}$ with $ux\in\BoxL{\cL}$,
$v(xy)\in\BoxL{\cL}$ and \m{\HasMatSize{2}{\gL{\cL}^{-1}(v(xy)-ux)}}; $\wl{j}{p}^{-1}x$ is near at $j$,
\m{\IsNear{\cL}{j}{\wl{j}{p}^{-1}x}}; $\wl{j}{q}^{-1}xy$ is near at $j$; and for every $a$ with $j\le a<\varrho$
there is a unit $c\in\floorValues{\cL}{a}{x}{y}$ with $c^{-1}\in\Gamma_1$. Then
\m{\HasMatSize{99n}{\gL{\cL}^{-1}(y-1)}}.
\end{lemma}
\begin{proof}
Put $X=ux$, $Y=v(xy)$; then $y-1=X^{-1}\bigl((Y-X)-(1-uv^{-1})Y\bigr)$. We bound $\gL{\cL}^{-1}(1-uv^{-1})$ by
Lemma~\ref{l:backsub} for the element $vu^{-1}\in\UT{T}=\Nb{\varrho}{T}$, $r=\gL{\cL}$ and $c=6$. For $a<j$ take
$w=ue^*_a$, the column of $u$ at $e^*_a$: by Lemma~\ref{l:minor}(iii) for $\Gamma_S$ and $\mathsf{K}=\BoxL{\cL}$ its
entries, and those of $ve^*_a$, lie in $(\Gamma_S)_{10n}$, so $\gL{\cL}^{-1}w$ and $\gL{\cL}^{-1}(u-v)e^*_a$ have
size at most $2$. For $a\ge j$ take the unit $c=w'_{e^*_a}$ of the hypothesis, $w'$ admissible, and $w=c^{-1}uw'$:
by (F3) $w_{e^*_a}=1$; $uw'=X(x^{-1}w')$ has size $3$ by (G3); and
$\gL{\cL}^{-1}(u-v)w'=\gL{\cL}^{-1}(X-Y)u_1+Y\gL{\cL}^{-1}(u_1-u_2)$ with $(u_1,u_2)$ the close pair of $w'$, of size
$5$. So $\gL{\cL}^{-1}(1-uv^{-1})$ has size $10+20\varrho$, and $\gL{\cL}^{-1}(y-1)$ has size at most
$5+(13+20\varrho)+1\le99n$.
\end{proof}
\leanline{\lean{hasMatSizesuboneoffloorValues}}

\sect{Masses of floors and integrals of weights}{s:masses}
\begin{lemma}[the mass of a floor through its quadric]\label{l:chart}
In Setting~\ref{set:gauge}, let $b\le\varrho$, $K\in[0,\infty]$, and $f\colon\GT{T}\to[0,\infty]$ measurable and
invariant under left translations by $\Ql{b+1}{T}$. Suppose that for every $g\in\BoxL{\cL}$ there is a measurable
$\Psi\colon V\to[0,\infty]$ with $\int_{\vecsOfSize{7}}\Psi\,\dd\quadric{\cL}{b}\le K$ and
$f(\sigma g)\le\Psi(\sigma^{-1}f_b)$ for all $\sigma\in\Hl{b}{T}$. Then \m{\FloorMassLe{\cL}{b}{K}{f}}.
\end{lemma}
\begin{proof}
Fix $g$. If $g$ is not near at $b$ the integral is over a null set. Otherwise $\xi g\in\BoxL{\cL}$ for some
$\xi\in\Ql{b}{T}$, and by Lemma~\ref{l:right} (with $S=\GT{T}$ and $s=\xi^{-1}$) we may replace $g$ by $\xi g$. For
$g\in\BoxL{\cL}$ and $\sigma\in\Hl{b}{T}$ with $\sigma g$ near at $b+1$, Lemma~\ref{l:three}(ii) gives
$\sigma^{-1}f_b\in\vecsOfSize{7}$; so the integral is at most
$\int\mathbf{1}_{\vecsOfSize{7}}(\sigma^{-1}f_b)\Psi(\sigma^{-1}f_b)\,\dd\floorMeasure{\cL}{b}(\sigma)$, which is
$\int_{\vecsOfSize{7}}\Psi\,\dd\quadric{\cL}{b}$ by Lemma~\ref{l:orbit}.
\end{proof}
\leanline{\lean{floorMassLeofchart}}

\begin{lemma}[two masses]\label{l:mass}
In Setting~\ref{set:gauge}: \textup{(i)} for $a<\varrho$ and $x\in\Posn{\cL}$, the function
$u(g)=\floorWeight{\cL}{a}{g}{x}/\floorWeight{\cL}{a+1}{g}{x}$ satisfies
\[\FloorMassLe{\cL}{a}{\Lambda(\Gamma)^{10^{4}n}}{u};\]
\textup{(ii)} for $b\le\varrho$ and $M\in[0,\infty]$ with \m{\HasFibreBound{\Gamma}{b}{\quadric{\cL}{b}}{7}{M}},
\[\FloorMassLe{\cL}{b}{M\cdot\vol(\Gamma_{11})^{|\Omega|}}{1}.\]
\end{lemma}
\begin{proof}
(i) Here $u(g)=\vol(\floorValues{\cL}{a}{g}{x})^{-s}$, $s=\volumeExponent{\cL}$, invariant under $\Ql{a+1}{T}$
(Lemma~\ref{l:valinv}). Let $g\in\BoxL{\cL}$, and take $\varphi$ from (G6) for $g$, the position $x$ and the matrix
$v=0$; put $\Psi(v')=\varphi(\orthonormalPart{\cL}^{*}\JT{T}v')^{s}$. For $\sigma\in\Hl{a}{T}$ and
$u'=\orthonormalPart{\cL}^{*}\JT{T}\sigma^{-1}f_a$, every $w\in W$ gives the vector $\sigma\orthonormalPart{\cL}w$,
which is admissible for $(\sigma g,x)$ and has the coordinate $h_T(\sigma^{-1}f_a,\orthonormalPart{\cL}w)=u'^{*}w$
at $e^*_a$; so $\{u'^{*}w:w\in W\}\subseteq\floorValues{\cL}{a}{\sigma g}{x}$ and
$u(\sigma g)\le\varphi(u')^{s}=\Psi(\sigma^{-1}f_a)$. By Lemma~\ref{l:fibint}(ii), Lemma~\ref{l:fibq}(ii) and (G6),
$\int_{\vecsOfSize{7}}\Psi\,\dd\quadric{\cL}{a}\le\Lambda(\Gamma)^{99n\cdot13}\Lambda(\Gamma)^{99n\cdot12}$. Apply
Lemma~\ref{l:chart}. (ii) Lemma~\ref{l:chart} with $f=1$, $\Psi=1$, and Lemma~\ref{l:fibint}(ii) with $\varphi=1$.
\end{proof}
\leanline{\lean{floorMassLefloorWeightdiv}\sep\lean{floorMassLeone}}

\begin{proof}[Proof of Theorem B\textup{(i)}]
\emph{Reduction to the box.} Let $\xi g\in\BoxL{\cL}$, $\xi=\gamma\nu_0$ with $\gamma\in\Hl{j}{T}$,
$\nu_0\in\Nb{j}{T}$. The weight of $(g,\cdot)$ is that of $(\xi g,\cdot)$, and
$g^{-1}\nu y=(\xi g)^{-1}(\gamma\nu_0\nu\gamma^{-1})\gamma y$; by \Inp{I3} the law of $\gamma\nu_0\nu\gamma^{-1}$ is
again $\etab{j}$. So we may assume $g\in\BoxL{\cL}$, for all $y$. \emph{Induction on $j$.} For $j=0$ both sides
are at most $1$. Let $a=j<\varrho$ and assume the bound for $j$ and all $y$. If the integral for $j+1$ is not $0$,
there is $\tau\in\Nb{j+1}{T}$ with $x_0=g^{-1}\tau y\in\positionsOfSize{\cL}{\Gamma}$, and by invariance we may
replace $y$ by $y_0=\tau y$. Write $\etab{j+1}=\etab{j}*\eta_j$ and
$\Phi(\mathrm{n})=\int\mathbf{1}_{P}(g^{-1}\nu\mathrm{n}y_0)\,\floorWeight{\cL}{j+1}{g}{g^{-1}\nu\mathrm{n}y_0}\,\dd\etab{j}(\nu)$,
$P=\positionsOfSize{\cL}{\Gamma}$. We apply Lemma~\ref{l:heisint} with $u=\gL{\cL}$, size $14$ and
$\Psi(v')=\varphi(\orthonormalPart{\cL}^{*}\JT{T}v')^{s}$, $\varphi$ from (G6) for $g$, the position $x_0$ and
$v=0$. Its hypotheses: if $g^{-1}\nu\mathrm{n}y_0\in P$, then by Lemma~\ref{l:three}(i) the element
$gx_0(g\cdot g^{-1}\nu\mathrm{n}y_0)^{-1}=\mathrm{n}^{-1}\nu^{-1}=\nu''\mathrm{n}^{-1}$, with $\nu''\in\Nb{j}{T}$, has
$\gL{\cL}^{-1}(\,\cdot-1)$ of size $14$. Next, with $x'=g^{-1}\nu\mathrm{n}y_0$, the weight at $j+1$ is the weight at
$j$ times $\vol(\floorValues{\cL}{j}{g}{x'})^{s}$, and $\floorValues{\cL}{j}{g}{x'}\subseteq\mathcal{C}_{u'}$ for
$u'=\orthonormalPart{\cL}^{*}\JT{T}\gL{\cL}^{-1}v$: for an admissible $w'$ with $c=w'_{e^*_j}$ and for $w\in W$,
the two close pairs give $h_T(w',\orthonormalPart{\cL}w)-h_T(w',\nu\mathrm{n}\orthonormalPart{\cL}w)$, and
$(1-\nu\mathrm{n})\orthonormalPart{\cL}w$ is $h_T(v,\orthonormalPart{\cL}w)e_j$ up to terms that $w'$ does not see
(F4), which gives $(cu')^{*}w$ after division by $\gL{\cL}$. So
$\vol(\floorValues{\cL}{j}{g}{x'})^{s}\le\varphi(u')^{s}$, and by the induction hypothesis (with $\mathrm{n}y_0$
for $y$) $\Phi(\mathrm{n})\le\prod_{l<j}\floorConst{\Gamma}{l}\cdot\Psi(\gL{\cL}^{-1}v)$. Finally
$\int_{\vecsOfSize{14}}\Psi\,\dd\vol_{j+1}\le\Lambda(\Gamma)^{99n\cdot20}\Lambda(\Gamma)^{99n\cdot19}$
(Lemmas~\ref{l:fibint}(ii), \ref{l:fibq}(iii), (G6)). Lemma~\ref{l:heisint} gives the bound
$\prod_{l<j}\floorConst{\Gamma}{l}\cdot\floorModulus{\cL}{j}{\gL{\cL}}\Lambda(\Gamma)^{16+3861n}
\le\prod_{l\le j}\floorConst{\Gamma}{l}$, as $c_A\ge1$.
\end{proof}

\begin{proof}[Proof of Theorem B\textup{(ii)}]
The function $\Phi$ is invariant under $\Nb{j}{T}$, a subgroup of $\Ql{b}{T}$. So the left side is
\[\int\Phi(y)\,I(y)\,\dd\rho(y)\qquad\text{with}\qquad I(y)=\int\invWeight{\cL}{j}{x}(\wl{j}{p}\nu y)\,\dd\etab{j}(\nu),\]
and only the $y$ near at $b$ count. We show $I(y)\le\prod_{a<j}\floorModulus{\cL}{a}{p_a^{-1}}\cdot\globalConst{\Gamma}$ for
such $y$. As in the proof of (i), with $\xi y\in\BoxL{\cL}$, $\xi\in\Ql{b}{T}\subseteq\Ql{j}{T}$, the invariance of
$\etab{j}$ and of the inverse weight (Lemma~\ref{l:valinv}) and (F5), we may assume $y\in\BoxL{\cL}$. The proof is
by induction on $j$, with $p$ a head of each length up to $j$; we show that floor $a$ contributes a factor
$\floorModulus{\cL}{a}{p_a^{-1}}\Lambda(\Gamma)^{10^6n^2}$. For the step from $a$ to $a+1$ write
$\etab{a+1}=\eta_a*\etab{a}$ and fix $\nu\in\Nb{a}{T}$. Lemma~\ref{l:heisint} is applied to the function
\[\mathrm{n}\mapsto\invWeight{\cL}{a+1}{x}(\wl{a+1}{p}\mathrm{n}\nu y),\]
with the unit $u=(-\bar p_a)^{-1}$ and $\Psi=\varphi(\orthonormalPart{\cL}^{*}\JT{T}\,\cdot\,)^{s}$, where $\varphi$ comes from (G6) for $y\in\BoxL{\cL}$,
the position $x$ and the matrix $v$ of the coordinates at $e_0,\dots,e_{\varrho-1}$ of the columns of
$\nu\orthonormalPart{\cL}$. If the inverse weight is not $0$, then $\wl{a+1}{p}\mathrm{n}\nu y$ is near at $a+1$,
and Lemma~\ref{l:minor}(iii) (for $\mathsf{K}=\BoxL{\cL}$, the head $\bar p$, the point $\mathrm{n}\nu y$ and
$s=(\mathrm{n}\nu)^{-1}$) bounds the column at $e^*_a$ of $\bar p_a(\nu^{-1}\mathrm{n}^{-1}-1)$, which is the size
condition; and the inverse weight at $a+1$ is the inverse weight at $a$ of $\wl{a}{p}\nu y$ times the power $-s$ of the
volume of $\floorValues{\cL}{a}{\wl{a+1}{p}\mathrm{n}\nu y}{x}$, where this set of floor values contains
$\{u'^{*}w:w\in W\}$ for the vector $u'$ attached to $u^{-1}v$ as in the proof of Lemma~\ref{l:mass}(i); so the
first inequality of (G6) bounds it by $\varphi(u')^{s}$. The integral of $\Psi$ is bounded as in (i), and the
exponents add up to at most $10^6n^2$ for each floor; $j\le n$ floors give $\globalConst{\Gamma}$.
\end{proof}

\sect{Plancherel}{s:plancherel}
\begin{proof}[Proof of Theorem C\textup{(i)}]
\emph{Support.} Let $\sigma\in\Hl{b}{T}$ with $F(\sigma)\ne0$. By (K5) $\sigma x$ is near at $b+1$, and so is
$(\sigma t^{-1})(qx)=(\sigma t^{-1}q\sigma^{-1})\sigma x$, because $\sigma(t^{-1}q)\sigma^{-1}\in\Nb{b}{T}
\subseteq\Ql{b+1}{T}$. Lemma~\ref{l:three}(ii) for $\sigma t^{-1}\in\Hl{b}{T}$ and $qx$ gives
$t\sigma^{-1}f_b\in\vecsOfSize{j+6}$.

\emph{Plancherel.} By Lemma~\ref{l:kinv}, $F$ is invariant under $\Ql{b+1}{T}\cap\Hl{b}{T}$, so on $\Hl{b}{T}$ it is
a measurable function of $\sigma^{-1}f_b$ \Inp{I5}; by Lemma~\ref{l:orbit}(ii) the inner integral is
\[\begin{gathered}G(v)=\int\psi_A\bigl(\xi\,h_T(v+y',v')\bigr)\,\alpha(v')\,\dd\quadric{\cL}{b}(v'),\\
\int|\alpha|^2\,\dd\quadric{\cL}{b}=\int|F|^2\,\dd\floorMeasure{\cL}{b},\end{gathered}\]
with $\alpha$ supported in $\vecsOfSize{j+6}$ by the first part. Write $\quadric{\cL}{b}$ by its fibre bound of
size $j+6$ (Lemma~\ref{l:fibq}(ii)): $v'=E_b\mathsf{k}(P',w',z)$, with fibre masses $\mu'(P')$,
$\int\mu'\le\Lambda(\Gamma)^{99n(j+12)}$. In the coordinates of $E_b$ the form $h_T$ pairs the orthonormal
coordinates with themselves and the fixed coordinates $(t,\pi')$ with the remaining ones $(t',\pi'')$. So for
$v=E_bx$ the phase $\psi_A(\xi h_T(v+y',v'))$ is the product of $\psi_A(\xi\,w(x)^{*}w')$, of a factor of absolute
value $1$ that depends on $x$ and $P'$ only, and of a factor that depends on $(P(x),P',w',z)$. Hence
$|G(E_bx)|\le\int|\widehat{\beta_{P(x),P'}}(w(x))|\,\dd P'$, where $\beta_{P,P'}(w')$ is the integral over $z$ of the
last factor times $\alpha$ and the hat is the Fourier transform $\int\psi_A(\xi\,w^{*}w')\,\cdot\,\dd w'$. We apply
Lemma~\ref{l:fibint}(i) to the measure $\mathsf{m}$ with $G(P,w)=(\int|\widehat{\beta_{P,P'}}(w)|\,\dd P')^{2}$. By the
Cauchy--Schwarz inequality in $P'$ with the weight $\mu'$, by \Inp{I6} in $w$, and by the Cauchy--Schwarz
inequality in $z$ on the fibre, whose mass is at most $\mu'(P')$ because $\alpha$ lives on $\vecsOfSize{j+6}$
(Lemma~\ref{l:basis}),
\[\begin{aligned}\int G(P,w)\,\dd w&\le\int\mu'\cdot\Bigl(\frac{c_A}{\modA(\xi)}\Bigr)^{|\Omega|}\iint\frac{|\beta_{P,P'}(w')|^{2}}{\mu'(P')}\,\dd w'\,\dd P'\\
&\le\Lambda(\Gamma)^{99n(j+12)}\Bigl(\frac{c_A}{\modA(\xi)}\Bigr)^{|\Omega|}\int|\alpha|^{2}\,\dd\quadric{\cL}{b}.\end{aligned}\]
The integrability that \Inp{I6} needs comes from Lemma~\ref{l:succ}(ii).
\end{proof}

\begin{proof}[Proof of Theorem C\textup{(ii)}]
We may assume $\int\Phi\,\dd\eta_a\ne0$; then $\Phi(m)\ne0$ for some $m\in\Nf{a}{T}$ \Inp{I3}, and $m\in S$ by
(o). Take $t$, $q$ from (h) for this $m$ and write $t=t_0r$ with $t_0\in\Hl{a+1}{T}$, $r\in\Nf{a}{T}$ (F1). By
\Inp{I3} and \Inp{I1}, $\eta_a$ is invariant under $\nu\mapsto r^{-1}\nu\,rm$ and under conjugation by $t_0$, so
\[\int\Phi(\nu)\,\dd\eta_a(\nu)=\int\Phi\bigl(t^{-1}\nu\,tm\bigr)\,\dd\eta_a(\nu).\]
We apply Lemma~\ref{l:heisint} to $\nu\mapsto\Phi(t^{-1}\nu tm)$, with $u$ and the size $j$. \emph{The size
condition.} If $\Phi(\nu_1)\ne0$ for $\nu_1=t^{-1}\nu tm$, then $\nu_1\in S$ by (o), and (h) gives
$\nu'\in\Nb{a}{T}$ with the column at $e^*_a$ of $u^{-1}(\nu'tm(t\nu_1)^{-1}-1)=u^{-1}(\nu'\nu^{-1}-1)$ in
$\vecsOfSize{j}$. \emph{The pointwise bound.} By (o), $\Phi(\nu_1)\le C|\upperPairing{\cL}{a+1}{k}{\nu_1^{\varepsilon}x}{\nu_1y}|^2$.
By Lemma~\ref{l:succ}(ii) this average is the integral over $\sigma\in\Hl{a+1}{T}$ of
$\upperPairing{\cL}{a+2}{k}{c^{\varepsilon}\sigma x}{c\,\sigma y}$ with $c=\sigma\nu_1\sigma^{-1}\in\Nf{a}{T}$,
which by (K2) and (e) is $\psi_A(\xi\,\sd(c))\cdot F_2(\sigma)$, $F_2(\sigma)=\upperPairing{\cL}{a+2}{k}{\sigma x}{\sigma y}$.
By (F6), $\sd(\sigma\nu_1\sigma^{-1})=-h_T(v_1,\sigma^{-1}f_{a+1})$, and the vector $v_1$ of $\nu_1$ is
$t_0^{-1}$ applied to the sum of the vector $v$ of $\nu$ and a vector that does not depend on $\nu$. So, with
$\xi'=-\xi\bar u$ and a fixed $y'\in V$,
\[\begin{gathered}\Phi(\nu_1)\le C\,\Psi(u^{-1}v),\\
\Psi(v')=\Bigl|\int\psi_A\bigl(\xi'\,h_T(v'+y',t_0\sigma^{-1}f_{a+1})\bigr)F_2(\sigma)\,\dd\floorMeasure{\cL}{a+1}(\sigma)\Bigr|^{2}.\end{gathered}\]
\emph{The integral of $\Psi$.} Theorem~C(i) at $b=a+1$, with $t_0$, the element $tq$ (for which
$t_0^{-1}tq=rq\in\Nb{a+1}{T}$ and \m{\HasMatSize{i}{tqx}}), the unit $\xi'$ and the measure $\vol_{a+1}$ with its
fibre bound (Lemma~\ref{l:fibq}(iii)), gives
\[\int_{\vecsOfSize{j}}\Psi\,\dd\vol_{a+1}\le\Lambda(\Gamma)^{99n(j+6)}\Lambda(\Gamma)^{99n(i+12)}\Bigl(\frac{c_A}{\modA(\xi u)}\Bigr)^{|\Omega|}\int|F_2|^2\dd\floorMeasure{\cL}{a+1}.\]
Lemma~\ref{l:heisint} multiplies this by $C\,\floorModulus{\cL}{a}{u}\Lambda(\Gamma)^{j+2}$, and
$99n(i+j+18)+j+2\le100n(i+j+31)$.
\end{proof}

\sect{Where the results are used}{s:used}
Note~L uses almost everything: Theorem~A in the descent between floors and in every change of variables; Theorems B
and~C and Lemmas~\ref{l:cs}, \ref{l:mass}, \ref{l:minor}, \ref{l:three}, \ref{l:subone} in the three bounds of the
local theorem and in the integration over positions. Note~C constructs the level and the gauges at the complex
place, proves there that the tower measure is a Haar measure (the hypothesis of Theorem~A(iii)), and verifies the
axioms (G1)--(G6). No other unit of the artifact uses this one directly except through these two.

\sect{Notes on the formal text}{s:notes}
\textbf{Implicit arguments.} The size is an implicit argument of $\operatorname{sz}$, $\mathrm{Vec}$ and
$\operatorname{Hd}$; a gauge is a size with further data, and for a gauge these predicates are applied to it. The
level is an explicit argument of all notions of Section~\ref{s:setting} except the volume exponent, whose only
argument it is.

\textbf{Argument lists not confirmed by Lean.} For 43 declarations of the five files the reader's dossier derived the
list of arguments by Lean's rules for section variables, because the census reads a name that the declaration binds
anew ($n$, $i$, $a$, $j$, $x$) as the section variable of that name, misses the section variables after a
\texttt{let} in a statement, or drops the variable \texttt{m} where a hypothesis has that name. Nothing was compiled,
so these lists are not confirmed by Lean; among them are the definitions of sizes, of vectors and matrices of a
size, of heads, of floor values, of the mass of a floor, of the averaged kernel, of pairing kernels, of the
self-norm, of gauges and of fibre bounds. The statements above follow the Lean text of each declaration. Three
statements continue after a \texttt{let} and were read to their end: Lemma~\ref{l:minor}(iii),
Lemma~\ref{l:heis} and Theorem~C(i).

\textbf{Differences from the natural statements.} (1) Averaged kernels are Bochner integrals, $0$ when the
integrand is not integrable; (K4) makes them honest integrals for pairing kernels. (2) In (L2) the third condition
is stated with $\UT{T}$, which is $\Ql{\varrho+2}{T}$; it is the bound \m{\FloorMassLe{\cL}{n}{1}{1}}. (3) The Weyl
law is defined with the exponent $k+1+2(\varrho-(l+1))$ in truncated subtraction for all $l$, and used for
$l<\varrho$. (4) A gauge allows $c_A=\infty$; the bounds are then empty. (5) In Lemma~\ref{l:row} and
Lemma~\ref{l:backsub} the hypothesis runs over the $a<\varrho$ with $a<j$; $j$ may exceed $\varrho$. (6) In
Lemma~\ref{l:rules}(d) the Lean text takes the first matrix with columns indexed by the basis of $V$ and any finite
set of rows. (7) Lemma~\ref{l:factor} takes two $s$-finite measures and does not ask that $\rho$ be a probability.
(8) In Theorem~B(i) the indicator of the positions of size $1$ and the weight are one function in the Lean text
(an indicator of the weight). (9) In Lemma~\ref{l:weylratio} the product runs over $a<j$ with exponents
$k+2(\varrho-a)$, the rank of $V_{a+1}$ plus one. (10) The count of the basis vectors of floor at least $a+1$,
$k+1+2(\varrho-(a+1))$, is a lemma of the formal text; it is used as the rank of $V_{a+1}$. (11) One line of the
file of the charts that the census gives to the invariance of the inverse weight belongs to the next declaration.

\textbf{Small results.} The declarations without a statement of their own are cited where they are used: the
matrix and the inverse of a Weyl ratio after Definition~\ref{d:wr}; the two rules for floor measures after
Definitions~\ref{d:fm} and~\ref{d:tower}; nearness after Definition~\ref{d:near}; the rules for weights after
Definition~\ref{d:weight} and Lemma~\ref{l:valinv}; measurability after Definition~\ref{d:consts}; the bound
$\Lambda\ge1$ after Definition~\ref{d:gauge}; and the count of (10) in Lemma~\ref{l:minor}.

\subhead{How this note was written.} From the structured proofs of the reader's dossier and from the Lean text. I
read the five Lean files completely, statements and proofs, and the definitions of the background library that they
use (shapes, split Gram matrices, the groups of the tower, Weyl and Heisenberg elements, places, flat charts,
quadric measures, hyperbolic bases, convolutions). The proofs of Sections~\ref{s:sizes}, \ref{s:tower},
\ref{s:near} and \ref{s:kernels}, of Lemmas~\ref{l:heis}, \ref{l:basis}, \ref{l:fibint}, \ref{l:heisint},
\ref{l:three}, \ref{l:subone}, \ref{l:chart}, \ref{l:mass} and of Theorem~C(ii) follow the formal proofs step by
step; the exponents in them were recomputed. The proofs flagged by the reader as least certain were read line by
line in the Lean text: the two integrals of Theorem~B, the two parts of Theorem~C, and Lemma~\ref{l:fibq}(i). For
these the text gives the argument as I could rebuild it from the terms: which lemma is applied to which objects,
and with which sizes. Three places are less than a full proof and say so here: in Lemma~\ref{l:fibq}(i) the
comparison of a fibre of the quadric with the cone uses a library lemma on measures of sets whose differences lie
in a given set, whose statement I took from its use; in Theorem~B the inclusions of sets of floor values in the
sets of (G6) are checked in the formal text by a computation with Heisenberg elements that I describe and did not
redo entry by entry, and the exponent per floor in Theorem~B(ii) was not recomputed; in Theorem~C(i) the splitting
of the phase in the coordinates of the floor basis is the formal text's computation. Lemma~\ref{l:weylratio} and
Lemma~\ref{l:orbit}(i) are written from the formal proofs, with the change of variables described and not carried
out. The program check of the unit's dossier covered \UnitSteps{} leaf steps of \UnitProofs{} proofs.

\head{Appendix A. Concordance with the Lean artifact}
The letters of this note and the names of the formal text: $\cL$ is \texttt{F}; $\Omega$ is \texttt{ι};
$\mathfrak{m}$ is \texttt{m}; $\mathcal{S}$ is \texttt{S}; $\Gamma_i$ and $\mathcal{S}_i$ are \texttt{ofSize i};
$\modA$, $\psi_A$, $c_A$, $\tau_A$, $\delta_A$, $\nu_A$ are the fields \texttt{md}, \texttt{ψ}, \texttt{c},
\texttt{τ}, \texttt{skewUnit}, \texttt{skewMeasure} of the place; $\Hl{l}{T}$, $\Ql{l}{T}$, $\Nb{l}{T}$,
$\Nf{l}{T}$, $\UT{T}$ are \texttt{GT}, \texttt{unitriGroup}, \texttt{unipotentBelow}, \texttt{floorUnipotent},
\texttt{besselU}; $\mathrm{n}_a$ is \texttt{heisenbergElem}; $\eta_l$, $\etab{j}$, $\varpi_l$ are
\texttt{heisenbergMeasure}, \texttt{ascConv heisenbergMeasure}, \texttt{weylLaw}; $\quadric{\cL}{b}$ and $\vol_b$
are \texttt{floorQuadric} and \texttt{floorVolume}; $E_b$ is \texttt{floorBasis}; $\sd$ is \texttt{simpleRootSum};
$\sigma^{-1}f_b$ is \texttt{orbitPoint}. Two statements as the artifact has them:
\quoted{stackMeasuremconvdirac}
\quoted{pivotModulussqle}
{\fontsize{8.6}{11.6}\selectfont\renewcommand{\leanfont}{\ttfamily\fontsize{7.7}{10}\selectfont\color{leancol}}
\setlength{\LTleft}{0pt}\setlength{\LTright}{\fill}\setlength{\tabcolsep}{0pt}
\begin{longtable}{@{}p{30mm}p{56mm}p{43mm}r@{}}
\textbf{In this note} & \textbf{Lean declaration} & \textbf{File} & \textbf{Lines}\\\midrule\endhead
\ConcordanceRows
\end{longtable}}

\head{Notation}
\begingroup\fontsize{8.6}{10.9}\selectfont\raggedright
\setlength{\columnsep}{9mm}\setlength{\parskip}{0pt}
\begin{multicols}{2}
\nt{$A$, $\bar x$, $\dd x$, $A^{-}$}{the place; involution; Haar measure; skew elements}
\nt{$\modA$, $\psi_A$, $c_A$, $\tau_A$, $\delta_A$, $\nu_A$}{module, character, Plancherel constant, trace-one element, skew unit, measure on $A^{-}$}
\nt{$\mathfrak{s}=(\varrho,\kappa,k_0)$, $n$, $k$, $r$}{the shape; $\varrho+1$; $|\kappa|$; $2\varrho+k+1$}
\nt{$V$, $V_l$, $e_a$, $e^*_a$, $x_i$, $z$, $f_a$}{the space, its floors, basis, frame}
\nt{$T$, $\JT{T}$, $h_T$, $\GT{T}$, $\ZT{T}{g}$}{form, split Gram matrix, hermitian form, unitary group, frame Gram matrix}
\nt{$\Hl{l}{T}$, $\Ql{l}{T}$, $\Nb{l}{T}$, $\Nf{l}{T}$, $\UT{T}$}{the subgroups of Section~\ref{s:space}}
\nt{$\sd$, $\wl{l}{p}$, $\wo{a}{c}$, $\mathrm{n}_a(v,\beta)$}{superdiagonal sum; Weyl elements; Heisenberg element}
\nt{$\Omega$, $\mathfrak{m}$, $E$, $E_b$, $\pi_b$}{parameters of a hyperbolic basis; basis matrices; index set of the extra pairs}
\nt{$\cL$, $E(\cL)$, $\gL{\cL}$, $\theta$, $\mu_H$}{a level and its data}
\nt{$\Ball{\cL}$, $\BoxL{\cL}$, $\Posn{\cL}$}{ball, box, positions}
\nt{$\operatorname{Lev}$, $\operatorname{Siz}$, $\operatorname{Gau}$}{levels, sizes, gauges}
\nt{$\floorMeasure{\cL}{l}$, $\eta_l$, $\etab{j}$, $\varpi_l$}{floor measure, Heisenberg measures, Weyl law}
\nt{$\stackMeasure{\cL}_{a,i}$, $\groupMeasure{\cL}$, $\upperMeasure{\cL}{l}$}{stack, tower and upper measures}
\nt{$\quadric{\cL}{b}$, $\vol_b$, $\vol$, $\delta_s$, $*$}{quadric measure, volume of $V_b$, volume, unit mass, convolution}
\nt{$\mathcal{S}$, $\mathcal{S}_i$, $\Gamma$, $\Gamma_i$, $\Lambda$, $\Gamma_S$}{a size, a gauge, its constant, a second gauge}
\nt{$\operatorname{sz}$, $\mathrm{Vec}_i$, $\operatorname{Hd}_j$}{size of a matrix, vectors of size $i$, heads}
\nt{$\operatorname{Near}_j$, $\mathsf{O}$, $\operatorname{Cl}$, $\operatorname{Val}_a$}{near points, orthonormal part, close pairs, floor values}
\nt{$\mathsf{s}$, $\mathsf{w}_j$, $\mathsf{w}^{-1}_j$, $\mathsf{d}_l$, $\operatorname{Mass}_j$}{volume exponent, weight, inverse weight, floor module, mass}
\nt{$\mathsf{Av}_j$, $\mathsf{Pv}_j$, $\operatorname{PK}_j$}{averaged kernel, pivot, pairing kernels}
\nt{$\mathsf{Sq}_j$, $\mathsf{Sm}_j$, $\operatorname{SN}_j$}{weighted square, its sum, self-norm}
\nt{$\operatorname{Pos}$, $C_l$, $C_{\mathrm{gl}}$, $\operatorname{Fib}$}{positions of size $1$, constants, fibre bound}
\nt{$W$, $\varphi$, $\mathcal{C}_u$}{the objects of axiom (G6)}
\nt{$P$, $w$, $t$, $t'$, $\pi'$, $\pi''$, $\mathsf{k}$, $Z$, $\zeta$, $\mu$, $\rho'$}{coordinates and data of a fibre bound}
\nt{$p$, $q$, $d$}{sequences of units; also $q$ a group element or a measure, $d$ the element of a pairing kernel or a size}
\nt{$g$, $x$, $y$, $s$, $u$, $v$, $m$, $h$}{group elements, vectors, units; $u$, $v$ also functions}
\nt{$\sigma$, $\tau$, $\nu$, $\gamma$, $\xi$, $\eta$, $\alpha$, $\beta$}{variables of integration, group elements, units, skew elements}
\nt{$i$, $j$, $l$, $a$, $b$, $c$, $N$}{natural numbers: sizes, floors, ranks; $c$ also an element of $A$}
\nt{$k$, $K$, $C$, $M$, $V$}{a kernel (also $k=|\kappa|$); constants}
\nt{$\mathsf{K}$, $\mathsf{m}$}{the set of Setting~\ref{set:two}; the measure on $V$ of Theorem~C(i)}
\nt{$E(D)$, $G_1(\mathfrak{s})$}{the CM field of the admissible data $D$ of the overview; the group at the complex place (only in Section~\ref{s:intro})}
\nt{$d$, $\lambda$, $K$, $f$, $\Theta$}{letters that are local to the notes of this unit and differ from those of the overview; this note uses $d$ (an element or a size), $\lambda$ (measures), $K$ (constants) and $f$ (functions)}
\nt{$f$, $F$, $G$, $\Phi$, $\Psi$, $D$, $X$, $Y$, $S$, $I$}{auxiliary functions, matrices, sets}
\nt{$\varepsilon$, $\rho$, $\lambda$, $\iota$, $\theta$, $\psi$, $\varrho$, $\kappa$, $\Lambda$, $\delta$}{an exponent; measures; an index set; letters of names}
\nt{$\mathbf{1}$, $\infty$, $L$, $H$, $O$, $Q$, $U$, $B$, $e$, $n$, $o$, $r$, $z$}{indicator; infinity; letters inside names}
\end{multicols}
\endgroup

\end{document}
